\documentclass[10pt,oneside]{book}
\usepackage[OT1]{fontenc}
\usepackage{geometry}


\usepackage[style = alphabetic, backend=biber, useprefix=true]{biblatex}
\addbibresource{refs.bib}
\renewcommand{\multicitedelim}{\addcomma\space}

\usepackage{hyperref}
\usepackage[dvipsnames]{xcolor}
\hypersetup{
    colorlinks = true,
    linktoc=page,
    urlcolor = WildStrawberry,
    citecolor = Turquoise,
    linkcolor = WildStrawberry,
    filecolor= WildStrawberry
}
\usepackage{enumitem}
\usepackage{tikz,tikz-cd}
\usepackage{changepage}   % for the adjustwidth environment

\usepackage{stmaryrd}
\usepackage{epigraph}
\usepackage{titlesec}
\usepackage{quiver}
\usepackage[mathscr]{euscript}
\usepackage{dsfont}
\usepackage{amssymb,amsthm,amsmath,mathtools}
\usepackage{wasysym}

\setlength\epigraphwidth{0.8\textwidth}
\setlength\epigraphrule{0pt}
\setlength{\columnsep}{1cm}
\parindent = 10pt

% useful formatting shortcuts
\def\~#1{\mathscr{#1}}
\def\cal#1{\mathcal{#1}}
\def\To#1{\xrightarrow{#1}}
\newcommand{\quadd}{\quad\quad}
\newcommand{\figref}[1]{Figure \ref{#1}}
\def\pre#1{\widehat{#1}}


% sets
\newcommand{\C}{\mathbb{C}}
\newcommand{\N}{\mathbb{N}}
\newcommand{\Q}{\mathbb{Q}}
\newcommand{\R}{\mathbb{R}}
\newcommand{\Z}{\mathbb{Z}}

% category theory symbols
\newcommand{\op}{\mathrm{op}}
\newcommand{\adj}{\dashv}
\newcommand{\dom}{\mathrm{dom}}
\newcommand{\colim}{\mathop{\mathrm{colim}}}
\newcommand{\ob}{\mathop{\mathrm{ob}}}
\newcommand{\cod}{\mathrm{cod}}
\newcommand{\Hom}{\mathrm{Hom}}
\newcommand{\isom}{\xrightarrow{\sim}}
\newcommand{\epi}{\twoheadrightarrow}
\newcommand{\mono}{\hookrightarrow}
\newcommand{\coker}{\mathop{\mathrm{coker}}}
\newcommand{\Ind}{\mathrm{Ind}}
\newcommand{\Coh}{\mathrm{Coh}}
\newcommand{\QCoh}{\mathrm{QCoh}}
\newcommand{\Sh}{\mathrm{Sh}}
\newcommand{\PSh}{\mathrm{PSh}}
\newcommand{\im}{\mathrm{im}}
\newcommand{\id}{\mathrm{id}}



% AG/algebra/number theory symbols
\newcommand{\Sym}{\mathrm{Sym}}
\newcommand{\lcm}{\mathrm{lcm}}
\newcommand{\Spec}{\mathrm{Spec}}
\newcommand{\Fun}{\mathrm{Fun}}
\newcommand{\ev}{\mathrm{ev}}



% all the categories
\newcommand{\cAlg}{\mathrm{c}\~{A}\mathrm{lg}}
\newcommand{\cAlgop}{\cAlg^\op}
\newcommand{\Set}{\~S\mathrm{et}}
\newcommand{\Rel}{\~R\mathrm{el}}
\newcommand{\Grp}{\~G\mathrm{rp}}
\newcommand{\Ab}{\~A\mathrm{b}}
\newcommand{\Ring}{\~R\mathrm{ing}}
\newcommand{\cRing}{\mathrm{c}\~{R}\mathrm{ing}}
\newcommand{\Mod}{\~M\mathrm{od}}
\newcommand{\Vect}{\~V\mathrm{ect}}
\newcommand{\fdVect}{\mathrm{fd}\~{V}\mathrm{ect}}
\newcommand{\Top}{\~T\mathrm{op}}
\newcommand{\Cat}{\~C\mathrm{at}}
\newcommand{\Mon}{\~M\mathrm{on}}
\newcommand{\Open}{\~{O}\mathrm{pen}}
\newcommand{\Rep}{\~R\mathrm{ep}}
\newcommand{\Field}{\~F\mathrm{ield}}
\newcommand{\Bilin}{\~{B}\mathrm{ilin}}

% misc
\newcommand{\vp}{\varphi}
\newcommand{\ve}{\varepsilon}
\newcommand{\sbullet}{\,\begin{picture}(-1,1)(-1,-3le*{3}\end{picture}\ }

\DeclareFontFamily{U}{dmjhira}{}
\DeclareFontShape{U}{dmjhira}{m}{n}{ <-> dmjhira }{}
\DeclareRobustCommand{\yo}{\text{\usefont{U}{dmjhira}{m}{n}\symbol{"48}}}


%theorem environments
\theoremstyle{plain}
\newtheorem{thm}{Theorem}[section]
\newtheorem*{thm*}{Theorem}
\newtheorem{prop}[thm]{Proposition}
\newtheorem{lem}[thm]{Lemma}
\newtheorem*{lem*}{Lemma}
\newtheorem{cor}[thm]{Corollary}
\newtheorem*{cor*}{Corollary}

\theoremstyle{definition}
\newtheorem{defn}[thm]{Definition}
\newtheorem*{defn*}{Definition}
\newtheorem{ex}[thm]{Example}
\AtEndEnvironment{ex}{\null\hfill$\Diamond$}
\newtheorem{exs}[thm]{Examples}
\AtEndEnvironment{exs}{\null\hfill$\Diamond$}
\theoremstyle{remark}
\newtheorem*{rmk}{Remark}
\newtheorem*{warn}{Warning}
\newtheorem*{rmks}{Remarks}
\newtheorem*{notation}{Notation}


%symbols 

\usepackage[automake]{glossaries-extra}
\GlsXtrEnableLinkCounting{general} 

\renewcommand*{\glslinkpresetkeys}{
 \setupglslink{hyper=false}} 


\makeglossaries

\newglossaryentry{0}%
{%
  name={$\mathbf{0}$},
  text={\mathbf{0}},
  description={empty category,},
  sort={0}
}
\newglossaryentry{I}%
{%
  name={$\mathbf{I}$},
  text={\mathbf{I}},
  description={interval category,},
  sort={I}
}
\newglossaryentry{Set}%
{%
  name={$\Set$},
  text={\Set},
  description={category of sets,},
  sort={S}
}
\newglossaryentry{adj}%
{%
  name={$\adj$},
  text={\adj},
  description={adjunction,},
  sort={!}
}

% simple environment to put a box around a given tikzcd diagram.
\newenvironment{boxedtikzcd} 
{\begin{lrbox}{\boxedtikzcdbox}\begin{tikzcd}}
{\end{tikzcd}\end{lrbox}\fbox{\usebox{\boxedtikzcdbox}}}
\newsavebox{\boxedtikzcdbox}

% change equation/figure numbering
\numberwithin{equation}{section}
\numberwithin{figure}{section}

% change autoref for sections
\renewcommand{\sectionautorefname}{\S}


% format sections
\usepackage{titlesec}
\titleformat{\section}
  {\normalfont\Large\bfseries}{\S\thesection}{1em}{}
\titleformat{\subsection}
  {\normalfont\large\itshape}{\thesubsection}{1em}{}
\begin{document}

\title{Category Theory}
\author{Elias Vandenberg}
\maketitle

\newpage

\tableofcontents 

% \chapter*{Introduction}

% The purpose of this book is to provide a guided overview of the key concepts of modern category theory. Our goal is to present, in a single volume, a self-contained account of most of the major topics in modern category theory, ranging from basic definitions to $n$-categories. For the most part, this text takes a homotopy-theory agnostic approach to category theory, which means concepts like model categories, simplicial sets, and $\infty$-categories will not be covered in this text. Instead, the presentation is focused on the concrete algebraic properties of categorical structures, and is organized so that the reader can achieve a rapid mastery of the essential topics. Indeed, many deep topics in modern category theory can be understood without much knowledge of homotopy theory, and it is these topics which we focus on here.

\section*{Prerequisites \& Conventions}

Very few prerequisites are assumed in this text. We assume the reader is familiar with the basic notions of algebra, such as groups, rings and fields, however, we do not expect readers to be experts on these subjects. Throughout this text, we provide numerous examples, some of which may come from areas such as topology or representation theory. Readers who are not comfortable with these subjects should feel free to skip over these examples and return to them as needed.

Since this book contains several lengthy examples and proofs, we make use of the following symbols to indicate the end of various items:\\

 $\Diamond$ indicates the end of an example

 $\Box$ indicates the end of a proof


\part{Basic definitions and constructions}
\chapter{Categories}

\epigraph{\itshape This concept of category is an omni-purpose affair... It frames a possible template for any mathematical theory: the theory should have nouns and verbs, i.e., objects, and morphisms, and there should be an explicit notion of composition related to the morphisms; the theory should,
in brief, be packaged by a category. There is hardly any species of mathematical object that doesn't fit into this convenient, and often enlightening, template.}{---Barry Mazur, \emph{When is One Thing Equal to Some Other Thing?} \cite{mazur2007}}

\section{Basic definitions and examples}

\begin{defn}
    A \emph{category} $\~{C}$ consists of:
    \begin{itemize}
        \item A collection $\ob (\~{C})$ of \emph{objects} $A, B, C, \cdots$
        
        \item For each $A$ and $B$ in $\ob (\~{C})$, a collection $\~C(A, B)$ or $\Hom_\~{C} (A, B)$ or $\Hom (A,B)$ of \emph{morphisms} or \emph{arrows} or \emph{maps} $f, g, h, \cdots$ from $A$ to $B$, such that every $f$ in $\~C(A, B)$ has specified \emph{domain} and \emph{codomain} objects $\dom(f)$ and $\cod(f)$ respectively. If $\dom(f) = A$ and $\cod(f) = B$, we write $f: A \to B$ or $A \To{f} B$.
        
        \item For each composable pair of morphisms  
        \[A \To{f} B \To{g} C\]
        $(\text{i.e., } \cod(f) = \dom(g))$, we require that there exist a third morphism $g \circ f$ with $\dom(g \circ f) = A$ and $\cod(g \circ f) = B$, called the \emph{composition} of $f$ and $g$ and often written as $gf$
        
        \item For each $A \in \ob (\~{C})$, we require the existence of an arrow $1_A: A \to A$ called the \emph{identity} on $A$.
    \end{itemize}
    These data are subject to the following axioms:
    \begin{enumerate}[label=(\roman*)]
        \item \emph{Associativity}: for any triple of arrows in the configuration 
        \[A \xrightarrow{f} B \xrightarrow{g} C \xrightarrow{h} D\]
        we always have 
        \[h \circ( g \circ f) = (h \circ g) \circ f\]
        \item \emph{Identity law}: for each $f \in \~C (A, B)$, we have $f \circ 1_A = 1_B \circ f = f$.
    \end{enumerate}
\end{defn}

\begin{rmks}
    \hfill
    \begin{itemize}
    
    \item The word ``collection'' can be thought of as meaning ``set'' in many contexts, however, if the reader is familiar with such things, it is better to interpret it as meaning ``class''.
        
    \item Other authors may use the notation $\text{Mor}(A, B)$ when referring to the collection of maps between $A$ and $B$. 
    
    
    \item We often write $A \in \~{C}$ to mean $A \in \ob(\~C)$.
    
    \item  In general, the definition of a category is set up so that from each string 
        \[A_0 \To{f_1} A_1 \To{f_2} \cdots \To{f_n} A_n\]
        of maps in $\~{C}$, we can construct exactly one map $A_0 \to A_n$, given by the $n$-fold composition $f_n \cdots f_1$.
    
    \item Composition in a category also gives rise to \emph{diagrams}, which are arrangements of morphisms like
    \[\begin{tikzcd}
        A & B \\
        & C
        \arrow["g", from=1-1, to=1-2]
        \arrow["f"', from=1-1, to=2-2]
        \arrow["h", from=1-2, to=2-2]
    \end{tikzcd}\]
    We say a diagram like the one above is \emph{commutative} if $f = h \circ g$. In this situation, we might also say that $f$ \emph{factors through} $g$ and $h$. In general, a diagram is said to commute if whenever there are two paths from an object $X$ to an object $Y$, the map from $X$ to $Y$ obtained by composing along one path is equal to the map obtained by composing along the other. When drawing commutative diagrams, an identity arrow $1_A: A \to A$ is usually denoted by an elongated `equals' sign, as in the following triangle:
    \[\begin{tikzcd}
        A & B \\
        & A
        \arrow["f", from=1-1, to=1-2]
        \arrow[Rightarrow, no head, from=1-1, to=2-2]
        \arrow["g", from=1-2, to=2-2]
    \end{tikzcd}\]
    \end{itemize}
\end{rmks}

Returning to the first remark above, we can classify a category $\~C$ based on how big $\ob(\~C)$ is, as well as by the size of $\~C(A, B)$ for each $A, B \in \ob(\~C)$.

\begin{defn}
    Let $\~C$ be a category. We say $\~C$ is \emph{small} if all of its objects form a set and all of its morphisms form a set. We say $\~C$ is \emph{locally small} if for all $A, B \in \ob(\~C)$, $\~C(A, B)$ contains only a set's worth of arrows. A category which is not small is (naturally) called \emph{large}.
\end{defn}

While the above definitions are sometimes ignored in settings outside of pure category theory, they will be significant to us. For instance, the fundamental concept of a limit is defined using the notion of a small category, and Yoneda's lemma, one of the most profound results in category theory, is stated in terms of locally small categories. It will sometimes be useful for us to distinguish between small and large categories, in which case we will write small categories using bold-face letters $\mathbf{A}, \mathbf{B}, \mathbf{C}, \ldots$

\begin{ex} \label{ex:all_the_categories} We now explore some well-known categories.
\begin{enumerate}[label=(\roman*)]
        \item The simplest categories to describe are very small categories, which are often presented simply as collections of objects and arrows. For instance, there is the category $\gls{0}$ with no objects and no morphisms, and $\mathbf{1}$ is the category with a single object $\ast$ and a single (identity) arrow. We often draw $\mathbf{1}$ like
        \[\begin{tikzcd}
        \ast \arrow[loop,swap,looseness=6] 
        \end{tikzcd}\]
        There is also a category \textbf{2} with two objects $A, B$, and two identity arrows $1_A$ and $1_B$. Going one step further, the \emph{interval category} $\gls{I}$ is the category with two objects $A$ and $B$ and a single (non-identity) arrow $A \to B$. We can draw $\mathbf{I}$ like
        \[\begin{tikzcd}
        A \arrow[loop left]{}{1_A} \arrow{r}{f} & B \arrow[loop right]{}{1_B}
        \end{tikzcd}\]
        It is common to omit identity arrows (and sometimes the object names) when drawing very small categories, so, for instance, we can also draw $\mathbf{I}$ above as
            \[\ast \to \ast \quad \text{or} \quad A \To{f} B\]
        It is easy to come up with more complicated examples. For instance, one can check that that the following are well-defined categories:
        \[\begin{tikzcd}
            \ast & \ast
            \arrow[shift left=1, from=1-1, to=1-2]
            \arrow[shift right=1, from=1-1, to=1-2]
        \end{tikzcd} \quad \begin{tikzcd}
        & B \\
        A && C
        \arrow["f", from=2-1, to=1-2]
        \arrow["g", from=1-2, to=2-3]
        \arrow["g \circ f"', from=2-1, to=2-3]
        \end{tikzcd} \quad \begin{tikzcd}
    & \ast & \ast \\
    \ast & \ast & \ast
    \arrow["kj"', from=1-2, to=2-1]
    \arrow["j"', from=1-2, to=2-2]
    \arrow["{hj=gf}"{description}, from=1-2, to=2-3]
    \arrow["f", from=1-2, to=1-3]
    \arrow["g", from=1-3, to=2-3]
    \arrow["h"', from=2-2, to=2-3]
    \arrow["k", from=2-2, to=2-1]
\end{tikzcd}\]
        
        \item The prototypical category is $\gls{Set}$. Its objects are sets $A, B, \cdots$, and its morphisms are just maps of sets (i.e., functions) $A \to B$. Composition is just given by function composition, and identities are just identity maps. 
        
        \item The category $\Rel$ has sets as objects and binary relations as morphisms. For objects $A, B \in \Rel$ there is a morphism $R: A \to B$ precisely when $R \subseteq A \times B$. Given two relations $A \To{R} B \To{S} C$, we define the relation $RS$ to be the set
        \[RS := \{(a, c) \in RS ~|~ \, \exists \, b \in B \text{ such that } (a, b) \in R \text{ and } (b, c) \in S \} \subseteq A \times C\]

        \item A \emph{preordered set} or \emph{poset} is a set $P$ together with a reflexive transitive relation $\leq$. We can view $P$ as a category where for any two $p, q \in P$, a map $p \to q$ means $p \leq q$. A similar concept is that of a \emph{partially ordered set}, which is a set $P'$ whose relation $\leq$ also needs to satisfy the antisymmetry property, which simply says that $p' \leq q' \text{ and } q' \leq p' \implies p' = q'$ for all $p', q' \in P'$. We can regard $P'$ as a category in exactly the same way as $P$. These are both examples of categories where there is at most one morphism between any two objects, i.e., where $\begin{tikzcd}
    X & Y
    \arrow["g"', shift right, from=1-1, to=1-2]
    \arrow["f", shift left, from=1-1, to=1-2]
\end{tikzcd}$ always implies  $f = g$. Categories with this property are called \emph{posetal} or \emph{thin}, and examples of them abound. For instance, posetal categories give rise to many fundamental examples of sheaves, one of the central objects of study in algebraic geometry. 
        
        \item The category $\Grp$ has groups as objects and group homomorphisms as arrows. There is also a category of abelian groups and their homomorphisms, denoted $\Ab$.
        
        \item The category $\Ring$ has unital rings as objects, and its arrows are ring homomorphisms $\vp: R \to S$ such that $\vp(1_R) = 1_S$ (where $1_R$ is the unity in $R$ and $1_S$ is the unity in $S$). Similarly, $\cRing$ has commutative unital rings as objects, and homomorphisms between these rings as maps. %There is also a category of non-unital rings, denoted \textbf{Rng}.
        
        \item The category $_R\Mod$ has left $R$-modules as objects and $R$-module homomorphisms as maps. For right $R$-modules, we have the category $\Mod_R$, defined analogously. 
        
        \item $\Mod$ is the category of all modules. An object in $\Mod$ is a pair $(R, M)$ consisting of a commutative ring $R$ and a left $R$-module $M$. An arrow $(R, M) \to (R', M')$ is a pair $(\vp, \mu)$ where $\vp: R \to R'$ is a ring homomorphism, and $\mu: M \to M'$ is a function that allows us to regard $M'$ as an $R$-module. The action of $R$ is given by
        \[\mu(r m) := \vp(r) \mu(m)\]
        for all $r \in R$, $m \in M$, where the multiplication on the left-hand-side is the action of $R$ on $M$, and the multiplication on the right is the action of $S$ on $M'$. Since $\vp(r) \in S$, this map (often called a \emph{restriction of scalars}) defines an action of $R$ on $M'$, allowing us to view $M'$ as an $R$-module. With objects and morphisms defined in this manner, one can easily check that $\Mod$ is a category.
        
        \item The category $\Vect_k$ has vector spaces over the field $k$ as objects, and linear maps as morphisms. This category is a special case of $\Mod_R$, where $R = k$ is a field.
        
        \item There is a category $\Top$ whose objects are topological spaces whose morphisms are continuous maps. It is sometimes useful to select a base point in a topological space, and there is indeed a category of based topological spaces $\Top_\ast$ with base point preserving continuous maps serving as the morphisms.
    \end{enumerate}
\end{ex}

\begin{defn}
    A category is \emph{discrete} if every arrow is an identity. 
\end{defn}

Every set $X$ can be considered as the set of objects of a discrete category, simply by assigning to each object $x \in X$ a morphism $x \to x$. If $\~D$ is a discrete category, then it is entirely determined by its set of objects, since any object $d$ of $\~D$ has only one possible morphism associated with it, namely $1_d: d \to d$. Because of this, we can view every discrete category as simply being a set, namely the set $\ob(\~D)$. Therefore, discrete categories are sets.

\section{Types of arrows in categories}

\subsection*{Inverses and isomorphisms}
\begin{defn}\label{defn: inverse_isom}
    Given an arrow $f: A \to B$ in a category $\~C$, an \emph{inverse} for $f$ is an arrow $g: B \to  A$ in $\~C$ such that $g \circ f = 1_A$ and $f \circ g = 1_B$. An arrow with an inverse is called an \emph{isomorphism}. If there is an isomorphism between $A$ and $B$, we say $A$ and $B$ are \emph{isomorphic}, denoted $A \cong B$. Arrows that are isomorphisms are often decorated with a tilde, as in $A \isom B$.
\end{defn}

Isomorphisms give us a notion of sameness between two objects in a category. Given an isomorphism $f: A \to B$, the inverse $g: B \to A$ allows us to pass freely between objects of $A$ and objects of $B$, simply by composing with $f$ on the right or left respectively.  

\begin{ex}
\label{ex:group}
A \emph{group} (in the sense of group theory) can be thought of as a category $\~G$ with one object, and where every arrow is an isomorphism. To see why this makes sense, denote by $G$ the unique object of $\~G$. Then $\~G$ contains the following data:
\begin{itemize}
    \item A set (or class) $\~G(G, G)$ of isomorphisms from $G$ to itself (the symmetries of the underlying object $G$),
    \item an associative composition map $\~G(G, G) \times \~G(G, G) \to \~G(G, G)$,
    \item and a two-sided unit $1_G: G \to G$.
\end{itemize}
Thus, the elements of $G$ are exactly the precisely the arrows in $\~G(G, G)$. Moreover, we have a associative binary operation, a unit $1_G$, and since every map in $\~G(G, G)$ is an isomorphism, each element of $g \in G$ has an inverse $h \in G$ such that $h \circ g = g \circ h = 1_G$. The category $\~G$ looks something like this:
\[\begin{tikzcd}
G \arrow[loop, distance=2em, in=35, out=325] \arrow[loop, distance=2em, in=215, out=145] \arrow[loop, distance=2em, in=305, out=235] \arrow[loop, distance=2em, in=125, out=55]
\end{tikzcd}\]
where each arrow $G \to G$ is a different element of the group $G$.
\end{ex}

We can also consider morphisms for which only one of the properties in definition \ref{defn: inverse_isom} holds.

\begin{defn}
    Let $f: A \to B$ be a morphism in a category $\~C$.
    \begin{itemize} 
        \item $f$ is called a \emph{section} or \emph{splitting} if there exists $B \To{g} A$ such that $g \circ f = 1_A$. That is, $f$ is a section if it is a right inverse for some morphism.
        \item $f$ is called a \emph{retraction} if there exists $B \To{g} A$ such that $f \circ g = 1_B$. That is, $f$ is a retraction if it is a left inverse for some morphism.
    \end{itemize}
\end{defn}

So, if we have a pair of arrows $r: A \to B$, $s: B \to A$ such that $r \circ s = 1_B$, as in the triangle
\[\begin{tikzcd}
    B & A \\
    & B
    \arrow["s", from=1-1, to=1-2]
    \arrow[Rightarrow, no head, from=1-1, to=2-2]
    \arrow["r", from=1-2, to=2-2]
\end{tikzcd}\]
then $s$ is a section of $r$ and $r$ is a retract of $s$.

\subsection*{Monics and epis}
\begin{defn}
Let $\~C$ be a category. 
\begin{itemize}
    \item A morphism $A \To{f} B$ in $\~C$ is called a \emph{monomorphism} if for all objects $X$ and all parallel arrows\begin{tikzcd}X \arrow[r, shift left]{}{x} \arrow[r, shift right, swap]{}{x'} & A\end{tikzcd}, we have
    \[f \circ x = f \circ x' \implies x = x'.\]
    We often write $f: A \mono B$ to signify that $f$ is a monomorphism.

    \item Similarly, $f$ is called an \emph{epimorphism} if for all objects $Y$ and all parallel arrows\begin{tikzcd}B \arrow[r, shift left]{}{y} \arrow[r, shift right, swap]{}{y'} & Y\end{tikzcd}, we have
    \[y \circ f = y' \circ f \implies y = y'.\] 
    We often write $f: A \epi B$ to signify that $f$ is a epimorphism.
\end{itemize}
\end{defn}

\begin{rmk}
    A monomorphism is often called a ``mono'' for short, or said to be \emph{monic}. Conversely, an epimorphism is sometimes called an ``epi'' for short or is described as being \emph{epic}.
\end{rmk}

Epimorphisms (resp. monomorphisms) are certainly related to surjective (resp. injective) maps, however, one should be careful not to think of the two as being ``the same''. Indeed, we will soon show that several key properties that are true of surjective and injective functions are not true of epimorphisms and monomorphisms. However, in the special context of $\Set$, the epimorphisms are exactly the surjective functions, and the monomorphism are indeed the injective functions. 

\begin{prop}
\label{prop:set_surjective_epi}
    If $f: X \to Y$ is an arrow in \textup{$\Set$}, then we have the following implications:
    \begin{enumerate}[label=(\roman*)]
        \item $f$ is surjective $\iff$ $f$ is an epimorphism
        \item $f$ is a injective $\iff$ $f$ is a monomorphism
    \end{enumerate}
\end{prop}

\begin{proof}
\hfill
\begin{enumerate}[label=(\roman*)]
    \item For the ``if'' direction, let $f: X \to Y$ be a surjective map in $\Set$, and suppose there exist parallel arrows \begin{tikzcd}Y \arrow[r, shift left]{}{g} \arrow[r, shift right, swap]{}{h} & Z\end{tikzcd} such that $g \circ f = h \circ f$. Let $y \in Y$. Since $f$ is surjective, there exists $x \in X$ such that $f(x) = y$, so we have
    \[g(f(x)) = h(f(x)) \implies g(y) = h(y)\]
    Since this is true for all $y \in Y$, we have $g = h$, so $f$ is an epimorphism. 

    For the reverse implication, suppose $f: X \to Y$ is an epimorphism in $\Set$. We need to find two arrows going out of $Y$ that will allow us to show surjectivity. Consider the maps $g_1, g_2: Y \to \{0, 1\}$ defined by
    \begin{align*}
        g_1: Y &\to \{0, 1\} \\
             y &\mapsto 1 \\
    \end{align*}
    and
    \[g_2(y) = \begin{cases}
                1 & y \in f(X) \\
                0 & \text{otherwise}
            \end{cases}\]
    Then for any $x \in X$ we see that 
    \[(g_1 \circ f)(x) = g_1(f(x)) = 1\]
    and since $f(x)$ is obviously in the image of $f$, we also have 
    \[(g_2 \circ f)(x) = g_2(f(x)) = 1\]
    so $g_1 \circ f = 1 = g_2 \circ f$, and since $f$ is an epi, we deduce that $g_1(y) = g_2(y)$ for all $y \in Y$. Since $g_1(y) = 1$ for all $y \in Y$, and $g_2(y) = 1$ only when $y \in f(X)$, equality of $g_1$ and $g_2$ implies that every $y \in Y$ belongs to $f(X)$, therefore $f$ is surjective.

    It is worth remarking that the choice of $g_2$ in the above proof, while seemingly magical, is actually a very natural map to consider. It is often referred to as the \emph{characteristic function} and is used in many proofs throughout mathematics, so it's an important map to keep in mind.

    \item First, let $f: X \to Y$ be injective, and suppose there exist parallel maps \begin{tikzcd}Z \arrow[r, shift left]{}{g} \arrow[r, shift right, swap]{}{h} & X \end{tikzcd} such that for all $z \in Z$ we have
    \[f(g(z)) = f(h(z))\]
    Then injectivity of $f$ allows us to immediately conclude that $g(z) = h(z)$, from which it follows that $g = h$, so $f$ is a monomorphism.

    To prove the converse, suppose $f$ is monic, and let $\{\ast\}$ be a one-element set. Let $x, x' \in X$, and define maps $g_1, g_2: \{\ast\} \to X$ by
    \[g_1(\ast) = x \,\text{ and }\, g_2(\ast) = x'\]
    Then $g_1(\ast) \neq g_2(\ast)$, and because $f$ is monic, we must have $f(g_1(\ast)) \neq f(g_2(\ast))$, i.e.,
    \[f(x) \neq f(x')\]
    Therefore $x \neq x' \implies f(x) \neq f(x')$, so $f$ is injective.
\end{enumerate}

\end{proof}
While the above proposition holds for sets, it is \emph{not} true in a general category. It may be tempting to think of epimorphisms and monomorphism as being ``the same'' as surjective and injective maps, however, this is not true in general. The next proposition illustrates that the best we can do in a general category is to relate sections (resp. retractions) to epimorphisms (resp. monomorphisms). 

\begin{thm}\label{prop:if_section_then_epi_etc}
    Let $\~C$ be a category and suppose $f: A \to B$ is a morphism in $\~C$. 
    \begin{enumerate}[label=(\roman*)]
        \item If $f$ is a section, then $f$ is a monomorphism. 
        \item If $f$ is a retraction, then $f$ is an epimorphism.
        \item If $f$ is both a section and a retraction, then $f$ is an isomorphism. 
    \end{enumerate}
\end{thm}

\begin{proof}
    \hfill
    \begin{enumerate}[label=(\roman*)]
        \item Let $f: A \to B$ be a section, and for some $X \in \~C$, suppose there exist morphisms \begin{tikzcd}
            X \arrow[shift left, r]{}{x} \arrow[shift right, swap, r]{}{x'} & A
        \end{tikzcd} such that $f \circ x = f \circ x'$. Since $f$ is a section, there exists $g: B \to A$ such that $g \circ f = 1_A$, so we have
            \begin{align*}
                (g \circ f) \circ x &= (g \circ f) \circ x \\
                1_A \circ x&= 1_A \circ x' \\
                \implies x &= x'
            \end{align*}
        so $f$ is a monomorphism.
        
        \item Let $f: A \to B$ be a retraction, and for some $Y \in \~C$, suppose there exist morphisms \begin{tikzcd}
            B \arrow[shift left, r]{}{y} \arrow[shift right, swap, r]{}{y'} & Y
        \end{tikzcd} such that $y \circ f = y' \circ f$. Since $f$ is a section, there exists $g: B \to A$ such that $f \circ g = 1_B$, so we have
            \begin{align*}
                y \circ (f \circ g) &= y' \circ (f \circ g) \\
                y \circ 1_B &= y' \circ 1_B \\
                \implies y &= y'
            \end{align*}
        and therefore $f$ is an epimorphism.

        \item Suppose that $f$ is both a section and a retraction. Then there exist $g: B \to A$ and $h: B \to A$ such that $g \circ f = 1_A$ and $f \circ h = 1_B$. Then we have
        \[(g \circ f) \circ h = 1_A \circ h = h\]
        and 
        \[g \circ (f \circ h) = g \circ 1_B = g\]
        So $h = g$, therefore $g \circ f = 1_A$ and $f \circ g = 1_B$, so $g$ is an inverse for $f$, hence $f$ is an isomorphism.
    \end{enumerate}
\end{proof}

In a general category, it is incredibly important to note that while we always have the implication
\[f \text{ is a section and a retraction} \implies f \text{ is an isomorphism}\]
a map that is both epic and monic need not be an isomorphism. The next example explores a category where the second implication does not hold.

\begin{ex}\label{ex:incl_of_Z_into_Q}
     The canonical inclusion $\iota$ of the integers into the rationals given by
    \begin{align*}
        \iota: \Z &\hookrightarrow \Q \\
        x &\mapsto x
    \end{align*}
    is a monomorphism \emph{and} an epimorphism in \textup{$\cRing$}, even though it is not an isomorphism.  It is easy to check that the former condition holds, so we will restrict ourselves to showing that it is an epimorphism. To that end, let $R$ be an object in $\cRing$, and suppose that there exist morphisms \begin{tikzcd} \Z \arrow[r, hook]{}{\iota} & \Q \arrow[r, shift left]{}{f} \arrow[r, shift right, swap]{}{f'} & R\end{tikzcd} such that $\iota f = \iota f'$. Then for any $k \in \Z$ we have 
    \[f(\iota(k)) = f'(\iota(k))\]
    which means that $f(k) = f'(k)$ for all $k \in \Z$. We now need to show this condition holds for all $q = \frac{a}{b} \in \Q$. So, take any $\frac{a}{b} \in \Q$. We'll now compute $f(\frac{a}{b})$
    \begin{align*}
        f\left(\frac{a}{b}\right) &= f\left(\frac{1}{b} \cdot a\right) \\
                                  &= f\left(\frac{1}{b}\right)f\left(a\right) \\
    \end{align*}
    Using the fact that $f(k) = f'(k)$ for all $k \in \Z$, we can write
    \begin{align*}
        f\left(\frac{a}{b}\right) &= f\left(ab^{-1}\right) \\
                                  &= f\left(a \right)f\left(b^{-1}\right) \\
                                  &= f\left(a\right)\left[f\left(b \right)\right]^{-1} \\
                                  &= f'\left(a\right)\left[f'\left(b \right)\right]^{-1} \\
                                  &= f'\left(ab^{-1}\right) \\
                                  &= f'\left(\frac{a}{b}\right)
    \end{align*}
    And so we see that $f(\frac{r}{m}) = f'(\frac{r}{m})$, so $f = f'$, and thus $\iota$ is an epimorphsim.
\end{ex}

One can use the proof strategy in the previous exercise to prove that the localization $S^{-1}R$ of $R$ by a multiplicative set $S$ obeys a similar property: namely, the natural inclusion
    \begin{align*}
        \iota: R &\hookrightarrow S^{-1}R\\
        r &\mapsto \frac{r}{1}
    \end{align*}
    is a monomorphism and an epimorphism in \textup{$\cRing$}.

While we need to be careful with regard to which implications are true or not in a general category, the category of sets has some particularly desirable properties, illustrated in the following proposition

\begin{prop}
    Given a morphism $f: X \to Y$ in \textup{$\Set$}, we have the following implications:
    \begin{enumerate}[label=(\roman*)]
        \item if $X$ is non-empty and $f$ is a monomorphism $\implies$ $f$ is a section
        \item $f$ is an epimorphism $\implies$ $f$ is a retraction
    \end{enumerate}
\end{prop}

\begin{proof}
    \hfill
    \begin{enumerate}[label=(\roman*)]
        \item Let $X$ be a non-empty set, and let $f$ be a monomorphism. Then $f$ is injective, so $f(x) = f(x') \implies x = x'$ for any $x, x' \in X$. Fix $x_0 \in X$, and define a map $g: Y \to X$ by
        \[g(y) = \begin{cases}
            x & \text{if } y = f(x) \text{ for some } x \in X \\
            x_0 & \text{otherwise}
        \end{cases}\]
        Note that we have used injectivity of $f$ implicitly in defining $g$, since otherwise we would not be able to identify a \emph{distinct} $x \in X$ for which $y = f(x)$. Now, take any $x \in X$. Computing $(g \circ f)(x)$, we get
        \begin{align*}
            (g \circ f)(x) &= g(f(x)) \\
                  &= x
        \end{align*}
        Therefore $g \circ f = 1_X$, so we see that $f$ is a section. 
        
        \item Suppose $f$ is an epimorphism. Then $f$ is surjective, so for every $y \in Y$ there exists $x \in X$ such that $f(x) = y$. We can therefore construct an inverse function $g$ as follows: for each $y \in Y$, choose $x \in X$ such that $y = f(x)$, and define $g(y) = g(f(x)) = x$. Then for any $y \in Y$, we have
        \begin{align*}
            (f \circ g)(y) &= f(g(y)) \\
                    &= f(g(f(x))) \\
                    &= f(x) \\
                    &= y
        \end{align*}
        Therefore $f \circ g = 1_Y$, so $f$ is a retraction. 
    \end{enumerate}
\end{proof}

The proof of (ii) is particularly interesting because $f$ might map multiple values in $X$ to the same $y \in Y$, so we need to \emph{choose} a particular $x$ from each fibre of $f$ such that $f(x) = y$. Since there may be uncountably many elements in $Y$, $f$ could very well have uncountably many fibres, meaning that for (ii) to hold, we must assume the axiom of choice. This gives us the following categorical definition of the axiom of choice.

\begin{defn}
    A category $\~C$ is said to \emph{satisfy the axiom of choice} if every epimorphism in $\~C$ is a retraction. 
\end{defn}

It is also important to note that while the above proposition holds in $\Set$, it is not true in a general category. For instance, the inclusion $\Z \hookrightarrow \Q$ that we studied in example \ref{ex:incl_of_Z_into_Q} is clearly a monomorphism, but is not a section -- there is certainly no map that can take \emph{every} $q \in \Q$ back to itself, so the map has no right inverse. As another example, in the interval category \textbf{I}:
    \[\mathbf{I} = \begin{boxedtikzcd} A \arrow[r, "f"] & B \end{boxedtikzcd}\]
$f$ is clearly an epimorphism, since the only other morphisms in \textbf{I} that are can compose with it are $1_A$ and $1_B$. However, it is also definitely not a retraction, since there is not even an arrow $B \to A$ in \textbf{I}. 

\section{Terminal and initial objects}

\begin{defn}
    \hfill
    \begin{itemize}
        \item An object $T$ in a category $\~C$ is called \emph{terminal} if for each $A$ in $\~C$, there is exactly one map $A \to T$.
        \item Similarly, an object $I \in \~C$ is called \emph{initial} if for each $A \in \~C$ there is exactly one map $I \to A$.
    \end{itemize}
    A terminal (resp. initial) object is called \emph{strict} if every map into (resp. out of) that object is an isomorphism. 
\end{defn}

\begin{ex}
\hfill
\begin{enumerate}[label=(\roman*)]

    \item $\varnothing$ is an initial object in $\Set$, since for any set $X$ there is exactly one map $\varnothing \to X$, which is the map that does nothing. As for terminal objects, any singleton set $\{\ast\}$ will be terminal in $\Set$, since for any set $X$, the only map $f: X \to \{\ast\}$ is $f(x) = \ast$ for all $x \in X$.
    
    \item It is easy to check that the ring $\{0\}$ is a terminal object in $\cRing$, since for any ring $R$ there is only one morphism $\vp: R \to \{0\}$, given by $\vp(r) = 0$ for all $r \in R$. Perhaps more surprisingly, $\Z$ is an initial object in $\cRing$. To see this, take any $k \in \Z$. Then for any ring $R$ we can define a homomorphism $\vp: \Z \to R$ by
        \[ \vp(k) = \begin{cases} 
          1_R + \cdots + 1_R \text{ ($k$ times)} & k > 0 \\
          0_R & k = 0 \\
          -1_R - \cdots - 1_R \text{ ($k$ times)} & k < 0
       \end{cases}\]
   To show this is the only map $\Z \to R$, suppose we had a different homomorphism $\psi: \Z \to R$. Since $\psi$ is a homomorphism of unital rings, $\psi(1)  = 1_R$, so given $k > 0$ we can write 
   \begin{align*}
       \psi(k) &= \psi(1 + 1 + \cdots + 1) & \text{($k$ times)} \\
               &= \psi(1) + \psi(1) + \cdots + \psi(1) & \text{(since $\psi$ is a ring homomorphism)} \\
               &= 1_R + 1_R + \cdots + 1_R & \text{($k$ times)} \\
               &= \vp(k)
   \end{align*}
   One can easily use the above technique to prove the case where $k < 0$, and we know $\psi(0) = 0_R$ since $\psi$ is a ring homomorphism. Therefore, $\vp = \psi$, so $\vp$ is indeed unique.

    \item The category $\Rel$ is interesting because it contains a special object that is both terminal and initial, namely the empty set, $\varnothing$. Recall that a morphism $A \To{R} B$ in $\Rel$ is a subset of $A \times B$. Thus, a morphism $A \to \varnothing$ is a subset of $A \times \varnothing = \varnothing$. But then the only subset of $A \times \varnothing$ is $\varnothing$, so there can be at most one map $A \to \varnothing$, therefore the empty set is terminal in this category. On the other hand, the same logic can easily be used to show $\varnothing$ is initial in $\Rel$. When an object $Z$ in a category $\~C$ is both initial and terminal, we call $Z$ a \emph{zero object} of $\~C$, and the category $\~C$ is sometimes said to be \emph{pointed.}

    \item In the category $\Cat$ of categories and functors, the one object category $\mathbf{1}$ is terminal: for any category $\~C$, there is precisely one functor $\~C \to \mathbf{1}$, given by mapping every object in $\~C$ to the unique object of $\mathbf{1}$ (and similarly for arrows). The initial object of $\Cat$ is the empty category $\mathbf{0}$, with no objects and no morphisms. 
\end{enumerate}
\end{ex}

Proving some basic results above terminal and initial objects is a good way to get a first look at the kind of arguments that are commonplace in category theory. The following propositions demonstrate an important property of terminal and initial objects, which is that they are determined up to a unique(!) isomorphism. We will come back to this property later when we discuss limits, which provide an alternative characterization of initial and terminal objects.

\begin{prop}
\label{prop:terminal_objects_are_isomorphic}
    Let $T$ and $T'$ be terminal objects in a category $\~C$. Then there is a unique isomorphism $T \to T'$.
\end{prop}

\begin{proof}
    Since $T'$ is terminal, there exists a unique map $f: T \to T'$. We must now produce $g: T' \to T$ such that $g \circ f = 1_T$ and $f \circ g = 1_{T'}$. Since $T'$ is terminal, $g$ will also be unique.
    
    Since $g \circ f: T \to T$ and $T$ is terminal, the map $g \circ f$ is unique. However, because $T$ is terminal, the only map $T \to T$ is $1_T$, so we must have 
    \[g \circ f = 1_T : T \to T\]
    Similarly, since $gf: T' \to T'$ and $T'$ is terminal, we deduce that 
    \[f \circ g = 1_{T'}: T' \to T'\]
    so $g$ is an inverse for $f$, and thus $f$ is an isomorphism.
\end{proof}

Using a similar proof strategy as above, one can also prove the following:
\begin{prop}
    Let $I$ and $I'$ be initial objects in a category $\~C$. Then there is a unique isomorphism $I \to I'$.
\end{prop}

\begin{prop}
    If $T$ is a terminal object in a category $\~C$ and $f: T \to T'$ is an isomorphism, then $T'$ is terminal in $\~C$.
\end{prop}

\begin{proof}
    Since $f$ is an isomorphism, there exists $g: T' \to T$ such that $g \circ f = 1_T$ and $f \circ g = 1_{T'}$. Since $T$ is terminal, $g$ is the only such map $T' \to T$. Let $X$ be an object in $\~C$. We must now show that there is only one arrow $X \to T'$. To that end, suppose we have two arrows\begin{tikzcd}X \arrow[r, shift left]{}{x} \arrow[r, shift right, swap]{}{x'} & T'\end{tikzcd}. Then certainly \begin{tikzcd}X \arrow[r, shift left]{}{g \circ x} \arrow[r, shift right, swap]{}{g \circ x'} & T\end{tikzcd}, but $T$ is terminal, so we have 
    \begin{align*}
        g \circ x &= g \circ x' \\
        (f \circ g) \circ x &= (f \circ g) \circ x' \\
        1_{T'} \circ x &= 1_{T'} \circ x'\\
        \implies x &= x'
    \end{align*}
    Thus, there is a unique arrow $X \to T'$ for all $X \in \~C$, so $T'$ is a terminal object.
\end{proof}

Again, we have a similar result for initial objects:

\begin{prop}
    If $I$ is a terminal object in a category $\~C$ and $f: T \to I'$ is an isomorphism, then $I'$ is terminal in $\~C$.
\end{prop}


\section{Functors}

\begin{defn}
    Let $\~A$ and $\~B$ be categories. A \emph{functor} $F: \~A \to \~B$ consists of
    \begin{itemize}
        \item A function
        \[\ob(\~A) \to \ob(\~B)\]
        written as $A \mapsto F(A)$.

        \item For each $A, A' \in \~A$, a function 
        \[\~A(A, A') \to \~B(F(A), F(A'))\]
        written as $f \mapsto F(f)$
    \end{itemize}
    These data must satisfy the following axioms
    \begin{itemize}
        \item $F(g \circ f) = F(g) \circ F(f)$ whenever $A \To{f} A' \To{g} A''$ in $\~A$
        \item $F(1_A) = 1_{F(A)}$ whenever $A \in \~A$
    \end{itemize}
\end{defn}
Given a functor $F$ and a map $g$, will almost always write $Fg = F(g)$.

With functors defined in this way, it is possible to construct a category $\Cat$, whose objects are \emph{small} categories and whose arrows are fuctors. The existence of such a category also tells us that functors can be composed. That is, given functors $F: \~A \to \~B$ and $G: \~B \to \~C$, we can construct a composite functor $G \circ F: \~A \to \~C$. This also tells us that for every category $\~A$, there is an identity functor $1_\~A : \~A \to \~A$.

\begin{ex}
\label{ex:forgetful_functors}
    Perhaps the easiest class of functors to describe are the so-called ``forgetful'' functors (this is an informal term with no precise definition). These are functors $F: \~A \to \~B$ that ``forget'' some of the structure of $\~A$. Here are a few examples
    \begin{enumerate}[label=(\roman*)]
        \item There is a functor $U: \Grp \to \Set$ defined as follows: if $G$ is a group then $U(G)$ is the underlying set of elements of $G$, and if $f: G \to H$ is a group homomorphism, $U(f): U(G) \to U(H)$ is just the underlying function from the set $U(G)$ to the set $U(H)$. So $U$ ``forgets'' a group's structure and forgets that group homomorphisms are homomorphisms.
        \item For $\Ring$ and $\Vect_k$ (where $k$ is a field), we have similar forgetful functors $\Ring \to \Set$ and $\Vect_k \to \Set$, which forget the ring and vector space structures, respectively.
        \item Forgetful functors don't have to forget everything. For instance, there is a functor $\Ring \to \Ab$ that forgets the multiplicative structure of a ring $R$, remembering only the underlying additive abelian group $(R, +)$ and its homomorphisms. Or, denote by $\Mon$ the category of monoids (i.e., groups without inverses). Then there is a functor $\Ring \to \Mon$ that takes a ring $R$ and forgets the additive structure, only remembering the underlying multiplicative monoid $(R, ~\cdot)$.
        \item There is an inclusion functor $U: \Ab \hookrightarrow \Grp$ defined by $U(A) = A$ for any abelian group $A$ and $U(f) = f$ for any homomorphism of abelian groups $f$. It forgets that abelian groups are abelian.
    \end{enumerate}
\end{ex}

\begin{ex}
\label{ex:free_functors}
    We'll now illustrate a type of functor that is often referred to as a \emph{free functor}. These are in some sense dual to forgetful functors, although they take a little more effort to construct. 
    \begin{enumerate}[label=(\roman*)]
        \item One of the simplest examples of a free functor is the construction of the free monoid on a set (in fact, whenever you see the word ``free'' in mathematics, there is a good chance that some kind of free functor is involved). Take any set $X$. For the sake of this example, let's say $X = \{a, b\}$. Then we can make \emph{words} from this set just by lining up elements of $X$ next to each other. For instance, $a$, $abaab$, and $bba$ are all completely valid words formed from $X$. The \emph{free monoid} on $X$ is the set
        \[F(X) := \{\text{words } x_1 x_2 \cdots x_k ~|~ x_i \in X\} \cup \{\varepsilon\}\]
        where $\varepsilon$ is the empty word (this will serve as the neutral element in our free monoid). If $\omega_0 = x_1 \cdots x_k$ and $w_1 = y_1 \cdots y_\ell$ are two words in $F(X)$, we can define a binary product $w_0 \ast w_1$ on $F(X)$ as follows
        \[w_0 \ast w_1 = x_1 \cdots x_k y_1 \cdots y_\ell\]
        
        Furthermore, if $M = F(X)$ and $N = F(Y)$ free monoids on the respective sets $X$ and $Y$, then a monoid morphism $f: M \to N$ is a map such that
        \[f(w_0 \ast_{M} w_1) = f(w_0) \ast_{N} f(w_1)\]
        for all words $w_0, w_1$ (where $\ast_{M}$ is the monoid product in $M$ and $\ast_N$ is the product in $N$) and
        \[f(\varepsilon) = \iota\]
        (where $\varepsilon$ is the neutral element in $M$ and $\iota$ is the neutral element in $N$).

        This free monoid construction defines a free functor $F: \Set \to \Mon$, where $F(X)$ is the free monoid on any set $X$, and a function $f: X \to Y$ between sets becomes a monoid morphism $F(f): F(X) \to F(Y)$. 

        \item Similarly to the construction above, we can take a set $X$ of elements and construct the \emph{free group} $F(X)$ on $X$, which indeed defines a functor $\Set \to \Grp$. The only difference between the free group and the free monoid is that, of course, we must account for inverses. 

        \item We can also construct a free commutative ring on any set $S$, therefore defining a functor $\Set \to \Ring$. In fact, this construction is quite familiar -- for each $s \in S$, we define an indeterminate $x_s$. The free commutative ring on $S$ is then defined as the ring of polynomials over $\Z$ in the commuting variables $x_s$, together with the usual addition and multiplication in polynomial rings. For instance, if $S$ is a two-element set, we have that $F(S) \cong \Z[x, y]$.

        \item We can also construct a free vector space quite easily. For a fixed field $k$, the free functor $F: \Set \to \Vect_k$ is defined by taking objects $F(S)$ to be vector spaces with bases in $F$. Roughly, given a set $S$, $F(S)$ is the set of all formal $k$-linear combinations of elements of $S$, i.e., expressions of the form
        \[\sum_{s \in S} \lambda_s s\]
        where each $\lambda_s$ is a scalar and there are only finitely many $s$ such that $\lambda_s \neq 0$ (we need this restriction because we can't take infinite sums in a vector space). We can add elements of $F(S)$
        \[\sum_{s \in S} \lambda_s s + \sum_{s \in S} \mu_s s = \sum_{s \in S}(\lambda_s + \mu_s)s\]
        and there is also scalar multiplication defined by 
        \[c \cdot \sum_{s \in S} \lambda_s s = \sum_{s \in S}(c \lambda_s)s\]
        for all $c \in k$. 

        Rings and vector spaces have the special property that it is relatively easy to write down an explicit formula for the free functor. The construction described  for monoids is much more typical.

    \end{enumerate}
\end{ex}

\begin{ex}
(For readers familiar with linear algebra) Let $V$ be a real vector space. The \emph{complexification} of $V$ is a vector space written $V_\C$ and defined as the tensor product of $V$ and $\C$:
        \[V_\C = V \otimes_\R \C\]
        Note that since we are tensoring over $\R$, it is not immediate that $V_\C$ is a vector space over $\C$. However, we can easily make $V_\C$ a vector space over $\C$ by defining an operation of scalar multiplication from $\C$ as follows
        \[\alpha(v \otimes \beta) = v \otimes (\alpha \beta) \quad \text{for all } v \in V \text{ and } \alpha, \beta \in \C\]
        Moreover, given a linear map $f: V \to W$, we may define a map $f_\C: V_\C \to W_\C$ by
        \[f_\C(v \otimes z) = f(v) \otimes z \quad \text{for all } v \in V \text{ and } z \in \C\]
        And from this we can easily verify that the complexification defines a functor
        \[(-)_\C: \Vect_\R \to \Vect_\C\]
        from the category of real vector spaces to the category of complex vector spaces.
\end{ex}

\begin{ex}
    Let $G$ and $H$ be groups (or monoids, if you prefer) regarded as one-object categories (i.e., groupoids) $\~G$ and $\~H$. A functor $F: \~G \to \~H$ must send the unique object of $\~G$ to the unique object of $\~H$ (i.e., $F(G) = G$ and $F(H) = H$), so it is completely determined by its effect on maps. Hence, a functor $F: \~G \to \~H$ is just a function $F: G \to H$ satisfying the additional functionality axioms, namely $F(g \circ g') = F(g) \circ F(g')$ for all $g, g' \in G$ and $F(1_G) = 1_H$. In other words, such a functor is just a group homomorphism $G \to H$.
\end{ex}

\begin{ex}
    Let $G$ be a group, regarded as a one-object category $\~G$, as in Example \ref{ex:group}. A functor $F: \~G \to \Set$ consists of a set $F(G) = S$ (the unique value of $F$ at $G$), together with a map of sets $F(g): S \to S$ for each $g \in \~G(G, G)$ (where recall that the arrows in $\~G(G, G)$ correspond to the elements of the group $G$). If we write $(F(g))(s) = g \cdot s$ for each $s \in S$, we see that the functor $F$ amounts to a set $S$ together with a function
    \begin{align*}
        G \times S &\to S \\
        (g, s) &\mapsto g \cdot s
    \end{align*}
    and functoriality tells us that $(F(g\circ g'))(s) = (F(g) \circ F(g'))(s)$ and $F(1_G) = 1_S$, i.e., $(gg') \cdot s = g \cdot (g' \cdot s)$ and $1 \cdot s = s$ for all $g, g' \in \~G(G,G)$ and $s \in S$. Therefore, a functor $F: \~G \to \Set$ is indeed a \emph{left $G$-set}: a set equipped with a left action from $G$.
\end{ex}

\begin{ex}
    Any system of polynomial equations like
    \begin{align}
        2x^2 + y^2 - 3z^2 &= 1 \\
        x^3 + x &= y^2
    \end{align}
    gives rise to a functor $\cRing \to \Set$. To see this, let $R$ be a commutative ring, and let $F(R)$ be the set of triples $(x, y, z) \in R \times R \times R$ satisfying equations (1.2) and (1.3). If $\vp: R \to S$ is a ring homomorphism and $(x, y, z) \in F(R)$, we have that the triple $(\vp(x), \vp(y), \vp(z))$ is in $F(S)$, so the map of rings $\vp: R \to S$ induces a map of sets $F(\vp): F(R) \to F(S)$. This defines a functor $F: \cRing \to \Set$. 
    
    In algebraic geometry, a \emph{scheme} is a functor $\cRing \to \Set$ with certain additional properties (this is not the most common way of defining schemes, but it is equivalent). The functor above is an easy example.
\end{ex}

\begin{ex}
    When $A$ and $B$ are (pre)ordered sets, a functor $F: A \to B$ between the corresponding categories is exactly an \emph{order-preserving map}, that is, a function $f: A \to B$ such that $a \leq a' \implies f(a) \leq f(a')$.
\end{ex}

\begin{defn}
    Let $\~A$ and $\~B$ be categories. A \emph{contravariant} functor from $\~A$ to $\~B$ is a functor $\~A^\op \to \~B$. 
\end{defn}

Contravariant functors are also called ``arrow-reversing'' because an arrow $F(A) \to F(A')$ in $\~B$ corresponds to an arrow $A \leftarrow A'$ in $\~B$. To avoid confusion, we will write ``contravariant functor from $\~A$ to $\~B$'' instead of ``contravariant functor $\~A \to \~B$''. Some authors call ordinary functors $\~A \to \~B$ \emph{covariant functors} for extra emphasis. One of the first places where one encounters contravariant functors is in the context of linear algebra:

\begin{ex}
\label{ex:double_dual}
    Fix a field $k$. For any two vector spaces $V$ and $W$ over $k$, there is a vector space
\[\Hom(V, W) = \{\text{linear maps } V \to W\}\]
where addition and scalar multiplication are defined \emph{pointwise}, i.e., for all $v \in V$ we have
\[(f + g)(v) = f(v) + g(v) \quad \text{and} \quad (cf)(v) = c\cdot f(v)\]
for all $f, g \in \Hom(V, W)$ and all $c \in k$. Let $W$ be any vector space. Given a linear map $q: V' \to W$, we can choose a map $f: V \to V'$ and compose it with $q$ to obtain a map $q \circ f: V \to W$:
\[V \To{f} V' \To{q} W\] 
This process induces a map
\begin{align*}
f^\ast : \Hom(V', W) &\to \Hom(V, W)\\
                   q &\mapsto q \circ f 
\end{align*}
and this in turn defines a functor
\[\Hom(-, W): \Vect_k^\op \to \Vect_k\]
This notation means that that $\Hom(-, W)$ is a functor whose value at a $k$-vector space $V$ is $\Hom(V, W)$. A particularly well-known case of this functor is when $W = k$, viewed as a one-dimensional vector space over itself. The vector space $\Hom(V, k)$ is known as the \emph{dual space} of $V$, and denoted by $V^\ast$. Thus, there is a functor
\begin{align*}
(~)^\ast = \Hom(-, k): \Vect_k^\op &\to \Vect_k \\
                        V &\mapsto V^\ast
\end{align*}
sending each vector space $V$ to its dual.
\end{ex} 
\begin{ex}
    Another well-known example of contravariant functors arises in topology. Given a topological space $X$, let $C(X)$ be the ring of continuous real-valued functions on $X$, as in the previous example, the ring operations are defined pointwise. So, for example, given continuous maps $p_1, p_2 : X \to \R$, there is a continuous map $p_1 + p_2: X \to \R$ defined by $(p_1 + p_2)(x) = p_1(x) + p_2(x)$. Moreover, given another topological space $Y$ and a continuous map $f: X \to Y$, there is a ring homomorphism $C(f): C(Y) \to C(X)$, where for each $q \in C(Y)$ we define $(C(f))(q)$ to be the composite map
    \[X \To{f}Y \To{q} \R\]   
    Observe that $C(f)$ goes in the opposite direction from $f$. We can therefore check that $C$ defines a functor 
    \[C: \Top^\op \to \Ring\]
    i.e., a contravariant functor from $\Top$ to $\Ring$. 
\end{ex}

Contravariant functors into $\Set$ are so important that we give them their own name
\begin{defn}
\label{defn:presheaf}
    Let $\~C$ be a category. A \emph{presheaf} on $\~C$ is a functor $F: \~C^\op \to \Set$.
\end{defn}

\begin{ex}
    Let $X$ be a topological space, and denote by $\~O(X)$ the poset of open subsets of $X$, ordered by inclusion. View this poset as a category. Thus, objects of $\~O(X)$ are the open subsets of $X$, and for $U, U' \in \~O(X)$, there is one map $U \to U'$ if $U \subseteq U'$, and none otherwise. A \emph{presheaf} on the space $X$ is a presheaf on the category $\~O(X)$. For example, a presheaf $F$ can be defined on $X$ by
    \[F(U) = \{\text{continuous functions } U \to \R\}\]
    If $U \subseteq U'$, then every continuous function defined on $U'$ is also defined on $U$, but there may be continuous functions defined on $U$ which are not defined on $U'$. Thus, we have the reverse inclusion $F(U) \supseteq F(U')$, so we can simply take the map $F(U') \to F(U)$ to be the restriction. A special class of presheaves called \emph{sheaves} are one of the central objects of study in algebraic geometry. Indeed, the presheaf $F$ in that we have just described is also a sheaf.
\end{ex}

\begin{defn}
    A functor $F: \~A \to \~B$ is called \emph{faithful} (resp. \emph{full}) if for all $A, A' \in \~A$, the function
    \begin{align*}
        \~A(A, A') &\to \~B(F(A), F(A')) \\
        f &\mapsto F(f)
    \end{align*}
    is injective (resp. surjective).
\end{defn}

\begin{rmk}
    It might not be obvious from this definition, but faithfulness does \emph{not} say that if $f_1$ and $f_2$ are distinct maps in $\~A$ then $F(f_1) \neq F(f_2)$. In the situation in Figure \ref{fig:full_faithful}, $F$ is faithful if for each $A, A'$ and $g$ as shown, there is at most one dotted arrow that $F$ sends to $g$. $F$ is full if for each $A, A'$ and $g$, there is at least one dotted arrow that $F$ sends to $g$.
\end{rmk}

\begin{figure}[h]
    \centering
        \[
            \underset{\mbox{\large $\~A$}}{
                \framebox[1.75in]{
                    \begin{tikzcd}
                    A \\
                    \\
                    {A'}
                    \arrow[dashed, from=1-1, to=3-1]
                    \end{tikzcd}                   
                }
            }
            \quadd \begin{tikzcd}
            {} & {}
            \arrow["F", from=1-1, to=1-2]
            \end{tikzcd} \quadd
            \underset{\mbox{\large $\~B$}}{
                \framebox[1.75in]{
                    \begin{tikzcd}
                    F(A) \\
                    \\
                    {F(A')}
                    \arrow[from=1-1, to=3-1]
                    \end{tikzcd}                   
                }
            }
        \]
    \caption{Full and faithful functors}
    \label{fig:full_faithful}
\end{figure}

\begin{defn}
    Let $\~C$ be a category. A \emph{subcategory} $\~S$ of $\~C$ consists of a subclass $\ob(\~S)$ of $\ob(\~C)$ together with a subclass $\Hom_\~S(S, S')$ of $\Hom_\~C(S, S')$ for all $S, S' \in \ob(\~S)$, such that $\~S$ is also a category -- that is, it is closed under composition and identities. $\~S$ is a \emph{full subcategory} if $\Hom_\~S(S, S') = \Hom_\~C(S, S')$ for all $S, S' \in \ob(\~S)$.
\end{defn}

\section{Natural transformations}

We have seen categories, and we have also seen functors, or arrows between categories. Perhaps surprisingly, in category theory, there is also a notion of arrows between functors. This notion applies when two functors have the same domain and codomain
\[\begin{tikzcd}
	{\~A} & {\~B}
	\arrow["F", shift left=1, from=1-1, to=1-2]
	\arrow["G"', shift right=1, from=1-1, to=1-2]
\end{tikzcd}\]
These arrows are called natural transformations, and to see how they might work, first, consider the simple case of the discrete category $\~A$ whose objects are natural numbers (recall that this means the only morphisms in $\~A$ are identities). A functor $F$ from $\~A$ to another category $\~B$ is simply a sequence $(F_0, F_1, F_2, \cdots)$ of objects in $\~B$. If we have another functor $G: \~A \to \~B$ with a corresponding sequence of objects $(G_0, G_1, G_2, \cdots)$ It would be reasonable to define a ``map from $F$ to $G$'' as a sequence of maps
\[\begin{tikzcd}
	{F_0} & {F_1} & {F_2} \\
	{G_0} & {G_1} & {G_2}
	\arrow["{\eta_0}", from=1-1, to=2-1]
	\arrow["{\eta_1}", from=1-2, to=2-2]
	\arrow["{\eta_2}", from=1-3, to=2-3]
\end{tikzcd} \quad \cdots \]
one for each natural number in $\~A$. This suggests that in the general case, a natural transformation between functors \begin{tikzcd}
{\~A} & {\~B}
\arrow["F", shift left=1, from=1-1, to=1-2]
\arrow["G"', shift right=1, from=1-1, to=1-2]
\end{tikzcd} should consist of maps $\eta_A: F(A) \to G(A)$, one for each $A \in \~A$. When we formally define natural transformations, we will also require some compatibility between the maps $\eta_A$ and arrows in the category $\~A$.

\begin{defn}
    Let $\~A$ and $\~B$ be categories and let \begin{tikzcd} {\~A} & {\~B} \arrow["F", shift left=1, from=1-1, to=1-2] \arrow["G"', shift right=1, from=1-1, to=1-2]\end{tikzcd} be functors. A \emph{natural transformation} $\alpha: F \to G$ is a family $\left(F(A) \To{\alpha_A} G(A)\right)_{A \in \~A}$ of maps in $\~B$ such that for every map $A \to A'$ in $\~A$, the square
        \[\begin{tikzcd}
    	{F(A)} & {F(A')} \\
    	{G(A)} & {G(A')}
    	\arrow["{\alpha_A}"', from=1-1, to=2-1]
    	\arrow["{\alpha_{A'}}", from=1-2, to=2-2]
    	\arrow["{F(f)}", from=1-1, to=1-2]
    	\arrow["{G(f)}"', from=2-1, to=2-2]
        \end{tikzcd}\]
    commutes, i.e. 
        \[F(f) \alpha_{A'} = \alpha_A G(f)\]
    The maps $\alpha_A$ are called the \emph{components} of $\alpha$. We write 
        \[\begin{tikzcd}
    	{\~A} \arrow[r, "F", bend left, ""{name=U, below}]
    	       \arrow[r, "G"', bend right, ""{name=D}]
            & {\~B}
    	   \arrow["~\alpha", Rightarrow, from=U, to=D]
    \end{tikzcd}\]
    to signify that $\alpha$ is a natural transformation from $F$ to $G$. 
\end{defn}

So, in our discrete category example above, a natural transformation $\alpha: F \to G$ is simply a family $\left(F(A) \To{\alpha_A} G(A)\right)_{A \in \~A}$ of maps in $\~B$. Indeed, this principle holds when $\~A$ is any discrete category, not just when $\ob(\~A) = \N$ like our example above.

\begin{ex}
    In this example, we'll show how the determinant of an $n \times n$ matrix can be interpreted as a natural transformation. Fix a natural number $n$. For any commutative ring $R$, the $n \times n$ matrices with entries in $R$ form a monoid $M_n(R)$ under matrix multiplication. Moreover, if we have a ring homomorphism $\vp: R \to S$, this induces a monoid homomorphism $M_n(R) \to M_n(S)$. For example, given a $2 \times 2$ matrix
    \[\begin{pmatrix}
        r_{11} & r_{12} \\
        r_{21} & r_{22}
    \end{pmatrix}\]
    in $M_2(R)$, we can define a matrix in $M_2(S)$ by 
    \[\begin{pmatrix}
        \vp(r_{11}) & \vp(r_{12}) \\
        \vp(r_{21}) & \vp(r_{22})
    \end{pmatrix}\]
    and since $\vp$ is a homomorphism of the underlying rings, this induces a monoid homomorphism $M_2(R) \to M_2(S)$. So we have now defined objects in $\Mon$ from objects in $\cRing$, and given a ring homomorphism $R \to S$, we can find a monoid homomorphism $M_n(R) \to M_n(S)$. This gives us a functor $M_n: \cRing \to \Mon$. Also, the elements of any ring $R$ form a monoid under multiplication, giving another functor $U: \cRing \to \Mon$. Every $n \times n$ matrix $X$ over a commutative ring $R$ has a determinant $\det_R(X)$, which is an element of $R$. Furthermore, thanks to familiar properties of the determinant 
        \[\mathrm{det}_R (XY) = \mathrm{det}_R(X)\mathrm{det}_R(Y), \quad \mathrm{det}_R(I) = 1\]
    it is easy to see that for each $R$, the function $\det_R: M_n(R) \to U(R)$ is a monoid homomorphism. So we have a family of maps 
    \[\left(F(R) \To{\det_R} G(R)\right)_{R \in \cRing}\]
    and we can ask whether they define a natural transformation 
\[\begin{tikzcd}
    \cRing & \Mon
    \arrow[""{name=0, anchor=center, inner sep=0}, "{M_n}", shift left, curve={height=-6pt}, from=1-1, to=1-2]
    \arrow[""{name=1, anchor=center, inner sep=0}, "U"', shift right, curve={height=6pt}, from=1-1, to=1-2]
    \arrow["\,\det", shift right, shorten <=6pt, shorten >=6pt, Rightarrow, from=0, to=1]
\end{tikzcd}\]
Indeed, they do. The fact that the naturality square commutes reflects the fact that $\det_R$ is defined in the same way no matter what $R$ is.
\end{ex}
\subsection*{Composition laws for natural transformations}
One straightforward question we can ask ourselves about natural transformations is whether they can be composed, and indeed, they can. Given categories $\~A$ and $\~B$ and natural transformations
\begin{equation}
\begin{tikzcd}
    {\~A} && {\~B}
    \arrow[""{name=0, anchor=center, inner sep=0}, "F", shift left, curve={height=-12pt}, from=1-1, to=1-3]
    \arrow[""{name=1, anchor=center, inner sep=0}, "H"', shift right, curve={height=12pt}, from=1-1, to=1-3]
    \arrow[""{name=2, anchor=center, inner sep=0}, "G"{description}, from=1-1, to=1-3]
    \arrow["\,\alpha", shorten <=3pt, shorten >=3pt, Rightarrow, from=0, to=2]
    \arrow["\,\beta", shorten <=3pt, shorten >=3pt, Rightarrow, from=2, to=1]
\end{tikzcd}
\end{equation}
there is a composite natural transformation
\[\begin{tikzcd}
    {\~A} & {\~B}
    \arrow[""{name=0, anchor=center, inner sep=0}, "G"', curve={height=12pt}, from=1-1, to=1-2]
    \arrow[""{name=1, anchor=center, inner sep=0}, "F", curve={height=-12pt}, from=1-1, to=1-2]
    \arrow["{\alpha \beta}", shorten <=4pt, shorten >=4pt, Rightarrow, from=1, to=0]
\end{tikzcd}\]
defined by $(\alpha \beta)_A = \alpha_A \beta_A$ for all $A \in \~A$. There is also an identity natural transformation
\[\begin{tikzcd}
    {\~A} & {\~A}
    \arrow[""{name=0, anchor=center, inner sep=0}, "F"', curve={height=12pt}, from=1-1, to=1-2]
    \arrow[""{name=1, anchor=center, inner sep=0}, "F", curve={height=-12pt}, from=1-1, to=1-2]
    \arrow["{1_F}", shorten <=4pt, shorten >=4pt, Rightarrow, from=1, to=0]
\end{tikzcd}\]
on any functor $F$, defined by $(1_F)_A = 1_{F(A)}$. So for any two categories $\~A$ and $\~B$, there is a category whose objects are are the functors from $\~A$ to $\~B$ and whose arrows are the natural transformations between those functors. This is called the \emph{functor category} from $\~A$ to $\~B$, written as $\Fun(\~A, \~B)$ or $[\~A, \~B]$ or $\~B^{\~A}$. Its underlying set of objects is equal to $\Cat(\~A, \~B)$. 

\begin{ex}
    Recall that a presheaf on a category $\~C$ is a functor $F: \~C^\op \to \Set$. Given two presheaves $F, G: \~C^\op \to \Set$ on the same category $\~C$, we can reasonably define a \emph{morphism of presheaves} to be a natural transformation \[\begin{tikzcd}[column sep = scriptsize]
        {\~C} & \Set
        \arrow[""{name=0, anchor=center, inner sep=0}, "F", bend left=50, from=1-1, to=1-2]
        \arrow[""{name=1, anchor=center, inner sep=0}, "G"', bend right=50, from=1-1, to=1-2]
        \arrow["\alpha"', shorten <=4pt, shorten >=4pt, Rightarrow, from=0, to=1]
    \end{tikzcd}\]
    In this way, presheaves on $\~C$ form a category, called the \emph{presheaf category} on $\~C$ and denoted $\pre{\~C}$. It is equal to the functor category $\Fun(\~C^\op, \Set)$.  
\end{ex}

The type of composition described above is called \emph{vertical composition}. There is also a \emph{horizontal composition}: given natural transformations
\[\begin{tikzcd}
    {\~A} & {\~B} & {\~C}
    \arrow[""{name=0, anchor=center, inner sep=0}, "F", curve={height=-12pt}, from=1-1, to=1-2]
    \arrow[""{name=1, anchor=center, inner sep=0}, "G"', curve={height=12pt}, from=1-1, to=1-2]
    \arrow[""{name=2, anchor=center, inner sep=0}, "{F'}", curve={height=-12pt}, from=1-2, to=1-3]
    \arrow[""{name=3, anchor=center, inner sep=0}, "{G'}"', curve={height=12pt}, from=1-2, to=1-3]
    \arrow["\alpha", shorten <=3pt, shorten >=3pt, Rightarrow, from=0, to=1]
    \arrow["{\alpha'}", shorten <=3pt, shorten >=3pt, Rightarrow, from=2, to=3]
\end{tikzcd}\]
the horizontal composition produces a natural transformation
\[\begin{tikzcd}
    {\~A} & {\~C}
    \arrow[""{name=0, anchor=center, inner sep=0}, "{F F'}", curve={height=-12pt}, from=1-1, to=1-2]
    \arrow[""{name=1, anchor=center, inner sep=0}, "{G G'}"', curve={height=12pt}, from=1-1, to=1-2]
    \arrow[shorten <=3pt, shorten >=3pt, Rightarrow, from=0, to=1]
\end{tikzcd}\]
which is traditionally written as $\alpha \ast \alpha'$. The component of $\alpha \ast \alpha'$ at $A \in \~A$ is defined to be the diagonal of the naturality square
\[\begin{tikzcd}
    {F'(F(A))} & {F'(G(A))} \\
    {G'(F(A))} & {G'(G(A))}
    \arrow["{F'(\alpha_A)}", from=1-1, to=1-2]
    \arrow["{\alpha'_{F(A)}}", from=1-1, to=2-1]
    \arrow["{\alpha'_{G(A)}}", from=1-2, to=2-2]
    \arrow["{G'(\alpha_A)}"', from=2-1, to=2-2]
\end{tikzcd}\]
This means that the component $(\alpha \ast \alpha')_A$ can be defined as either $F'(\alpha_A) \alpha'_{G(A)}$ or $\alpha'_{F(A)} G'(\alpha_A)$. It makes no difference which, since they are equal to each other. A particularly important case is when one of $\alpha$ or $\alpha'$ is the identity transformation. In this case
\[\begin{tikzcd}
    {\~A} & {\~B} & {\~C}
    \arrow["F", from=1-1, to=1-2]
    \arrow[""{name=0, anchor=center, inner sep=0}, "{G'}"', curve={height=12pt}, from=1-2, to=1-3]
    \arrow[""{name=1, anchor=center, inner sep=0}, "{F'}", curve={height=-12pt}, from=1-2, to=1-3]
    \arrow["{\alpha'}", shorten <=3pt, shorten >=3pt, Rightarrow, from=1, to=0]
\end{tikzcd} \quad \text{gives rise to} \quad \begin{tikzcd}
    {\~A} & {\~C}
    \arrow[""{name=0, anchor=center, inner sep=0}, "{FF'}", curve={height=-12pt}, from=1-1, to=1-2]
    \arrow[""{name=1, anchor=center, inner sep=0}, "{FG'}"', curve={height=12pt}, from=1-1, to=1-2]
    \arrow["{F\alpha'}", shift right, shorten <=3pt, shorten >=3pt, Rightarrow, from=0, to=1]
\end{tikzcd}\]
where $(F \alpha')_A = \alpha'_{F(A)}$, and
\[\begin{tikzcd}
    {\~A} & {\~B} & {\~C}
    \arrow[""{name=0, anchor=center, inner sep=0}, "F", curve={height=-12pt}, from=1-1, to=1-2]
    \arrow[""{name=1, anchor=center, inner sep=0}, "G"', curve={height=12pt}, from=1-1, to=1-2]
    \arrow["{F'}", from=1-2, to=1-3]
    \arrow["\alpha", shorten <=3pt, shorten >=3pt, Rightarrow, from=0, to=1]
\end{tikzcd} \quad \text{gives rise to} \quad \begin{tikzcd}
    {\~A} & {\~C}
    \arrow[""{name=0, anchor=center, inner sep=0}, "{FF'}", curve={height=-12pt}, from=1-1, to=1-2]
    \arrow[""{name=1, anchor=center, inner sep=0}, "{GF'}"', curve={height=12pt}, from=1-1, to=1-2]
    \arrow["{\alpha F'}", shift right, shorten <=3pt, shorten >=3pt, Rightarrow, from=0, to=1]
\end{tikzcd}\]
where $(\alpha F')_A = F'(\alpha_A)$. Some authors refer to this special type of horizontal composition as \emph{whiskering} (although we will not adopt this convention in this text). 

Vertical and horizontal composition play well with one another: natural transformations
\[\begin{tikzcd}
    {\~A} & {\~B} & {\~C}
    \arrow[""{name=0, anchor=center, inner sep=0}, "G"{description}, from=1-1, to=1-2]
    \arrow[""{name=1, anchor=center, inner sep=0}, "H"', curve={height=25pt}, from=1-1, to=1-2]
    \arrow[""{name=2, anchor=center, inner sep=0}, "F", curve={height=-25pt}, from=1-1, to=1-2]
    \arrow[""{name=3, anchor=center, inner sep=0}, "{H'}"', curve={height=25pt}, from=1-2, to=1-3]
    \arrow[""{name=4, anchor=center, inner sep=0}, "{G'}"{description}, from=1-2, to=1-3]
    \arrow[""{name=5, anchor=center, inner sep=0}, "{F'}", curve={height=-25pt}, from=1-2, to=1-3]
    \arrow["\beta", shorten <=4pt, shorten >=4pt, Rightarrow, from=0, to=1]
    \arrow["\alpha", shorten <=4pt, shorten >=4pt, Rightarrow, from=2, to=0]
    \arrow["{\beta'}", shorten <=4pt, shorten >=4pt, Rightarrow, from=4, to=3]
    \arrow["{\alpha'}", shorten <=4pt, shorten >=4pt, Rightarrow, from=5, to=4]
\end{tikzcd}\]
obey the \emph{interchange law}:
\begin{equation}
(\alpha \beta) \ast (\alpha' \beta') = (\alpha \ast \alpha') (\beta \ast \beta'): F F' \to H H' \quad \text{and} \quad 1_F \ast 1_{F'} = 1_{F F'}
\end{equation}

Whenever we study natural transformations, we are implicitly working with \emph{$2$-categories}; categories with an additional notion of ``morphisms between morphisms''. The category $\Cat$ of all small categories is the prototypical example of a $2$-category: there are the regular morphisms between objects in $\Cat$ (``1-morphisms''), which are given by functors, and there are also arrows between these 1-morphisms (``2-morphisms''), which are given by natural transformations. Since natural transformations are one of the basic building blocks of category theory, we will eventually need to study 2-categories in greater detail. However, once we discuss 2-categories, it becomes tempting to discuss 3-categories, 4-categories, and so on, until eventually we are discussing $\infty$-categories. We will not wade into the waters of $\infty$-category theory in these notes, however, we will explore 2-categories in some detail in a later chapter. 

\begin{ex}
    Let $\mathbf{2}$ be the discrete category with two objects. A functor from $\mathbf{2}$ to another category $\~B$ is a pair of objects in $\~B$, and a natural transformation is a pair of maps. The functor category $\Fun(\mathbf{2}, \~B)$ is therefore isomorphic to the product category $\~B \times \~B$, which fits well with the alternative notion $\~B^\mathbf{2}$ for the functor category.
\end{ex}

\begin{ex}[For readers familiar with representation theory]
Let $G$ be a group, interpreted as a one-object category $\~G$, and consider the category $\Vect_k$ of vector spaces over a field $k$. A functor $\rho: \~G \to \Vect_k$ consists of a single vector space $V$ (the value of $\rho$ at the unique object $G$ of $\~G$), together with a $k$-linear map $\rho_g := \rho(g): V \to V$ for each map $g: G \to G$. Since maps $g: G \to G$ represent elements of the group $G$, a functor $\rho: \~G \to \Vect_k$ is an assignment of each element $g \in G$ to a $k$-linear map in $\Hom(V, V) = GL(V)$ which respects composition (the group operation) and identity. In other words, such a functor is simply a group homomorphism $G \to GL(V)$, i.e., a representation of $G$! The category $\Fun(\~G, \Vect_k)$ of all functors $\~G \to \Vect_k$ therefore describes all possible representations of $G$. If $\rho, \sigma \in \Fun(\~G, \Vect_k)$ are representations of $G$ with $\rho(G) = V$ and $\sigma(G) = W$, a morphism from $\rho$ to $\sigma$ is simply a natural transformation between the associated functors $\rho$ and $\sigma$. This consists of a single $k$-linear map $T: \rho(G) = V \to W = \sigma(G)$ such that for every map $g: G \to G$ (i.e., every element of the group $G$), the diagram
\[\begin{tikzcd}
    {V} & {V} \\
    {W} & {W}
    \arrow["{T}"', from=1-1, to=2-1]
    \arrow["{T}", from=1-2, to=2-2]
    \arrow["{\rho_g}", from=1-1, to=1-2]
    \arrow["{\sigma_g}"', from=2-1, to=2-2]
\end{tikzcd}\]
is commutative. Thus, an element of $\Fun(\~G, \Vect_k)$ is a representation of $G$, and a morphism in this category is precisely a $G$-equivariant map. Given a group $G$ and a vector space $V$ over the field $k$, $\Rep(G)$ denotes the category whose objects are $k$-linear representations of $G$, and whose arrows are $G$-equivariant maps. This is equivalent to the functor category $\Fun(\~G, \Vect_k)$ described above. When $G$ is a Lie group, we require that all the representations be smooth.
\end{ex}

\begin{defn}
    Let $\~A$ and $\~B$ be categories. A \emph{natural isomorphism} between $\~A$ and $\~B$ is an isomorphism in $\Fun(\~A, \~B)$. 
\end{defn}

The following equivalent form of the definition is also useful

\begin{lem}
\label{lem:nat_isom_condition}
    Let \begin{tikzcd}
        {\~A} \arrow[r, "F", bend left, ""{name=U, below}]
               \arrow[r, "G"', bend right, ""{name=D}]
            & {\~B}
           \arrow["~\alpha", Rightarrow, from=U, to=D]
    \end{tikzcd} be a natural transformation. Then $\alpha$ is a natural isomorphism if and only if $\alpha_A: F(A) \to G(A)$ is an isomorphism for all $A \in \~A$.
\end{lem}

\begin{defn}
\label{defn:naturally}
    Given functors \begin{tikzcd}
    {\~A} & {\~B}
    \arrow["F", shift left=1, from=1-1, to=1-2]
    \arrow["G"', shift right=1, from=1-1, to=1-2]
\end{tikzcd}, we say that
\[F(A) \cong G(A) \emph{ naturally in } A\]
if $F$ and $G$ are naturally isomorphic.
\end{defn}

The motivation for category theory was not the study of categories nor functors, but rather of natural transformations. As described in their foundational paper \emph{General Theory of Natural Equivalences} \cite{eilenberg_maclane}, Samuel Eilenberg and Saunders Mac Lane founded category theory as a means of systematically studying the notion of an isomorphism being ``natural'' or ``independent of choice''. Such isomorphisms have long been ubiquitous in mathematics, however, prior to the advent of category theory, mathematicians had no means of formally defining the concept of naturality. This idea is best illustrated through examples:

\begin{ex}
    Our first example of a natural isomorphism is the first one considered by Eilenberg and Mac Lane in \cite{eilenberg_maclane}
    Let $\fdVect_k$ be the category of finite-dimensional vector spaces over the field $k$. The dual vector space construction of Example \ref{ex:double_dual} defines a contravariant functor from $\fdVect_k$ to itself:
    \begin{align*}
    (~)^\ast = \Hom(-, k): \Vect_k^\op &\to \Vect_k \\
                            V &\mapsto V^\ast
\end{align*}
Moreover, for each $W \in \fdVect_k$, there is a canonical isomorphism $\ev_W: W \to W^{**}$ given by the ``evaluation at $w \in W$''; that is, $\ev_W(w): W^* \to k$ maps $\vp \in W^*$ to $\vp(w) \in k$. It is a standard exercise in linear algebra to verify that the evaluation map defines a linear isomorphism $\ev_W: W \isom W^{**}$. This defines a natural transformation 
\[\begin{tikzcd}
    {\fdVect_k} & {\fdVect_k}
    \arrow[""{name=0, anchor=center, inner sep=0}, "1_{\fdVect_k}", shift left, curve={height=-6pt}, from=1-1, to=1-2]
    \arrow[""{name=1, anchor=center, inner sep=0}, "{(~)^{**}}"', shift right, curve={height=6pt}, from=1-1, to=1-2]
    \arrow["\ev", shorten <=6pt, shorten >=6pt, Rightarrow, from=0, to=1]
\end{tikzcd}\]
from the identity to the double dual functor. Since $\ev_W$ is an isomorphism for each $W \in \fdVect_k$, this defines a natural isomorphism, so $1_{\fdVect_k} \cong (~)^{**}$. In the language of Definition \ref{defn:naturally}, $V \cong V^{**}$ naturally in $V$.

It is an important fact in linear algebra that $V \cong V^*$, but that this isomorphism is never ``canonical''; we must first fix a basis for $V$ to define the isomorphism, and the isomorphism is therefore dependent on the choice of basis. In other words, there is no natural isomorphism from $V$ to $V^*$.
\end{ex}
\subsection{Equivalence of categories}
Our discussion of naturally isomorphisms inevitably leads to the notion of an \emph{equivalence of categories}. While two elements of a set are either equal or not, equality in a general category is a more subtle affair: two objects in a category can be equal, not equal but isomorphic, or not isomorphic. The notion of equality in a category is often unreasonably strict; we usually care more about isomorphism. This means that the right notion of sameness in a category is isomorphism, while for sets, the right notion of sameness is equality. In the special case of the functor category $\Fun(\~A, \~B)$, this tells us that the right notion of sameness for functors is natural isomorphism. 

But what about categories themselves? Isomorphism is unreasonably strict, since if $\~A \cong \~B$ there are functors
\[\begin{tikzcd}
    {\~A} & {\~B}
    \arrow["F", shift left, from=1-1, to=1-2]
    \arrow["G", shift left, from=1-2, to=1-1]
\end{tikzcd}\]
such that $FG = 1_{\~A}$ and $GF = 1_{\~B}$, and we have just seen that equality between functors is too strict. To remedy this, we replace the strict equalities with isomorphisms to get
\[FG \cong 1_{\~A} \quad \text{and} \quad GF \cong  1_{\~B}\]
This is precisely how we define an equivalence of categories:
\begin{defn}
    An \emph{equivalence} between categories $\~A$ and $\~B$ consists of a pair of functors \begin{tikzcd}
	{\~A} & {\~B}
	\arrow["F", shift left=1, from=1-1, to=1-2]
	\arrow["G"', shift right=1, tail reversed, no head, from=1-1, to=1-2]
\end{tikzcd} together with natural isomorphisms 
    \[\eta: 1_\~A \isom FG, \quad \ve: GF \isom 1_\~B\]
If there exists an equivalence between $\~A$ and $\~B$, we say $\~A$ and $\~B$ are \emph{equivalent}, and write $\~A \simeq \~B$. We also call the functors $F$ and $G$ \emph{equivalences}.
\end{defn}

\begin{warn}
One should be especially careful to distinguish between the symbol $\cong$, used for isomorphisms (including isomorphisms between objects of $\Cat$), and the symbol $\simeq$, which most category theorists use to denote an equivalence of categories. 
\end{warn}

\begin{defn}
    A functor $F: \~A \to \~B$ is \emph{essentially surjective on objects} if for all $B \in \~B$, there exists $A \in \~A$ such that $F(A) \cong B$
\end{defn}

\begin{prop}
    A functor is an equivalence if and only if it is full, faithful and essentially surjective on all objects.
\end{prop}

\begin{cor}
    Let $F: \~C \to \~D$ be a full and faithful functor. Then $\~C$ is equivalent to the full subcategory $\~C'$ of $\~D$ whose objects are of the form $F(C)$ for some $C \in \~C$.
\end{cor}

\begin{proof}
    Since the functor $F': \~C \to \~C'$ defined by $F'(C) = F(C)$ is full and faithful (because $F$ is) and is essentially surjective on objects (by definition of $\~C'$), it is an equivalence of the categories $\~C$ and $\~C'$, giving the result.
\end{proof}

\section{Constructions on categories}


The \emph{principal of duality} says that every definition, theorem, or proof has a \emph{dual}, obtained by reversing all the arrows. Duality is one of the fundamental principles when studying categories. 

\begin{defn}
    Every category $\~C$ has an \emph{opposite} or \emph{dual} category $\~C^{\op}$, which is defined by reversing all of the arrows in $\~C$. Identities in $\~C^{\op}$ are the same as in $\~C$. Composition is simply defined with the arguments reversed -- if $C \To{g} B \To{f} A$ are maps in $\~C^{\op}$, then there are maps $C \xleftarrow{g} B \xleftarrow{f} A$ in $\~C$. This gives rise to a composite map $C \xleftarrow{fg} A$ in $\~C$, and the composite of the original pair of maps is the simply corresponding map $C \To{gf} A$ in $\~C^{\op}$.
\end{defn}

We can rephrase some of our previous definitions in terms of opposite categories:
\begin{itemize}
    \item \emph{An element $T$ of a category $\~C$ is terminal iff it is initial in $\~C^{\op}$.}

    \item \emph{A morphism $f: A \to B$ in a category $\~C$ is a monomorphism (or equivalently, a section) iff it is an epimorphism (equivalently, a retraction) in $\~C^{\op}$.}
\end{itemize}
Many of the proofs we have done up until this point can also be shortened with the help of duality. The following is an example of this.
\begin{prop}
Let $A \To{f} B$ and $B \To{g} C$ be two morphisms in a category $\~C$.
\begin{enumerate}[label=(\roman*)]
    \item If $f$ and $g$ are both monomorphisms, then $fg$ is a monomorphism. Similarly, if $f$ and $g$ are both epimorphisms, then $fg$ is an epimorphism.
    \item If $f$ and $g$ are both sections, then $fg$ is a section. Similarly, if $f$ and $g$ are both retractions, then $fg$ is a retraction.
\end{enumerate}
\end{prop}

\begin{proof}
\hfill
\begin{enumerate}[label=(\roman*)]
    \item We will prove the first implication, and then conclude the second by duality. To that end, let $f$ and $g$ be monomorphisms in $\~C$, and suppose there exist morphisms $a, a': A \to X$ such that $a(fg) = a'(fg)$. From this we get $(af)g = a'(fg)$, and $g$ being monic, we deduce $af = a'f$. Since $f$ is also monic, we have $a = a'$, therefore $fg$ is a monomorphism. Reversing all the arrows in this proof, ``monomorphism'' becomes ``epimorphism'', so if $f$ and $g$ are both epimorphisms, we conclude that $fg$ is an epimorphism.
    \item Suppose $f$ and $g$ are both sections. Then there exist $Y \To{f'} X$ and $Z \To{g'} Y$ such that $ff' = 1_X$ and $gg' = 1_Y$. We want a right inverse for $fg$. Consider $g'f': Z\to X$. We have
    \begin{align*}
        fg(g'f') &= f(gg')f' \\
                 &= ff' \\
                 &= 1_X
    \end{align*}
    Using duality, we may again conclude that when $f$ and $g$ are both retractions, $fg$ is a retraction.
\end{enumerate}        
\end{proof}

\subsection*{Elements}
In $\Set$ or other categories whose objects are sets, it is easy to ask about \emph{elements} $x \in X$ of a given object $X$. However, we have met many perfectly categories whose objects are not set-like. It is therefore reasonable 

\begin{defn}
If a category $\~C$ has a terminal object $T$, and $X$ is an object of $\~C$, an \emph{element} of $X$ is a map $T \to X$. 
\end{defn}

\begin{ex}
Indeed, if we are working in $\Set$, then an element of an object $X$ (in the category-theoretic context) is the same as an element of $X$ (in the set-theoretic context). Recall that in $\Set$, any singleton set $\{\ast\}$ is a terminal object. A categorical element of the set $X$ is thus a map
\begin{align*}
\{\ast\} &\to X\\
\ast &\mapsto x
\end{align*} 
corresponding to a choice of a single $x \in X$. 
\end{ex}


\begin{ex}
\label{ex:elements_in_cat}
    Similarly, in $\Cat$, a map $\mathbf{1} \to \~C$ is essentially the same as an object of the category $\~C$.
\end{ex}

Elements lead to many fascinating constructions. One particularly interesting example of this is the category $\cRing^\op$.

\begin{ex}
    What are the elements of $\Z[x, y] / \langle y^2 - x^2 - x^3 \rangle$ as an object of $\cRing^\op$? \\
    
    Since $\Z$ is initial in $\cRing$, it is terminal in $\cRing^\op$. Therefore, if $R \in \cRing^\op$, an element of $R$ is a ring homomorphism $\Z \to R$. However, for any ring $R$, there is exactly one homomorphism $\Z \To{f_0} R$, given by $f_0(k) = k$ for any $k \in \Z$ (where $1_R$ is the identity in $R$). Therefore, an element of $\Z[x, y] / \langle y^2 - x^2 - x^3 \rangle$ in $\cRing^{\op}$ is a (unique) homomorphism 
    \[\Z \To{f} \Z[x, y] / \langle y^2 - x^2 - x^3 \rangle\]
    Now we need to figure out what $f$ ``looks like''. We have that $x = f(k) = k$ and $y = f(\ell) = \ell$ for some $k, \ell \in \Z$, however, $k$ and $\ell$ are not just any integers, since we also require that 
    \[y^2 - x^2 - x^3 = 0\]
    so $x = k$ and $y = \ell$ must be \emph{solutions} to the above equation. 
    
    This means that given an equation $p(x, y) = 0$ in $\Z[x,y]$, we can find zeroes of the equation by looking at elements (in the categorical sense) of the quotient ring $\Z[x, y]/\langle p(x, y) \rangle$. Indeed, a similar procedure can be used to find solutions in polynomial rings over $\Z$ with more coefficients, however, more care must be taken when approaching polynomials over an arbitrary ring.
\end{ex}

\subsection*{Product of categories}

\begin{defn}
    Given categories $\~A$ and $\~B$, we can define  a \emph{product category} $\~A \times \~B$, in which 
    \begin{align*}
        \ob(\~A \times \~B) &= \ob(\~A) \times \ob(\~B) \\
        \~A \times \~B((A, B), (A', B')) &= \~A(A, A') \times \~B(B, B')\
    \end{align*}
    That is, an object in $\~A \times \~B$ is a pair $\langle A, B \rangle$ where $A \in \~A$ and $B \in \~B$. An arrow $\langle A, B\rangle \to \langle A', B' \rangle$ is a pair $\langle f, g \rangle$ where $f: A \to A'$ in $\~A$ and $g: B \to B'$ in $\~B$. Composition is defined by
    \[ \langle f, g \rangle  \langle f', g' \rangle := \langle f f', g g' \rangle\]
    and for the identity, we simply have
    \[1_{\~A \times \~B} := \langle 1_\~A, 1_\~B \rangle\]
\end{defn}

There are functors
\[\begin{tikzcd}
    & {\~A \times \~B} \\
    {\~A} && {\~B}
    \arrow["P"', from=1-2, to=2-1]
    \arrow["Q", from=1-2, to=2-3]
\end{tikzcd}\]
called \emph{projections}, defined on arrows by
\[P\langle f, g \rangle = f, \quad Q\langle f, g \rangle = g\]
and similarly for objects. The projections have the following property: given any category $\~C$ and functors 
\[\begin{tikzcd}
    & {\~C} \\
    {\~A} && {\~B}
    \arrow["R"', from=1-2, to=2-1]
    \arrow["S", from=1-2, to=2-3]
\end{tikzcd}\]
There exists a unique functor $F: \~C \to \~A \times \~B$ with $FP = R$ and $QF = T$, as in the following commutative diagram of functors:
\[\begin{tikzcd}
    & {\~C} \\
    \\
    & {\~A \times \~B} \\
    {\~A} && {\~B}
    \arrow["F", dashed, from=1-2, to=3-2]
    \arrow["R"', curve={height=6pt}, from=1-2, to=4-1]
    \arrow["S", curve={height=-6pt}, from=1-2, to=4-3]
    \arrow["P", from=3-2, to=4-1]
    \arrow["Q"', from=3-2, to=4-3]
\end{tikzcd}\]
That is, $P$ and $Q$ are ``universal'' among pairs of functors from $\~A$ to $\~B$. These two conditions require that for all arrows $h: C \to C'$ in $\~C$, we have $F(h) = \langle R(h), T(h) \rangle$; conversely, setting $F(h) = \langle R(h), T(h) \rangle$, we see that $F$ defines a functor with the desired properties. 

An important class of functors are those of the form $\~A \times \~B \to \~C$; such functors are known as \emph{bifunctor}, and examples abound in mathematics. For example, there is a bifunctor $\Set \times \Set \to \Set$ sending a pair of sets $\langle X, Y \rangle$ to their cartesian product $X \times Y$.


\subsection*{Comma categories}
Our next goal is to introduce an important construction known as the \emph{comma category}. This concept was first introduced by Lawvere in 1963 (\cite{lawvere1963}, p. 36), although it did not adopt widespread use until several years later. 

\begin{defn}
    Suppose we have categories and functors 
    \begin{equation}
    \label{eq:comma_cat}
    \begin{tikzcd}
    & {\~B} \\
    {\~A} & {\~C}
    \arrow["Q", from=1-2, to=2-2]
    \arrow["P"', from=2-1, to=2-2]
\end{tikzcd}\end{equation}
The \emph{comma category} $(P \downarrow Q)$ is the category defined as follows:
\begin{itemize}
    \item objects are triples $(A, h, B)$ with $A \in \~A$, $B \in \~B$, and $h: P(A) \to Q(B)$ in $\~C$;
    \item an arrow $(A, h , B) \to (A', h', B')$ is a pair $(f: A \to A', \, g: B \to B')$ of arrows such that the square 
    \[\begin{tikzcd}
    {P(A)} & {P(A')} \\
    {Q(B)} & {Q(B')}
    \arrow["{P(f)}", from=1-1, to=1-2]
    \arrow["h"', from=1-1, to=2-1]
    \arrow["{h'}", from=1-2, to=2-2]
    \arrow["{Q(g)}"', from=2-1, to=2-2]
\end{tikzcd}\]
commutes.
\end{itemize}
\end{defn}

\begin{rmk}
    Given categories and functors (\ref{eq:comma_cat}), there are canonical functors and a canonical natural transformation as shown:
    \[\begin{tikzcd}
    {(P \downarrow Q)} & {\~B} \\
    {\~A} & {\~C}
    \arrow[from=1-1, to=1-2]
    \arrow[from=1-1, to=2-1]
    \arrow["Q",from=1-2, to=2-2]
    \arrow[shorten <=14pt, shorten >=14pt, Rightarrow, from=2-1, to=1-2]
    \arrow["P"',from=2-1, to=2-2]
\end{tikzcd}\]
in a suitable $2$-categorical sense, $(P \downarrow Q)$ is universal with this property.
\end{rmk}

An important special case of comma categories are \emph{slice categories} and the dual notion of \emph{coslice categories}. The structure of these categories is detailed in the following definition:

\begin{defn}
Let $\~C$ be a category, and let $X \in \~C$. 

\begin{enumerate}[label=(\roman*)]
    \item We define the \emph{slice category} $\~C_{/X}$ of $\~C$ over $X$ as follows:
\begin{itemize}
    \item Objects are maps $f: A \to X$, for some $A \in \~C$, which we represent formally as a pair $(A, f)$.
    \item Given two objects $(A, f)$ and $(A', f')$ in $\~C_{/X}$, an arrow $g: (A, f) \to (A', f')$ in $\~C_{/X}$ is a map $g: A \to A'$ making
        \[\begin{tikzcd}[column sep=small] 
            A \arrow[dr, "f", swap] \arrow[rr, "g"]
            & & A' \arrow[dl, "f'"] \\
            & X
        \end{tikzcd}\]
        commute.
\end{itemize}

    \item The dual concept of the slice category is called the \emph{coslice category} $X_{/\~C}$. Explicitly, this means $X_{/\~C}$ contains the following data:
\begin{itemize}
    \item Objects in $X_{/\~C}$ are maps $X \To{f} A$ for some $A \in \~C$.
    \item Given a pair of objects $X \To{f} A$, $X \To{f'} A'$ in $X_{/\~C}$, a morphism in $X_{/\~C}$ is a map $A \To{g} A'$ making
    \[\begin{tikzcd}[column sep=small]
        & X \arrow[dl, "f", swap] \arrow[dr, "f'"] & \\
        A \arrow[rr, "g", swap] & & A'
    \end{tikzcd}\]
        commute.
\end{itemize}
\end{enumerate}
\end{defn}

Recall from Example \ref{ex:elements_in_cat} that a map $X: \mathbf{1} \to \~C$ is essentially the same as an object $X \in \~C$. Then the slice category $\~C_{/X}$ is equivalent to the comma category $(1_\~C \downarrow X)$, as in the diagram
    \[\begin{tikzcd}
    & {\mathbf{1}} \\
    {\~C} & {\~C}
    \arrow["{X}", from=1-2, to=2-2]
    \arrow["{1_{\~C}}"', from=2-1, to=2-2]
\end{tikzcd}\]
In principle, an object of $(1_\~C \downarrow X)$ is  a triple $(A, f, \ast)$, where $h: A \to X$. However, since $\ast$ is a fixed object, this is essentially the same as a pair $(A, f)$ where $h: X \to A$, so $(1_{\~C} \downarrow X)$ has the same objects as $\~C_{/X}$. Moreover, $(1_{\~C} \downarrow X)$ has the same maps as $\~C/X$: given two objects $(A, f)$ and $(A', f')$ in $(1_\~C \downarrow X)$, a map $(A, f) \to (A', f')$ consists of a map $g: A \to A'$ such that the square
\[\begin{tikzcd}
    A & {A'} \\
    X & X
    \arrow["g", from=1-1, to=1-2]
    \arrow["f"', from=1-1, to=2-1]
    \arrow["{f'}", from=1-2, to=2-2]
    \arrow[Rightarrow, no head, from=2-1, to=2-2]
\end{tikzcd}\]
commutes. But this is the same as asking for a commutative triangle 
\[\begin{tikzcd}[column sep=small] 
            A \arrow[dr, "f", swap] \arrow[rr, "g"]
            & & A' \arrow[dl, "f'"] \\
            & X
        \end{tikzcd}\]
So $(1_{\~C} \downarrow X)$ indeed has the same maps as $\~C_{/X}$. This gives $(1_{\~C} \downarrow X) \cong \~C_{/X}$, and duality gives us a similar statement for coslice categories: $X_{/\~C} \cong (X \downarrow 1_{\~C})$.


% The slice category has enough structure that it is worth verifying that the above definition indeed makes $\~C_{/X}$ into a category. Let $X \in \~C$, and let $(A_0, f_0), (A_1, f_1), (A_2, f_2)$, and $(A_3, f_3)$ be objects in the slice category $\~C_{/X}$. Let's first check that the composition is well-defined. Suppose we have arrows $g: (A_0, f_0) \to (A_1, f_1)$ and $h: (A_1, f_1) \to (A_2, f_2)$. Then the following triangles commute
%     \[
%         \begin{tikzcd}[column sep=small] 
%             A_0 \arrow[dr, "f_0", swap] \arrow[rr, "g"]
%             & & A_1 \arrow[dl, "f_1"] \\
%             & X
%         \end{tikzcd} 
%         \quadd 
%         \begin{tikzcd}[column sep=small] 
%             A_1 \arrow[dr, "f_1", swap] \arrow[rr, "h"]
%             & & A_2 \arrow[dl, "f_2"] \\
%             & X
%         \end{tikzcd}
%     \]
% that is, $gf_1 = f_0$ and $h f_2 = f_1$. From this we see that $(g h) f_2 = f_0$, so the diagram
%     \[\begin{tikzcd}[column sep=small] 
%             A_0 \arrow[dr, "f_0", swap] \arrow[rr, "gh"]
%             & & A_2 \arrow[dl, "f_2"] \\
%             & X
%         \end{tikzcd}\]
% also commutes, and therefore the composition is well-defined. Let's now check associativity. Suppose we have an arrow $i: (A_2, f_2) \to (A_3, f_3)$ (in addition to $g$ and $h$ above). Then the triangle
%     \[\begin{tikzcd}[column sep=small] 
%             A_2 \arrow[dr, "f_2", swap] \arrow[rr, "i"]
%             & & A_3 \arrow[dl, "f_3"] \\
%             & X
%         \end{tikzcd}\]
% commutes. We previously found that $(g h) f_2 = f_0$, and substituting $f_2 = i f_3$ gives $(g h) i f_3 = f_0$
% Now, the map $h i: (A_1, f_1) \to (A_3, f_3)$ makes 
%     \[\begin{tikzcd}[column sep=small] 
%             A_1 \arrow[dr, "f_1", swap] \arrow[rr, "h i"]
%             & & A_3 \arrow[dl, "f_3"] \\
%             & X
%         \end{tikzcd}\]
% commute, so $(h i) f_3 = f_1$. We also have $g: (A_0, f_0) \to (A_1, f_1)$, so $g f_1 = f_0$, and substituting $f_1 = (h i) f_3$ in this expression, we find that
%     \[g (h i) f_3 = f_0 = (g h) i f_3\]
% therefore the composition in $\~C_{/X}$ is associative. 

% We now turn our attention to identities. Given an object $(A, f)$ in $\~C_{/X}$, the identity on $(A, f)$ is the map $1_A: A \to A$, as in the commutative triangle
%     \[\begin{tikzcd}[column sep=small]
%     A && A \\
%     & X
%     \arrow[Rightarrow, no head, from=1-1, to=1-3]
%     \arrow["f"', from=1-1, to=2-2]
%     \arrow["f", from=1-3, to=2-2]
% \end{tikzcd}\]
% Given an arrow $g: (A_0, f_0) \to (A_1, f_1)$ in $\~C_{/X}$, it is clear that $(1_{A_0})g = g = g (1_{A_1})$, so the identity law holds, so $\~C_{/X}$ is a category.

% Upon first glance, the definition of a slice category may seem slightly random, however, it is useful for defining structures where we want to focus on certain maps into an object $X \in \~C$. For maps \emph{out} of $X$, we have a dual notion, known as the coslice category. 

% \begin{defn}
% Let $\~C$ be a category, and let $X \in \~C$. The \emph{coslice category} $X_{/~C}$ consists of the following data:
% \begin{itemize}
%     \item Objects in $X_{/~C}$ are maps $X \To{f} A$ for some $A \in \~C$.
%     \item Given a pair of objects $X \To{f} A$, $X \To{f'} A'$ in $X_{/~C}$, a morphism in $X_{/~C}$ is a map $A \To{g} A'$ making
%     \[\begin{tikzcd}[column sep=small]
%         & X \arrow[dl, "f", swap] \arrow[dr, "f'"] & \\
%         A \arrow[rr, "g", swap] & & A'
%     \end{tikzcd}\]
%         commute.
% \end{itemize}
% \end{defn}

\begin{ex}
    Let $\ast$ denote the one-point set, viewed as a discrete category. Then the coslice category $\ast_{/\Set}$ is called the \emph{category of pointed sets}, and denoted by $\Set_\ast$. Objects in this category are pairs $(X, x)$, where $X$ is a set and $x: \ast \to X$ is a map of sets. However, the object $\ast$ is terminal in $\Set$, so the map $x : \ast \to X$ is the same as choosing an element $x \in X$. Therefore, an object in $\Set_\ast$ can be represented as a pair $(X, x)$ consisting of a set $S$ and an object $x \in X$. Given pointed sets $(X, x)$ and $(Y, y)$, an arrow $(X, x) \to (Y, y)$ is a function $f: X \to Y$ rendering the triangle
    \[\begin{tikzcd}[column sep=small]
        & \ast \arrow[dl, "x", swap] \arrow[dr, "y"] & \\
        X \arrow[rr, "f", swap] & & Y
    \end{tikzcd}\] 
    commutative, i.e., $f(x) = y$.
\end{ex}

\begin{ex}
    Let $R$ be a commutative ring with 1. We claim that $R/\cRing = \cAlg_R$, the category of commutative $R$-algebras. To see that this is the case, recall that an element of $R/\cRing$ is a map $R \To{f} A$ for some $A \in \cRing$. Since $A$ is commutative, the homomorphism of rings $f: R \to A$ makes $A$ into a commutative $R$-algebra. Recall that this means $A$ has a natural $R$-module structure defined by 
    \[r \cdot a = a \cdot r := f(r)a\]
    we will be careful to only use the notation $r \cdot a$ when referring to this scalar multiplication. Now that we've concluded objects of $R/\cRing$ are indeed commutative $R$-algebras, we must also verify that a morphism in $R/\cRing$ is a morphism in $\cAlg_R$ (i.e., an $R$-algebra homomorphism). To that end, suppose we have two objects $R \To{f} A$, $R \To{f'} B$ in $R/\cRing$. Then a morphism $f \To{\vp} f'$ in $R/\cRing$ is a commutative triangle
        \begin{equation}
            \label{eq:calgebra_triangle}
            \begin{tikzcd}[column sep=small]
                & R \arrow[dl, "f", swap] \arrow[dr, "f'"] & \\
                A \arrow[rr, "\vp", swap] & & B
            \end{tikzcd}
        \end{equation}
    and since $A \To{\vp} B$ is a morphism in $\cRing$, it is clear that $\vp(1_A) = 1_B$. 
    Now, we need only check that $\vp(r \cdot_A a) = r \cdot_B \vp(a)$, where $\cdot_A$ and $\cdot_B$ denote the scalar multiplication in the $R$-algebras $A$ and $B$, respectively. To do this, we first write
        \[\vp(r \cdot_A a) = \vp(f(r)a) = f \vp(r) \cdot_B \vp(a)\]
    But note that since the triangle (\ref{eq:calgebra_triangle}) commutes, we have $f\vp(r) = f'(r)$ for all $r \in R$, so we get
        \[\vp(r \cdot_A a) = f'(r)\vp(a)\]
    Now, $f'(r)\vp(A)$ is just the multiplication of $\vp(A) \in B$ by the scalar $r$ in $R$-algebra $B$, so we obtain
        \[\vp(r \cdot_A a) = r \cdot_B \vp(A)\]
    as desired. From this, we see that the objects and morphisms in $R/\cRing$ are the same as the objects and morphisms in $\cAlg_R$, which gives the result.
\end{ex}

\section{Groups and monoids in category theory}

This short section explores two basic categorical notions that appear in numerous areas of mathematics, those being the notions of monoids and groups \emph{internal to a category} $\~C$

\section{Remarks on the theory of sets}


\chapter{Adjoint functors}
\label{chap:3}

\epigraph{\itshape Adjoint functors arise everywhere.}{---Saunders Mac Lane, \emph{Categories for the Working Mathematician} \cite{maclane}}

\section{Definition and first consequences}

\begin{notation}
    Henceforth, we write $f g$ to denote the composition $f \circ g$, and we frequently write $Fg$ to denote the value of a functor $F$ at a morphism $g$.
\end{notation}

Consider a pair of functors $F: \~A \to \~B$ and $G: \~B \to \~A$ in opposite directions. In plain language, $F$ is said to be \emph{left adjoint} to $G$ if, whenever $A \in \~A$ and $B \in \~B$, maps $F(A) \to B$ are essentially the same as maps $A \to G(B)$. To make this notion precise, we need to be clear about what we mean by ``essentially the same''.
\begin{defn}
\label{defn:adjunction}
Let $\begin{tikzcd}
    {\~A} & {\~B}
    \arrow["F", shift left, from=1-1, to=1-2]
    \arrow["G", shift left, from=1-2, to=1-1]
\end{tikzcd}$ be functors between categories. We say that $F$ is \emph{left adjoint} to $G$, and $G$ is \emph{right adjoint} to $F$, and write $F \dashv G$, if
\begin{equation}
\label{eq:adjoint}
\~B(F(A), B) \cong \~A(A, G(B))
\end{equation}
naturally in $\~A$ in $\~B$ (we define what ``naturally'' means below). An \emph{adjunction} between $F$ and $G$ is a choice of natural isomorphism (\ref{eq:adjoint}). 
\end{defn}
To describe what ``naturally in $\~A$ in $\~B$'' means in the above definition, we start by constructing a bijection between maps $F(A) \to B$ and maps $A \to G(B)$.   Given functors $F$ and $G$ as above, we can define a correspondence $(\bar{\,\,\,}): \~B(F(A), B) \isom \~A(A, G(B))$ (`bar') between maps $F(A) \to B$ and $A \to G(B)$ by
\begin{align*}
\left(F(A) \To{g} B \right) &\mapsto \left(A \To{\bar{g}} G(B) \right) \\
\left(F(A) \To{\bar{f}} B \right) &\mapsfrom \left(A \To{f} G(B) \right) 
\end{align*}
so $\bar{\bar{g}} = g$ and $\bar{\bar{f}} = f$. The map $\bar{g}$ is called the \emph{transpose} of $g$, and similarly for $f$. With this notion established, the naturality axiom can be stated as:
\begin{equation}
\label{eq:bar_axiom1}
\overline{\left(F(A) \To{g} B \To{s} B' \right)} = \left(A \To{\bar{g}} G(B) \To{G(s)} G(B') \right)
\end{equation}
that is, $\overline{sg} = G(s)\bar{g} $ for all $g: F(A) \to B$ and $s: B \to B'$ in $\~B$, and
\begin{equation}
\label{eq:bar_axiom2}
\overline{\left(A' \To{t} A \To{f} G(B) \right)} = \left(F(A') \To{F(t)} F(A) \To{\bar{f}} B \right)
\end{equation}
i.e., $\overline{ft} =  \bar{f}F(t) $ for all $f:A \to G(B)$ and $t: A' \to A$ in $\~A$.

A given functor $F$ may not have an adjoint, but when it does, this adjoint is always unique, so we may speak of \emph{the} adjoint of $F$. As was the case for limits, this uniqueness property turns out to be one of the most crucial properties of adjoints. We will wait until Chapter \ref{chap:5} (Proposition ) to prove that adjoints are unique, since this uniqueness turns out to be an easy consequence of Yoneda's lemma (which is the central result in Chapter \ref{chap:5}).

\begin{rmk}
    Suppose we have categories $\~A$ and $\~B$ and adjoint functors $\begin{tikzcd}
    {\~A} & {\~B}
    \arrow[""{name=0, anchor=center, inner sep=0}, "F", shift left=2, from=1-1, to=1-2]
    \arrow[""{name=1, anchor=center, inner sep=0}, "G", shift left=2, from=1-2, to=1-1]
    \arrow["\dashv"{anchor=center, rotate=-90}, draw=none, from=0, to=1]
\end{tikzcd}$. Then given sequences 
\[A_0 \to A_1 \to \cdots \to A_n  \text{ in } \~A \]
and
\[B_0 \to B_1 \to \cdots \to B_m \text{ in } \~B\]
and a map $f \in \~B(F(A_n), B_0)$, we can form the composite map
\[F(A_n) \To{f} B_0 \to B_1 \to \cdots \to B_m\]
The naturality axiom implies that exists a unique map $A_0 \to G(B_m)$, given by the composite 
\[A_0 \to \cdots \to A_n \To{\bar{f}} G(B_0) \to \cdots \to G(B_m)\]
This is similar to the remarks we made when first describing categories, functors, and natural transformations.
\end{rmk}

\begin{exs}
\label{ex:adjoints}
\hfill
    \begin{enumerate}[label=(\roman*)]
        \item \label{ex:adjoint_vspace} Let $k$ be a field. Then the forgetful functor $U: \Vect_k \to \Set$ of Example \ref{ex:forgetful_functors} is right adjoint to the free functor $F: \Set \to \Vect_k$ of Example \ref{ex:free_functors}:
        \[\begin{tikzcd}
            {\Set} & {\Vect_k}
            \arrow[""{name=0, anchor=center, inner sep=0}, "F", shift left=2, from=1-1, to=1-2]
            \arrow[""{name=1, anchor=center, inner sep=0}, "U", shift left=2, from=1-2, to=1-1]
            \arrow["\dashv"{anchor=center, rotate=-90}, draw=none, from=0, to=1]
        \end{tikzcd}\]
        To see that this is the case, let $S$ be a set and $V$ a vector space. Given a linear map $g: F(S) \to V$, we can define a map of sets $\bar{g}: S \to U(V)$ by $\bar{g}(s) = g(s)$ for all $s \in S$. This gives us a function 
        \begin{align*}
            (\bar{\,\,\,}): \Vect_k(F(S), V) &\to \Set(S, U(V)) \\
            g &\mapsto \bar{g}
        \end{align*}
        Similarly, given a map of sets $f: S \to U(V)$, we can define a linear map $\bar{f}: F(S) \to V$ by 
        \[\bar{f}\left(\sum_{s \in S} \lambda_s s \right) = \sum_{s \in S} \lambda_s f(s) \quad \text{for all} \,  \sum_{s \in S} \lambda_s s  \in F(S)\]
        This assignment gives us a function 
        \begin{align*}
            (\bar{\,\,\,}): \Set \left(S, U(V) \right) &\to \Vect_k \left(F(S), V \right) \\
            f &\mapsto \bar{f}
        \end{align*}
        And these two 'bar' functions are mutually inverse: if $g: F(S) \to V$ is a linear map, then for any $\sum_{s \in S}\lambda_s s \in F(S)$, we have
        \begin{align*}
            \bar{\bar{g}} \left(\sum_{s \in S}\lambda_s s\right) &= \sum_{s \in S} \lambda_s \bar{g}(s) \\
            &= \sum_{s \in S} \lambda_s g(s) \\
            &= g \left(\sum_{s \in S}\lambda_s s\right)
        \end{align*}
        and if $f: S \to U(V)$ is a map in $\Set(S, U(V))$, then we have
        \[\bar{\bar{f}}(s) = \bar{f}(s) = f(s)\]
        This gives the desired isomorphism $\Vect_k(F(S), V) \cong \Set(S, U(V))$.


        \item Similarly, we have an adjunction 
        \[\begin{tikzcd}
            {\Set} & {\Grp}
            \arrow[""{name=0, anchor=center, inner sep=0}, "F", shift left=2, from=1-1, to=1-2]
            \arrow[""{name=1, anchor=center, inner sep=0}, "U", shift left=2, from=1-2, to=1-1]
            \arrow["\dashv"{anchor=center, rotate=-90}, draw=none, from=0, to=1]
        \end{tikzcd}\]
        Where $U$ is the forgetful functor $\Grp \to \Set$ and $F: \Set \to \Grp$ is the free group functor, which is trickier to describe.

        \item Consider the inclusion functor $U: \Ab \mono \Grp$ defined by $U(A) = A$ for any abelian group $A$ and $U(f) = f$ for a homomorphism $f$ of abelian groups. There is an adjunction
        \[\begin{tikzcd}
            {\Grp} & {\Ab}
            \arrow[""{name=0, anchor=center, inner sep=0}, "F", shift left=2, from=1-1, to=1-2]
            \arrow[""{name=1, anchor=center, inner sep=0}, "U", shift left=2, from=1-2, to=1-1, hook']
            \arrow["\dashv"{anchor=center, rotate=-90}, draw=none, from=0, to=1]
        \end{tikzcd}\]
        Here $F(G)$ is the \emph{abelianization} of the group $G$. This is an abelian quotient group of $G$ such that every map from $G$ to an abelian group $A$ factors uniquely through $F(G)$:
        \begin{equation}
        \label{eq:abelianization}
        \begin{tikzcd}
            G & {F(G)} \\
            & {\forall A}
            \arrow["\pi", from=1-1, to=1-2]
            \arrow["{\forall \varphi}"', from=1-1, to=2-2]
            \arrow["{\exists! \bar{\varphi}}", dashed, from=1-2, to=2-2]
        \end{tikzcd}
        \end{equation}
        Where $\pi: G \to F(G)$ is the usual projection map from a group to its quotient. To define the bijection 
        \[\Ab(F(G), A) \isom \Grp(G, U(A))\]
        we'll produce two maps that are inverses of one another. In the first direction, given a map of abelian groups $f: F(G) \to A$, define a map $\bar{f}: G \to U(A)$ by
        \[\bar{f} = f \pi\]
        where $\pi: G \to F(G)$ is the top arrow in the triangle (\ref{eq:abelianization}). Conversely, given a group homomorphism $g: G \to U(A)$, define $\bar{g}$ to be the unique abelian group homomorphism $F(G) \to U(A)$ making the triangle
        \[\begin{tikzcd}
            G & {F(G)} \\
            & {U(A)}
            \arrow["\pi", from=1-1, to=1-2]
            \arrow["{g}"', from=1-1, to=2-2]
            \arrow["{\bar{g}}", dashed, from=1-2, to=2-2]
        \end{tikzcd}\]
        commute. 

        \item The inclusion functor $\Grp \mono \Mon$ is interesting because it has both a left and a right adjoint:
        \[\begin{tikzcd}
            \Grp & \Mon
            \arrow[""{name=0, anchor=center, inner sep=0}, hook, from=1-1, to=1-2]
            \arrow[""{name=1, anchor=center, inner sep=0}, "R", shift left=4, from=1-2, to=1-1]
            \arrow[""{name=2, anchor=center, inner sep=0}, "F"', shift right=4, from=1-2, to=1-1]
            \arrow["\dashv"{anchor=center, rotate=-90}, draw=none, from=0, to=1]
            \arrow["\dashv"{anchor=center, rotate=-90}, draw=none, from=2, to=0]
        \end{tikzcd}\]
        Here the middle functor is the inclusion, and the functor $R$ is straightforward: for a monoid $M$, $R(M)$ is the group consisting of all invertible elements. The functor $F$ is trickier to describe, and involves formally adding inverses for each non-invertible monoid element. We say that the category $\Grp$ is both a \emph{reflexive} and \emph{coreflexive} subcategory of $\Mon$, meaning that the inclusion functor $\Grp \mono \Mon$ has both a left and a right adjoint.

        \item Let $\Field$ denote the category whose objects are fields, with ring homomorphisms as arrows. The forgetful functor $\Field \to \Set$ does \emph{not} have a left adjoint. This is because the theory of fields is different from groups, rings, etc., in that it is not an \emph{algebraic theory}. 

        To explain this terminology, let's first consider the theory of groups. A group consists of a set $G$ together with an associative composition law $\ast: G^2 = G \times G \to G$, an inverse $(~)^{-1}: G \to G$ for each element, and an identity $e$, which can be described as a map $e: 1 \to G$, where $1$ denotes the one-element set (and by definition $G^0 = 1$). Thus, a group can be described as a collection of 3 maps
        \[\ast: G^2 \to G, \quad (~)^{-1}: G^1 \to G, \quad e: G^0 \to G\]
        with 2, 1, and 0 inputs respectively. Of course, we require that these operations satisfy the familiar group equations. The theory of groups which we have just described is a specific example of a larger phenomenon known as an \emph{algebraic theory}. An algebraic theory consists of two things: first, a collection of operations, each with a specified number of inputs or \emph{arity}, and second, a set of equations. So for example, the theory of groups has three operations: one of arity 2, one of arity 1, and one of arity 0. 

        An \emph{algebra} or \emph{model} for an algebraic theory consists of a set $X$ together with a specified map $X^n \to X$ for each operation of arity $n$, such that the desired equations hold everywhere. For instance, a model for the algebraic theory of groups is exactly a group. 
        The main property of an algebraic theory is that the equations hold \emph{everywhere}. This is why the theory of fields is not an algebraic theory: the inversion map $x \mapsto x^{-1}$ is not defined for \emph{all} $x$, only $x \ne 0$.

        \item There are adjunctions
        \[\begin{tikzcd}
            \Top & \Set
            \arrow[""{name=0, anchor=center, inner sep=0}, hook, from=1-1, to=1-2]
            \arrow[""{name=1, anchor=center, inner sep=0}, "I", shift left=4, from=1-2, to=1-1]
            \arrow[""{name=2, anchor=center, inner sep=0}, "D"', shift right=4, from=1-2, to=1-1]
            \arrow["\dashv"{anchor=center, rotate=-90}, draw=none, from=0, to=1]
            \arrow["\dashv"{anchor=center, rotate=-90}, draw=none, from=2, to=0]
        \end{tikzcd}\]
        where the middle functor is the inclusion $\Top \mono \Set$ (sending a topological space to its underlying set of points), the functor $D$ equips a set with the discrete topology, and the functor $I$ equips a set with the indiscreet topology.

        \item  \label{ex:adjoints6} Fix a set $B$. Then for any set $A$, we can construct the cartesian product of $A$ and $B$. The process of taking the cartesian product of $B$ with another set is functorial: there is a functor
        \begin{align*}
            - \times B : \Set &\to \Set \\
            A &\mapsto A \times B
        \end{align*}
        There also get a functor $\Set \to \Set$ defined by sending a set $C$ to the set $C^B := \Set(B, C)$ of functions $B \to C$:
        \begin{align*}
            \Set(B, -): \Set &\to \Set \\
            C &\mapsto C^B
        \end{align*}
        Indeed, we can show that there is a canonical bijection 
        \[\Set(A \times B, C) \cong \Set(A, C^B)\]
        for all sets $A$ and $C$. This bijection is quite simple to define: given a map $g: A \times B \to C$, define a new map $\bar{g}: A \to C^B$ by
        \[\left(\bar{g}(a)\right)(b) = g(a, b), \quad \text{for all } a \in A, \, b \in B\]
        Conversely, given a map $f: A \to C^B$, define $\bar{f}: A \times B \to C$ by 
        \[\bar{f}(a, b) = \left(f(a)\right)(b)\]
        These two `bar' functions define an adjunction
        \[\begin{tikzcd}
            {\Set} & {\Set}
            \arrow[""{name=0, anchor=center, inner sep=0}, "- \times B", shift left=2, from=1-1, to=1-2]
            \arrow[""{name=1, anchor=center, inner sep=0}, "(~)^B", shift left=2, from=1-2, to=1-1]
            \arrow["\dashv"{anchor=center, rotate=-90}, draw=none, from=0, to=1]
        \end{tikzcd}\]
        for any set $B$.
    \end{enumerate}
\end{exs}

Initial and terminal objects can also be described using adjoints: Let $\~C$ be a category and let $\mathbf{1}$ denote the category with one object $\ast$ and one map $1_\ast: \ast \to \ast$. Then there is precisely one functor $\~C \to \mathbf{1}$, and an element $X \in \~C$ can be described via the functor $\mathbf{1} \to \~C$ which maps $\ast$ to $X$. Viewing objects of $\~C$ as functors $\mathbf{1} \to \~C$, a left adjoint to the unique functor $\~C \to \mathbf{1}$ is exactly an initial object of $A$. Similarly, a right adjoint to the functor $\~C \to \mathbf{1}$ is a terminal object of $\~C$. 

\section{Units and counits}

Suppose $\begin{tikzcd}
    {\~A} & {\~B}
    \arrow[""{name=0, anchor=center, inner sep=0}, "F", shift left=2, from=1-1, to=1-2]
    \arrow[""{name=1, anchor=center, inner sep=0}, "G", shift left=2, from=1-2, to=1-1]
    \arrow["\dashv"{anchor=center, rotate=-90}, draw=none, from=0, to=1]
\end{tikzcd}$ is an adjunction. This naturally gives rise to the following composite functors:
\[GF: \~A \to \~A \quad \text{and} \quad FG: \~B \to \~B\]
Now, for each $A \in \~A$, we can define a map $\eta_A: A \to GF(A)$ by 
\[\eta_A = \overline{\left(F(A) \To{1_{F(A)}} F(A)\right)}\]
and similarly, for each $B \in \~B$, we can define a map $\ve_B: FG(B) \to B$ by
\[\ve_B = \overline{\left(G(B) \To{1_{G(B)}} G(B)\right)}\]
The maps $\eta$ and $\ve$ define natural transformations
\[\begin{tikzcd}
    {\~A} & {\~A}
    \arrow[""{name=0, anchor=center, inner sep=0}, "GF"', curve={height=12pt}, from=1-1, to=1-2]
    \arrow[""{name=1, anchor=center, inner sep=0}, "1", curve={height=-12pt}, from=1-1, to=1-2]
    \arrow["\eta", shorten <=3pt, shorten >=3pt, Rightarrow, from=1, to=0]
\end{tikzcd} \quad \text{and} \quad \begin{tikzcd}
    {\~B} & {\~B}
    \arrow[""{name=0, anchor=center, inner sep=0}, "1"', curve={height=12pt}, from=1-1, to=1-2]
    \arrow[""{name=1, anchor=center, inner sep=0}, "FG", curve={height=-12pt}, from=1-1, to=1-2]
    \arrow["\ve", shorten <=3pt, shorten >=3pt, Rightarrow, from=1, to=0]
\end{tikzcd}\]
called the \emph{unit} and \emph{counit} of the adjunction respectively.

\begin{lem}[Triangle identities]
\label{lem:tri_identities}
    Given an adjunction $\begin{tikzcd}
    {\~A} & {\~B}
    \arrow[""{name=0, anchor=center, inner sep=0}, "F", shift left=2, from=1-1, to=1-2]
    \arrow[""{name=1, anchor=center, inner sep=0}, "G", shift left=2, from=1-2, to=1-1]
    \arrow["\dashv"{anchor=center, rotate=-90}, draw=none, from=0, to=1]
\end{tikzcd}$ with unit $\eta$ and counit $\ve$, the triangles
\[\begin{tikzcd}
    F & FGF \\
    & F
    \arrow["{F \eta}", from=1-1, to=1-2]
    \arrow[Rightarrow, no head, from=1-1, to=2-2]
    \arrow["\ve F", from=1-2, to=2-2]
\end{tikzcd} \qquad \begin{tikzcd}
    G & GFG \\
    & G
    \arrow["{\eta G}", from=1-1, to=1-2]
    \arrow[Rightarrow, no head, from=1-1, to=2-2]
    \arrow["{G \ve}", from=1-2, to=2-2]
\end{tikzcd}\]
commute in the functor categories $\Fun(\~A, \~B)$ and $\Fun(\~B, \~A)$ respectively. 
\end{lem}

An equivalent statement is that the triangles 
\begin{equation}
\label{eq:triangle_ids}
\begin{tikzcd}
    F(A) & FGF(A) \\
    & F(A)
    \arrow["{F(\eta_A)}", from=1-1, to=1-2]
    \arrow[Rightarrow, no head, from=1-1, to=2-2]
    \arrow["\ve_{F(A)}", from=1-2, to=2-2]
\end{tikzcd} \qquad \begin{tikzcd}
    G(B) & GFG(B) \\
    & G(B)
    \arrow["{\eta_{G(B)}}", from=1-1, to=1-2]
    \arrow[Rightarrow, no head, from=1-1, to=2-2]
    \arrow["{G(\ve_B)}", from=1-2, to=2-2]
\end{tikzcd} \end{equation}
commute for all $A \in \~A$ and $B \in \~B$. 

% \begin{warn}
%     When composing several arrows using diagrammatic order, new arrows are added to the \emph{back} instead of the front. For example, in our diagrammatic notation, we have
%     $\eta_A: A \to FG(A)$, and applying the functor $F$ gives $F(\eta_A): F(A) \to FGF(A)$ (\emph{not} $F(\eta_A): F(A) \to FFG(A)$ -- the composite in the codomain doesn't even make sense in this case).
% \end{warn}

\begin{proof}
    We prove that the triangles (\ref{eq:triangle_ids}) commute. Take any $A \in \~A$. We have
    \begin{align*}
        \overline{\eta_A} = \overline{\left(A \To{\eta_A} GF(A) \To{1_{GF(A)}} GF(A)\right)} &= F(A) \To{F(\eta_A)} FGF(A) \To{\overline{1_{GF(A)}}} F(A) & \left(\text{using (\ref{eq:bar_axiom2})} \right) \\
        &= F(A) \To{F(\eta_A)} FGF(A) \To{\ve_{F(A)}} F(A) &\left(\text{since } \overline{1_{GF(A)}} = \ve_{F(A)} \right)
    \end{align*}
    but $\overline{\eta_A} = \overline{\overline{1_{F(A)}}} = 1_{F(A)}$, so $1_{F(A)} = \ve_{F(A)} F(\eta_A) $, which gives the first triangle identity. The second identity follows by duality.
\end{proof}

The amazing part about units and counits is that they completely characterize the transpose (and therefore the entire adjunction), even though they are only defined in terms of the transpose of the identities. The next set of results spell this out in full detail.

\begin{lem}
\label{lem:transpose}
If $\begin{tikzcd}
    {\~A} & {\~B}
    \arrow[""{name=0, anchor=center, inner sep=0}, "F", shift left=2, from=1-1, to=1-2]
    \arrow[""{name=1, anchor=center, inner sep=0}, "G", shift left=2, from=1-2, to=1-1]
    \arrow["\dashv"{anchor=center, rotate=-90}, draw=none, from=0, to=1]
\end{tikzcd}$ is an adjunction with unit $\eta$ and counit $\ve$, then the underlying isomorphism between $\~B(F(A), B)$ and $\~A(A, G(B))$ (i.e., the transpose) is given by
\[\bar{g} = G(g)\eta_A  \quad \text{for all } g: F(A) \to B\]
and
\[\bar{f} = \ve_B F(f)  \quad \text{for all } f: A \to G(B)\]
\end{lem}
\begin{proof}
    For any map $g: F(A) \to B$, we have 
    \begin{align*}
        \overline{\left( F(A) \To{g} B \right)} &= \overline{\left( F(A)\To{1_{F(A)}} F(A) \To{g} B \right)} \\
        &= A \To{\overline{1_{F(A)}}} GF(A) \To{G(g)} G(B) \\
        &= A \To{\eta_A} GF(A) \To{G(g)} G(B) & \text{Since } \overline{1_{F(A)}} = \eta_A
    \end{align*}
    which gives $\bar{g} = G(g)\eta_A$, as claimed. The fact that $\bar{f} =\ve_B F(f)$ follows by duality.
\end{proof}

\begin{thm}
\label{thm:adj_triangles}
    Take categories and functors $\begin{tikzcd}
    {\~A} & {\~B}
    \arrow["F", shift left, from=1-1, to=1-2]
    \arrow["G", shift left, from=1-2, to=1-1]
\end{tikzcd}$. There is a one-to-one correspondence between:
\begin{enumerate}[label=(\roman*)]
    \item adjunctions $F \dashv G$;
    \item pairs $(1_{\~A} \To{\eta} GF, FG \To{\ve} 1_{\~B})$ of natural transformations that satisfy the triangle identities.
\end{enumerate}
\end{thm}

\begin{proof}
    We have already seen (Lemma \ref{lem:tri_identities}) that an adjunction $\begin{tikzcd}
    {\~A} & {\~B}
    \arrow[""{name=0, anchor=center, inner sep=0}, "F", shift left=2, from=1-1, to=1-2]
    \arrow[""{name=1, anchor=center, inner sep=0}, "G", shift left=2, from=1-2, to=1-1]
    \arrow["\dashv"{anchor=center, rotate=-90}, draw=none, from=0, to=1]
\end{tikzcd}$ gives rise to a pair $(\eta, \ve)$ satisfying the triangle identities. We must show that there this process is bijective. To that end, let $(\eta, \ve)$ be a pair of natural transformations that satisfy the triangle identities. We must show that there is a unique adjunction between $F$ and $G$ (i.e., a choice of isomorphism (\ref{eq:adjoint}) satisfying the naturality equations (\ref{eq:bar_axiom1}) and (\ref{eq:bar_axiom2})) with unit $\eta$ and counit $\ve$.

Uniqueness follows from Lemma $(\ref{lem:transpose})$, since the adjunction is uniquely defined by $\bar{g} = G(g) \eta_A$ for all $g: F(A) \to B$ and $\bar{f} = F(f) $ for all $f: A \to G(B)$. For existence, take natural transformations $\eta: 1_{\~A} \to FG$ and $\ve: GF \to 1_{\~B}$ that satisfy the triangle identities. For each $A \in \~A$ and $B \in \~B$, define maps
\begin{equation}
\label{eq:functions}
\~B(F(A), B) \rightleftarrows \~A(A, G(B))
\end{equation} 
both denoted by a bar, as follows. For each $g \in \~B(F(A), B)$, set $\bar{g} = G(g)\eta_A$. Conversely, for each $f \in \~A(A, G(B))$, set $\bar{f} = \ve_B F(f)$. We need to show that the maps $g \mapsto \bar{g}$ and $f \mapsto \bar{f}$ are inverses of one another. To see that this is the case, let $g: F(A) \to B$ be a map in $\~B$, and consider the following commutative diagram:
\[\begin{tikzcd}
    {F(A)} & {FGF(A)} & {FG(B)} \\
    & {F(A)} & B
    \arrow["{F(\eta_A)}", from=1-1, to=1-2]
    \arrow[Rightarrow, no head, from=1-1, to=2-2]
    \arrow["{FG(g)}", from=1-2, to=1-3]
    \arrow["{\ve_{F(A)}}", from=1-2, to=2-2]
    \arrow["{\ve_B}", from=1-3, to=2-3]
    \arrow["g"', from=2-2, to=2-3]
\end{tikzcd}\]
where the triangle on the left commutes by the triangle identities and the right-hand square is simply the naturality square for $\ve$ with the objects $F(A)$ and $B$. Going one way around the outside of the diagram, we get
\begin{align*}
\ve_B \circ FG(g) \circ F(\eta_A) &= \ve_B \circ F(G(g)) \circ F(\eta_A)  \\
                    &= \ve_B \circ F\left(G(g) \circ \eta_A \right)\\
                    &= \ve_B \circ F(\bar{g})\\
                    &= \bar{\bar{g}}
\end{align*}
but going around the other way, we see that this is equal to $1_{F(A)}g = g$, so $\bar{\bar{g}} = g$, as claimed. Dually, for any map $f: A \to G(B)$, we have $\bar{\bar{f}} = f$. 

Next, we check that the naturality equations (\ref{eq:bar_axiom1}) and (\ref{eq:bar_axiom2}) hold. For the first equation, let $B, B' \in \~B$ and suppose we have a map $s: B \to B'$. Then for any $g: F(A) \to B$, we have
\begin{align*}
    \overline{sg} &= G(sg) \circ \eta_A  \\
    &= G(s) \circ G(g) \circ \eta_A   \\
    &= G(s) \circ \bar{g} 
\end{align*}
which gives the first equation. The second equation is similarly straightforward to check. Therefore the functions (\ref{eq:functions}) define an adjunction. Finally, the unit and counit of this adjunction are indeed $\eta$ and $\ve$, since for any $A \in \~A$ we have
\begin{align*}
\overline{1_{F(A)}} &= G(1_{F(A)}) \circ \eta_A  \\
                    &= 1 \circ \eta_A \\
                    &= \eta_A
\end{align*}
and dually for the counit.
\end{proof} 

\begin{cor}
    Let $\begin{tikzcd}
    {\~A} & {\~B}
    \arrow["F", shift left, from=1-1, to=1-2]
    \arrow["G", shift left, from=1-2, to=1-1]
\end{tikzcd}$ be categories and functors. Then $F \dashv G$ if and only if there exist natural transformations $\eta: 1_{\~A} \to GF$ and $\ve: FG \to 1_{\~B}$ satisfying the triangle identities.
\end{cor}

\section{Adjoints via initial objects}

\begin{ex}[Universal arrows]
Let $F: \~C' \to \~C$ be a functor and let $X \in \~C$. We can form the comma category $(X \downarrow F)$, as in the diagram
\[\begin{tikzcd}
    & {\~C'} \\
    {\mathbf{1}} & {\~C}
    \arrow["{F}", from=1-2, to=2-2]
    \arrow["{X}"', from=2-1, to=2-2]
\end{tikzcd}\]
Its objects are pairs $\left(A \in \~C, f: X \to F(A)\right)$, and a map $(A, f) \to (A', f')$ is a map $t: A \to A'$ making the triangle 
\[\begin{tikzcd}
    X & {F(A)} \\
    & {F(A')}
    \arrow["f", from=1-1, to=1-2]
    \arrow["{f'}"', from=1-1, to=2-2]
    \arrow["{F(t)}", from=1-2, to=2-2]
\end{tikzcd}\]
commute. 

Now let $B \in \~C$, and suppose that the object $(B, g: X \to F(B))$ is initial in the comma category $(X \downarrow F)$. This means that there is only one map \emph{out of} $(B, g)$. More precisely, for every object $(C, h)$ in $(X \downarrow F)$, there is exactly one map $u: B \to C$ rendering the following triangle commutative:
 \begin{equation}
 \label{eq:universal_triangle}
 \begin{tikzcd}
    X & {F(B)} \\
    & {F(C)}
    \arrow["g", from=1-1, to=1-2]
    \arrow["{h}"', from=1-1, to=2-2]
    \arrow["{F(u)}", from=1-2, to=2-2, dashed]
\end{tikzcd}
\end{equation}
Thus, an initial object in the comma category $(X \downarrow 1_{\~C})$ is precisely a universal arrow making the diagram (\ref{eq:universal_triangle}) commute.

For example, if $F = 1_{\~C}$ is the identity on $\~C$, then asking for $(B, g)$ to be an initial object in the comma category $(X \downarrow 1_{\~C})$ is the same as asking for a unique map $u: B \to C$ such that
\[\begin{tikzcd}
X & {B} \\
& {C}
\arrow["g", from=1-1, to=1-2]
\arrow["{h}"', from=1-1, to=2-2]
\arrow["{u}", from=1-2, to=2-2, dashed]
\end{tikzcd}\]
commutes for all $B \in \~C$. 
\end{ex}

\begin{ex}
To see how comma categories arise in the context of adjoint functors, consider the adjunction
\[\begin{tikzcd}
            {\Set} & {\Vect_k}
            \arrow[""{name=0, anchor=center, inner sep=0}, "F", shift left=2, from=1-1, to=1-2]
            \arrow[""{name=1, anchor=center, inner sep=0}, "U", shift left=2, from=1-2, to=1-1]
            \arrow["\dashv"{anchor=center, rotate=-90}, draw=none, from=0, to=1]
        \end{tikzcd}\]
from Example \ref{ex:adjoints} (\ref{ex:adjoint_vspace}). For any set $S$, the vector space $F(S)$ can be defined by the following universal property:\\

``\emph{given a vector space $V$, any function $f: S \to V$ extends uniquely to a linear map $\bar{f}: F(S) \to V$}''\\

This is the universal property of $F(S)$ as one would commonly find it stated in a textbook on vector spaces, but it is worth noting that here the forgetful functor $U$ has been forgotten: strictly speaking, the universal property above should read ``$f: S \to U(V)$'' instead of ``$f: S \to V$''. The word ``extends'' refers to the embedding of $S$ into the (potentially larger) set $U(F(S))$, given by the unit of the adjunction:
\begin{align*}
    \eta_S: S &\to U(F(S)) \\
    s &\mapsto s
\end{align*}
Thus, the universal property above is more precisely stated as:\\
\begin{adjustwidth}{0.5cm}{}
``\emph{for any vector space $V$ and $f: S \to V$, there exists a unique linear map $\bar{f}: F(S) \to V$ such that the diagram}
\[\begin{tikzcd}
    S & {U(F(S))} \\
    & {U(V)}
    \arrow["{\eta_S}", from=1-1, to=1-2]
    \arrow["f"', from=1-1, to=2-2]
    \arrow["{U(\bar{f})}", dashed, from=1-2, to=2-2]
\end{tikzcd}\] 
\emph{commutes}''
\end{adjustwidth}
In the language of comma categories, the unit $\eta_S$ is initial in $(S \downarrow U)$. This turns out to be a property of any adjunction, and is explained explained in our next result.
\end{ex}

\begin{lem}
\label{lem:initial}
    Let $\begin{tikzcd}
    {\~A} & {\~B}
    \arrow[""{name=0, anchor=center, inner sep=0}, "F", shift left=2, from=1-1, to=1-2]
    \arrow[""{name=1, anchor=center, inner sep=0}, "G", shift left=2, from=1-2, to=1-1]
    \arrow["\dashv"{anchor=center, rotate=-90}, draw=none, from=0, to=1]
\end{tikzcd}$ be an adjunction. Then for any $A \in \~A$, the unit map $\eta_A: A \to FG(A)$ is initial in $(A \downarrow G)$.
\end{lem}

\begin{proof}
    Let $(B, f: A \to G(B))$ be an object in $(A \downarrow G)$. We must show that there is exactly one map from $(F(A), \eta_A)$ to $(B, f)$. Such a map is an arrow $t: F(A) \to B$ in $\~B$ such that 
\begin{equation}
\label{eq:adj_triangle}
    \begin{tikzcd}
    A & {G(F(A))} \\
    & {G(B)}
    \arrow["{\eta_A}", from=1-1, to=1-2]
    \arrow["f"', from=1-1, to=2-2]
    \arrow["{G(t)}", dashed, from=1-2, to=2-2]
\end{tikzcd}
\end{equation} 
commutes. But we have $\bar{t} = \eta_A G(t)$ by Lemma \ref{lem:transpose}, so we must have $f = \bar{t}$, and hence $\bar{f} = t$. Therefore, the diagram (\ref{eq:adj_triangle}) commutes if and only if $t = \bar{f}$, and so $\bar{f}$ is the unique map in $(A \downarrow G)$ making (\ref{eq:adj_triangle}) commute. This proves the claim.
\end{proof}

We can now introduce our third and final formulation of adjoint functors:

\begin{thm}
\label{thm:adjoint_initial}
Take categories and functors $\begin{tikzcd}
    {\~A} & {\~B}
    \arrow["F", shift left, from=1-1, to=1-2]
    \arrow["G", shift left, from=1-2, to=1-1]
\end{tikzcd}$. There is a one-to-one correspondence between:
\begin{enumerate}[label=(\roman*)]
    \item  adjunctions $F \dashv G$; 
    \item \label{thm:initial(ii)} natural transformations $\eta: 1_{\~A} \to FG$ such that $\eta_A: A \to G(F(A))$ is initial in $(A \downarrow G)$ for every $A  \in \~A$. 
\end{enumerate} 
\end{thm}

\begin{proof}
    We have already seen (Lemma \ref{lem:initial}) that an adjunction $F \dashv G$ gives rise to a natural transformation $\eta$ satisfying the second property. We now have to show that every $\eta$ satisfying property (\ref{thm:initial(ii)}) above is the unit of exactly one adjunction $F \dashv G$. By Theorem \ref{thm:adj_triangles}, an adjunction $F \dashv G$ is simply a pair $(\eta, \ve)$ of natural transformations satisfying the triangle identities. We will prove that for every $\eta$ satisfying property (\ref{thm:initial(ii)}), there exists a unique natural transformation $\ve: GF \to 1_{\~B}$ such that the pair $(\eta, \ve)$ satisfies the triangle identities. To that end, let $\eta: 1_{\~A} \to FG$ be a natural transformation with the property stated in (\ref{thm:initial(ii)}).

    For uniqueness, suppose that $\ve, \ve': GF \to 1_{\~B}$ are natural transformations such that both $(\eta, \ve)$ and $(\eta, \ve')$ satisfy the triangle identities. One of the triangle identites states that for all $B \in \~B$, the triangle 
    \begin{equation}
    \label{eq:initial_triangle}
    \begin{tikzcd}[column sep = large, row sep = large]
        {G(B)} & {G(F(G(B)))} \\
        & {G(B)}
        \arrow["{\eta_{G(B)}}", from=1-1, to=1-2]
        \arrow["{G(\ve_B)}", from=1-2, to=2-2]
        \arrow[Rightarrow, no head, from=1-1, to=2-2]
    \end{tikzcd}
    \end{equation}
    commutes. This tells us that in the comma category $(G(B) \downarrow G)$, $\ve_B$ is a map 
    \[ \left( F(G(B)), G(B) \To{\eta_{G(B)}} G(F(G(B))) \right) \to \left(B, G(B) \To{1} G(B) \right)\]
    The same is true of $\ve'_B$. But $\eta_{G(B)}$ is initial, so there is only one such map, so $\ve_B = \ve'_B$. Since this is true for all $B \in \~B$, we have $\ve = \ve'$, which proves uniqueness.

    We now prove that there exists a natural transformation $\ve: FG \to 1_{\~B}$ such that $(\eta, \ve)$ satisfies the triangle identities. For each $B \in \~B$, define the component $\ve_B: F(G(B)) \to B$ of $\ve$ at $B$ to be the unique map
    \[\left(F(G(B)), \eta_{G(B)}\right) \to \left(B, 1_{G(B)}\right)\]
    in $(G(B) \downarrow G)$ (so the triangle (\ref{eq:initial_triangle}) commutes). To show that $\ve$ defines a natural transformation from $FG$ to $1_{\~B}$, let $q: B \to B'$ be an arrow in $\~B$. We need to prove that the square
    \[\begin{tikzcd}[column sep = large]
    {F(G(B))} & {F(G(B'))} \\
    B & {B'}
    \arrow["{F(G(q))}", from=1-1, to=1-2]
    \arrow["{\ve_B}"', from=1-1, to=2-1]
    \arrow["{\ve_{B'}}", from=1-2, to=2-2]
    \arrow["q"', from=2-1, to=2-2]
    \end{tikzcd}\]
    commutes, i.e., $F(G(q))\ve_{B'} = \ve_{B} q$. To see that this is the case, consider the following pair of commutative diagrams:
    \[\begin{tikzcd}[column sep = large, row sep = large]
    {G(B)} & {G(F(G(B)))} \\
    & {G(B)} \\
    & {G(B')}
    \arrow["{\eta_{G(B)}}", from=1-1, to=1-2]
    \arrow[Rightarrow, no head, from=1-1, to=2-2]
    \arrow["{G(q)}"', from=1-1, to=3-2]
    \arrow["{G(\ve_B)}", from=1-2, to=2-2]
    \arrow["{G(q)}", from=2-2, to=3-2]
\end{tikzcd} \qquad \begin{tikzcd}[row sep = large, column sep = scriptsize]
    {G(B)} && {G(F(G(B)))} \\
    {G(B')} && {G(F(G(B')))} \\
    && {G(B')}
    \arrow["{\eta_{G(B)}}", from=1-1, to=1-3]
    \arrow["{G(q)}", from=1-1, to=2-1]
    \arrow["{G(q)}"', bend right = 80, looseness = 2, from=1-1, to=3-3]
    \arrow["{G(F(G(q)))}", from=1-3, to=2-3]
    \arrow["{\eta_{G(B')}}", from=2-1, to=2-3]
    \arrow[Rightarrow, no head, from=2-1, to=3-3]
    \arrow["{G(\ve_{B'})}", from=2-3, to=3-3]
\end{tikzcd}\]
These diagrams tell us that $\ve_B q$ and $F(G(q))$ are both maps $\eta_{G(B)} \to G(q)$ in $(G(B) \downarrow G)$, and since $\eta_{G(B)}$ is initial, they must be equal. Hence $\ve$ is a natural transformation. 

We have already observed one of the triangle identities (equation (\ref{eq:initial_triangle})) holds. The other triangle identity states that
\[\begin{tikzcd}[column sep = large, row sep = large]
    F(A) & F(G(F(A))) \\
    & F(A)
    \arrow["{F(\eta_A)}", from=1-1, to=1-2]
    \arrow[Rightarrow, no head, from=1-1, to=2-2]
    \arrow["\ve_{F(A)}", from=1-2, to=2-2]
\end{tikzcd}\]
commutes for all $A \in \~A$. To prove that this triangle commutes, we use a similar technique as before: we have commutative diagrams 
\[\begin{tikzcd}[column sep = large, row sep = large]
    A & {G(F(A))} \\
    & {G(F(A))}
    \arrow["{\eta_A}", from=1-1, to=1-2]
    \arrow["{\eta_A}"', from=1-1, to=2-2]
    \arrow["{G(1_{F(A)})}", from=1-2, to=2-2]
\end{tikzcd} \qquad  \begin{tikzcd}[row sep = large, column sep = scriptsize]
    A && {G(F(A))} \\
    {GF(A)} && {G(F(G(F(A))))} \\
    && {G(F(A))}
    \arrow["{\eta_A}", from=1-1, to=1-3]
    \arrow["{\eta_A}", from=1-1, to=2-1]
    \arrow["{\eta_A}"', bend right = 80, looseness = 2, from=1-1, to=3-3]
    \arrow["{GF(\eta_A)}", from=1-3, to=2-3]
    \arrow["{\eta_{G(F(A))}}", from=2-1, to=2-3]
    \arrow[Rightarrow, no head, from=2-1, to=3-3]
    \arrow["{G(\ve_{F(A)})}", from=2-3, to=3-3]
\end{tikzcd}\]
so $1_F(A)$ and $\ve_{F(A)} \circ F(\eta_{A})$ are both maps $\eta_A \to \eta_A$ in $(A \downarrow G)$, but since $\eta_A$ is initial, we have $\ve_{F(A)} \circ F(\eta_A) = 1_{F(A)}$, as desired.
\end{proof}

In the coming sections, we will meet the adjoint functor theorems, which give explicit conditions under which a functor has a left adjoint. An important starting point for these theorems is the following corollary:

\begin{cor}
    Let $G: \~B \to \~A$ be a functor. Then $G$ has a left adjoint if and only if for each $A \in \~A$, the category $(A \downarrow G)$ has an initial object.
\end{cor}

\begin{proof}
    Lemma \ref{lem:initial} proves the ``if'' direction. To prove ``only if'', choose for each $A \in \~A$ an initial object of $(A \downarrow G)$, call it $\left(F(A), \eta_A: A \to G(F(A))\right)$. For each arrow $f: A \to A'$ in $\~A$, let $F(f): F(A) \to F(A')$ be the unique map $\eta_A \to f\eta_A$ in $(A \downarrow G)$; that is, we have a commutative diagram
    \[\begin{tikzcd}[sep  = scriptsize]
        A && {G(F(A))} \\
        & {A'} \\
        && {G(F(A'))}
        \arrow["{\eta_A}", from=1-1, to=1-3]
        \arrow["f"', from=1-1, to=2-2]
        \arrow["{G(F(f))}", from=1-3, to=3-3]
        \arrow["{\eta_{A'}}"', from=2-2, to=3-3]
    \end{tikzcd}\]
    It is easy to see that $F$ defines a functor $\~A \to \~B$, and the diagram above tells us that $\eta$ is a natural transformation $1 \to FG$. Therefore $F$ is left adjoint to $G$ by Theorem \ref{thm:adjoint_initial}.
\end{proof}

\begin{prop}
Let$\begin{tikzcd}[sep = scriptsize]
    {\~C} & {\~D}
    \arrow[""{name=0, anchor=center, inner sep=0}, "F", shift left=2, from=1-1, to=1-2]
    \arrow[""{name=1, anchor=center, inner sep=0}, "G", shift left=2, from=1-2, to=1-1]
    \arrow["\dashv"{anchor=center, rotate=-90}, draw=none, from=0, to=1]
\end{tikzcd}$ be a pair of adjoint functors between locally small categories. Then $G$ is faithful if and only if the counit morphism $\ve_D: FG(D) \to D$ is an epimorphism for every object $D$ of $\~D$.
\end{prop}
\begin{proof}
    As $F \dashv G$, we have a natural isomorphism $\Hom_{\~C}(F(C), D) \cong \Hom_{\~C}(C, G(D))$. Denote by $\eta$ the unit of this adjunction. To prove ``if'', suppose $\ve_D: FG(D) \epi D$ is an epimorphism for all $D \in \~D$, and take arrows $f,g: D \to D'$ in $\~D$ with $G(f) = G(g)$. By naturality of the counit, we have $f \circ \ve_D = \ve_{D'} \circ FG(f)$ and $g \circ \ve_D = \ve_{D'} \circ FG(g)$. As $G(f) = G(g)$, we have $FG(f) = FG(g)$, hence
    \[f \circ \ve_D = \varepsilon_{D'} \circ FG(f) = \varepsilon_{D'} \circ FG(g) = g \circ \ve_D,\]
    and since $\ve_D$ is epi, we deduce $f = g$, so $G$ is faithful. To prove ``only if'', assume $G$ is faithful, and take any $D \in \~D$. Suppose we are given arrows $f,g: D \to D'$ such that $f \circ \ve_D = g \circ \ve_D$. Then $G(f \circ \ve_D) = G(g \circ \ve_D)$, so $G(f) \circ G(\ve_D) = G(g) \circ G(\ve_D)$. We have $G(\ve_D) \circ \eta_{G(D)} = 1_{G(D)}$ by the triangle identities (\ref{eq:triangle_ids}), so $G(\ve_D)$ has a right inverse and is therefore monic. As $G(f) \circ G(\ve_D) = G(g) \circ G(\ve_D)$, we have 
\begin{align*}G(f) \circ G(\varepsilon_D) \circ \eta_{G(D)} &= G(g) \circ G(\varepsilon_D) \circ \eta_{G(D)}\\G(f) \circ 1_{G(D)} &= G(g) \circ 1_{G(D)}\\G(f) &= G(g)\end{align*}
    therefore $f = g$ as $G$ is faithful, so $\ve_D$ is an epimorphism.
\end{proof}
\chapter{Representable functors \& Yoneda's lemma}

\section{Basic definitions and examples}

Let $\~C$ be a category and fix an object $A$ in $\~C$. For each $B \in \~C$, there is an associated set (or class) $\~C(A, B)$ of maps $A \to B$. Our first observation is that this assignment is functorial: given $B' \in \~C$ and a map $B \to B'$, we obtain a function $\~C(A, B) \to \~C(A, B')$. This is the main content of our first definition.

\begin{defn}
Suppose $\~C$ is a locally small category, so that $\~C(A, B)$ is an object in the category $\Set$ of sets. Then for each $A \in \~C$, we define a \emph{ covariant hom-functor}
    \[H^A = \~C(A, -): \~C \to \Set\]
    which sends an object $B \in \~C$ to the set $\~C(A,B)$, and an arrow $g: B \to B'$ in $\~C$ to the function
    \begin{align*}
    H^A(g) = \~C(A, g): \~C(A, B) &\to \~C(A, B')\\
                    (A \To{f} B) &\mapsto (A \To{gf} B')
    \end{align*}

    There is also a \emph{contravariant hom-functor}
    \[H_B = \~C(-, B): \~C^\op \to \Set\]
    which sends an object $A \in \~C$ to the set $\~C(A, B)$, and an arrow $k: A \to A'$ to the function
    \begin{align*}
        H_B(k) = \~C(k, B): \~C(A', B) &\to \~C(A, B)\\
                    (A' \To{f} B) &\mapsto (A \To{fk} B)
    \end{align*}
\end{defn}
The above definition gives us two functions, defined at $f: A' \to B$ by
    \[H^A (g) f = gf, \quad H_B(k)f = fk\]

    For any two arrows $g: A \to A'$ and $k: B \to B'$, we have a diagram
    \begin{equation}
    \label{eq:hom_square}
        \begin{tikzcd}
            {\~C(A',B)} & {\~C(A,B)} \\
            {\~C(A',B')} & {\~C(A,B')}
            \arrow["{H_B(k)}", from=1-1, to=1-2]
            \arrow["{H^{A'} (g) }"', from=1-1, to=2-1]
            \arrow["{H^A (g) }", from=1-2, to=2-2]
            \arrow["{H_{B'}(k)}"', from=2-1, to=2-2]
        \end{tikzcd}
    \end{equation}
    which commutes in set, since both paths send $f: A' \to B$ to $gfk: A \to B'$.


\begin{defn}
    Let $\~C$ be a locally small category. A functor $F: \~C \to \Set$ is \emph{representable} if $F \cong H^A$ for some $A \in \~C$. A \emph{representation} of $F$ is a choice of object $A \in \~C$ and an isomorphism between $H^A$ and $F$.

    Dually, a functor $G: \~C^\op \to \Set$ is \emph{representable} if $G \cong H_B$ for some $B \in \~C$. A \emph{representation} of $G$ is a choice of object $B \in \~C$ and an isomorphism between $H_B$ and $G$.
\end{defn}

Representable functors are sometimes called ``representables''. Note that only set-valued functors can be representable. 

\begin{ex}
    Consider $H^1: \Set \to \Set$, where $1$ is the one-element set. If $X$ is a set, a map $1 \to X$ is simply an element of $X$, we have
    \[H^1(X) \cong X\]
    for each $X \in \Set$. We can easily check that this isomorphism is natural in $X$, so $H^1$ is isomorphic to the identity functor $1_{\Set}$, and therefore $1_{\Set}$ is representable.
\end{ex}

\begin{ex}
    There is a functor $\ob: \Cat \to \Set$ sending a small category to its underlying set of objects, and this functor is representable. To see that this is the case, consider the terminal category $\mathbf{1}$ (with one object and one identity arrow). A functor from $\mathbf{1}$ to a category $\~C$ is the same as an object of $\~C$, so 
    \[H^1(\~C) \cong \ob \~C\]
    Again, it is easy to see that this isomorphism is natural in $\~C$, so indeed, we find that $\ob \cong \Cat(\mathbf{1}, -)$. Similarly, we can show that the functor $\mathrm{map}: \Cat \to \Set$ sending a small category to its set of maps is representable.
\end{ex}

\begin{ex}
    Let $G$ be a group, regarded as a one-object category. Recall that a functor $G \to \Set$ is simply a $G$-set: a set equipped with a left $G$-action. Since $G$ has only one object as a category, there is only one choice of representable (up to isomorphism). This unique representable is the so-called \emph{left regular representation} of $G$, which will be familiar to those readers with knowledge of representation theory. It is the underlying set of $G$ acted on by left multiplication.
\end{ex}

\begin{ex}
    Fix a field $k$ and vector spaces $U, V$, and $W$ over $k$. Then there is a functor
    \[\Bilin(U, V; -): \Vect_k \to \Set\]
    whose value $\Bilin(U, V; W)$ at $W \in \Vect_k$ is the set of bilinear maps $U \times V \to W$. It can be shown that this functor is representable; there is some $T \in \Vect_k$ such that
    \[\Bilin(U, V; W) \cong \Vect_k(T, W)\]
    Indeed, this $T$ is the tensor product $U \otimes V$.
\end{ex}

\begin{ex}
    There is a functor 
    \[\~P: \Set^\op \to \Set\]
    sending each set $X$ to its power set $\~P(X)$. It is defined on maps $f: X' \to X$ by $(\~P(f))(U) = f^{-1}(U)$ for all $U \in \~P(X)$ (here $f^{-1}(U)$ denotes the preimage of $U$: $f^{-1}(U) = \{x' \in X' : f(x') \in U\}$). Moreover, a subset is the same as a map into the two-point set $2$ (via the characteristic function which we met in Proposition \ref{prop:set_surjective_epi}), so indeed, we have $\~P \cong H_2$.
\end{ex}

\begin{ex}
    Similarly, there is a functor 
    \[\~O: \Top^\op \to \Set\]
    defined on objects $X$ by taking $\~O(X)$ to be the set of open subsets of $X$. If $S$ denotes the Sierpi\'nski space (the two-point space in which exactly one of the two singleton subsets is open), then continuous maps $X \to S$ correspond naturally to open subsets of $X$. Therefore $\~O \cong H_S$, and $\~O$ is representable.
\end{ex}

\begin{ex}
    We previously defined a functor $C: \Top^\op \to \Ring$, assigning to each space the ring of continuous real-valued functions on it. The composite functor 
    \[\Top^\op \To{C} \Ring \To{U} \Set\]
    is representable, since by definition, $U(C(X)) = \Top(X, \R)$ for topological spaces $X$.
\end{ex}

\begin{lem} 
Let $\begin{tikzcd}
    {\~A} & {\~B}
    \arrow[""{name=0, anchor=center, inner sep=0}, "F", shift left=2, from=1-1, to=1-2]
    \arrow[""{name=1, anchor=center, inner sep=0}, "G", shift left=2, from=1-2, to=1-1]
    \arrow["\dashv"{anchor=center, rotate=-90}, draw=none, from=0, to=1]
\end{tikzcd}$ be an adjunction between locally small categories, and let $A \in \~A$. Then the functor 
\[\~A(A, G(-)): \~B \to \Set\]
(that is, the composite $\~B \To{G} \~A \To{H^A} \Set$) is representable.
\end{lem}

\begin{proof}
    By hypothesis, we have
    \[\~A(A, G(B)) \cong \~B(F(A), B)\]
    for any $B \in \~B$. If we can show that this isomorphism is natural in $B \in \~B$, we'll have shown that $\~A(A, G(-)) \cong H^{F(A)}$, from which we may conclude that $\~A(A, G(-))$ is representable. To show that this is natural, we need to show that the isomorphism defined by the adjunction gives rise to a natural transformation from $\~A(A, G(-))$ to $\~B(F(A), -)$. So, let $q: B \to B'$ be an arrow in $B$. We must show that the square 
    \[\begin{tikzcd}
        {\~A(A, G(B))} && {\~B(F(A), B)} \\
        {\~A(A, G(B'))} && {\~B(F(A), B')}
        \arrow["\sim", from=1-1, to=1-3]
        \arrow["{H^A(Gq)}"', from=1-1, to=2-1]
        \arrow["{H^{F(A)}(q)}", from=1-3, to=2-3]
        \arrow["\sim"', from=2-1, to=2-3]
    \end{tikzcd}\]
    commutes, where the horizontal arrows are the bijections provided by the adjunction. Suppose we have an arrow $f: A \to G(B)$. Going right then down around the diagram above gives $q\bar{f}$, while going down then right gives $\overline{G(q)f}$. Therefore the rectangle above will commute if and only if $\overline{G(q)f} = q\bar{f}$. But this is just the naturality condition \ref{eq:bar_axiom1} (with $g = \bar{f}$), so we are done.
\end{proof}

\begin{prop}
    Any set-valued functor with a left adjoint is representable.
\end{prop}

\begin{proof}
    Let $G: \~C \to \Set$ be a functor with a left adjoint $F$. Write $1$ for the one-point set. Then 
    \[G(A) \cong \Set(1, G(A))\]
    naturally in $A \in \~C$, so by the Lemma above, $G$ is representable, indeed, $G \cong H^{F(1)}$.
\end{proof}

So far we have defined, for each $A \in \~C$, a functor $H^A \in \Fun(\~C, \Set)$. This functor tells us how the object $A$ sees the surrounding ``world'' of the category $\~C$.  Dually, the functor $H_A$ tells us not how $A$ sees other objects, but instead how $A$ is seen by other objects of $\~C$. However, there are potentially many objects in the category $\~C$, so it makes sense to ask about what happens to $H^A$ as $A$ varies. In some sense, every object in $\~C$ has a different view of this category, so as $A \in \~C$ varies, so does the overall view of $\~C$.  On the other hand, each object perceives the same world, so the family $(H^A)_{A \in \~C}$ of ``views'' has some consistency to it. This means that whenever there is an arrow between objects $A$ and $A'$, there is also an arrow between $H^A$ and $H^{A'}$.

In detail, a map $f: A' \to A$ induces a natural transformation 
\[\begin{tikzcd}
    {\~C} && \Set
    \arrow[""{name=0, anchor=center, inner sep=0}, "{H^A}", curve={height=-12pt}, from=1-1, to=1-3]
    \arrow[""{name=1, anchor=center, inner sep=0}, "{H^{A'}}"', curve={height=12pt}, from=1-1, to=1-3]
    \arrow["{\,H^f}", shorten <=3pt, shorten >=3pt, Rightarrow, from=0, to=1]
\end{tikzcd}\]
the $B$-component of which is given by
\begin{align*}
    H^A(B) = \~C(A, B) &\to H^{A'}(B) = \~C(A', B) \\
    g &\mapsto gf
\end{align*}
Note the reversed direction! Each functor $H^A$ is covariant, but they come together to form a contravariant functor. This is outlined in the following definition:
\begin{defn}
    Let $\~C$ be a locally small category. The functor 
    \[H^\bullet: \~C^\op \to \Fun(\~C, \Set)\]
    is defined on objects by $H^\bullet(A) = H^A$ and on arrows $f$ by $H^\bullet(f) = H^f$.
\end{defn}
We can do something similar for $H_A$: any map $f: A \to A'$ induces a natural transformation
\[\begin{tikzcd}
    {\~C^\op} && \Set
    \arrow[""{name=0, anchor=center, inner sep=0}, "{H_A}", curve={height=-12pt}, from=1-1, to=1-3]
    \arrow[""{name=1, anchor=center, inner sep=0}, "{H_{A'}}"', curve={height=12pt}, from=1-1, to=1-3]
    \arrow["{\,H_f}", shorten <=3pt, shorten >=3pt, Rightarrow, from=0, to=1]
\end{tikzcd}\]
whose component at $B \in \~C$ is 
\begin{align*}
    H_A(B) = \~C(B, A) &\to H_{A'}(B) = \~C(B, A') \\
    g &\mapsto fg
\end{align*}
This motivates the following definition:
\begin{defn}
    Let $\~C$ be a locally small category. The \emph{Yoneda embedding} of $\~C$ is the functor
    \[H_\bullet: \~C \to \Fun(\~C^\op, \Set)\]
    defined on objects $A$ by $H_\bullet(A) = H_A$ and maps $f$ by $H_\bullet(f) = H_f$.
\end{defn}

To summarize, here are the key definitions so far:
\begin{itemize}
    \item For each $A \in \~C$, we have a functor $H^A: \~C \to \Set$, and putting these all together, we get a functor $H^\bullet: \~C^\op \to \Fun(\~C,\Set)$.
    \item For each $A \in \~C$, we have a functor $H_A: \~C^\op \to \Set$, and putting these all together, we get a functor $H_\bullet: \~C \to \Fun(\~C^\op,\Set)$.
\end{itemize}

\begin{prop}
    The Yoneda embedding $H_\bullet$ is injective on isomorphism classes of objects. That is, if $\~C$ is a locally small category with $A, B \in \~C$ and $H_A \cong H_{B}$, then $A \cong B$.
\end{prop}

\begin{proof}
    Suppose that $H_A \cong H_{B}$. Then there exists a natural isomorphism $\alpha: H_A \isom H_{B}$ in $\Fun(\~C^\op, \Set)$. This means that for each $C \in \~C$, there is a map $\alpha_C: H_A(C) \to H_{B}(C)$ with an inverse $\overline{\alpha_C}: H_{B}(C) \to H_A(C)$. Now, for each $C \in \~C$, we have that $H_A(C) \cong H_{B}(C)$, i.e.,
    \[\Hom(C, A) \cong \Hom(C, B)\]
    In particular, we have that $\Hom(A, A) \cong \Hom(A, B)$ and $\Hom(B, A) \cong \Hom(B, B)$. We now need to define maps $f: A \to B$ and $g: B \to A$ and show that they are inverses of one another. Set $f = \alpha_A(1_A)$ and $g = \overline{\alpha_B}(1_B)$. Then the fact that the naturality square
    \[\begin{tikzcd}
    {H_A(A)} & {H_A(B)} \\
    {H_{B}(A)} & {H_{B}(B)}
    \arrow["{H_A(g)}", from=1-1, to=1-2]
    \arrow["{\alpha_A}"', from=1-1, to=2-1]
    \arrow["{\alpha_{B}}", from=1-2, to=2-2]
    \arrow["{H_{B}(g)}"', from=2-1, to=2-2]
\end{tikzcd}\]
commutes tells us that $fg = 1_{B}$, since chasing $1_A: A \to A$ around this square gives
\[\begin{tikzcd}
    {1_A} & g \\
    {\alpha_A(1_A) = f} & {\begin{tabular}{c} $\alpha_{B}(g) = 1_{B}$ \\ \hline $fg$ \end{tabular}}
    \arrow[maps to, from=1-1, to=1-2]
    \arrow[maps to, from=1-1, to=2-1]
    \arrow[maps to, from=1-2, to=2-2]
    \arrow[maps to, from=2-1, to=2-2]
\end{tikzcd}\]
Similarly, by chasing $1_B: B \to B$ around the naturality square
\[\begin{tikzcd}
    {H_B(B)} & {H_B(A)} \\
    {H_A(B)} & {H_A(A)}
    \arrow["{H_A(f)}", from=1-1, to=1-2]
    \arrow["{\overline{\alpha_B}}"', from=1-1, to=2-1]
    \arrow["{\overline{\alpha_A}}", from=1-2, to=2-2]
    \arrow["{H_B(f)}"', from=2-1, to=2-2]
\end{tikzcd}\]
we deduce that $gf = 1_A$, which completes the proof.
\end{proof}

Strictly speaking, it does not make much difference whether we work with $H^\bullet$ or $H_\bullet$; anything that is true of the one is also true of another thanks to duality. In a sense we will explain shortly, $H_\bullet$ ``embeds'' the category $\~C$ into the functor category $\Fun(\~C^\op, \Set)$. While this initially seems like a slightly strange choice of embedding, it is actually incredibly useful, since the category $\Fun(\~C^\op, \Set)$ has some nice properties that $\~C$ might not have. There is also a third functor which unifies the two hom-functors $H_A$ and $H^A$. It is the final functor we will introduce in this section.

\begin{defn}
\label{defn:hom_bifun}
    Let $\~C$ be a locally small category. The \emph{hom bifunctor} 
    \[\Hom_{\~C}: \~C^\op \times \~C \to \Set\]
    is defined by
    \[\begin{tikzcd}
    {(A, B)} & \mapsto & {\~C(A, B)} \\
    & \mapsto \\
    {(A', B')} & \mapsto & {\~C(A', B')}
    \arrow["g", shift left=2, from=1-1, to=3-1]
    \arrow["{g \circ - \circ f}", from=1-3, to=3-3]
    \arrow["f", shift left=2, from=3-1, to=1-1]
\end{tikzcd}\]
That is, $\Hom_\~C(A, B) = \~C(A, B)$, and $\Hom_\~C(A' \To{f} A, B \To{g} B')$ is equal to the function $f(-)g: \~C(A, B) \to \~C(A', B')$ whose value at $p \in \~C(A, B)$ is the composite $gpf: A' \to B'$ (i.e., $\left(\Hom_\~C(f, g)\right)(p) = gpf$ for all $p: A \to B$ in $\~C$)
\end{defn}

In his book \cite{leinster2016}, Leinster compares the functor $\Hom_\~C$ to the metric in a metric space $(X, d)$, writing that ``the existence of the functor $\Hom_\~C$ is something like the fact that for a metric space $(X, d)$, the metric itself is a continuous map $d: X \times X \to \R$'' (p. 91).
\begin{rmks}
\hfill
\begin{enumerate}[label=(\roman*)]
    \item In Example \ref{ex:adjoints}(\ref{ex:adjoints6}), we saw that there is an adjunction $(- \times B) \dashv \Set(B, -)$ of functors for any set $B$. In a similar manner, there is an adjunction $(-\times \~B) \dashv \Fun(\~B, -)$ of functors $\Cat \to \Cat$, i.e., a canonical bijection
    \[\Cat(\~A \times \~B, \~C) \cong \Cat(\~A, \Fun(\~B, \~C))\]
    for $\~A, \~B, \~C \in \Cat$. Under this bijection, the functors 
    \[\Hom_\~A : \~A^\op \times \~A \to \Set, \quad H^\bullet: \~A^\op \to \Fun(\~A, \Set)\]
    correspond to one another. So the functor $\Hom_\~A$ carries the same information as $H^\bullet$ (or $H_\bullet$), just presented differently. 

    \item The functor $\Hom_\~A$ also helps us to explain the naturality axiom in the definition of adjunction (Definition \ref{defn:adjunction}). Suppose we have categories and functors $\begin{tikzcd}
    {\~A} & {\~B}
    \arrow["F", shift left, from=1-1, to=1-2]
    \arrow["G", shift left, from=1-2, to=1-1]
\end{tikzcd}$. They give rise to functors 
\[\begin{tikzcd}
    {\~A^\op \times \~B} & {\~A^\op \times \~A} \\
    {\~B^\op \times \~B} & \Set
    \arrow["{1 \times G}", from=1-1, to=1-2]
    \arrow["{F^\op \times 1}"', from=1-1, to=2-1]
    \arrow["{\Hom_\~A}", from=1-2, to=2-2]
    \arrow["{\Hom_\~B}"', from=2-1, to=2-2]
\end{tikzcd}\]
Going right then down sends the pair $(A, B)$ to $\~B(F(A), B)$, and can therefore be written as $\~B(F(-), -)$. Going down then right sends $A, B$ to $\~A(A, G(B))$, and can be written as $\~A(-,G(-))$. These two functors are naturally isomorphic if and only if $F \dashv G$.
\end{enumerate}
\end{rmks}

\section{The Yoneda lemma}

Fix a locally small category $\~C$. Recall from Definition \ref{defn:presheaf} that a \emph{presheaf} on $\~C$ is a functor $F: \~C^\op \to \Set$. The functor category $\Fun(\~C^\op, \Set)$ is therefore called the \emph{category of presheaves on $\~C$}, henceforth denoted by $\pre{\~C}$. 

Previously, we have defined, for each $A \in \~C$, a representable presheaf $H_A: \~C^\op  \to \Set$. The question we aim to answer in this section is deceptively simple: what do arrows in the presheaf category $\pre{\~C}$ look like from the point of view of $H_A$? That is, if $X$ is another presheaf on $\~C$, what are the arrows $H_A \to X$? That is, what are the natural transformations $\begin{tikzcd}
    {\~C^\op} & \Set
    \arrow[""{name=0, anchor=center, inner sep=0}, "{H_A}", bend left, from=1-1, to=1-2]
    \arrow[""{name=1, anchor=center, inner sep=0}, "X"', bend right, from=1-1, to=1-2]
    \arrow[shorten <=2pt, shorten >=2pt, Rightarrow, from=0, to=1]
\end{tikzcd}$? We will denote the set of all such natural transformations by $\pre{\~C}(H_A, X)$. Our goal in this section is to study this set in greater detail. 


An informal principle in category theory that Leinster describes in his book \cite{leinster2016} states that, given any collection of input data of the same type, there ought to be exactly one output of another type. For instance, when we defined composition of arrows in a category, we remarked that for any collection of objects and arrows $A_0 \to A_1 \to A_2 \to \cdots \to A_n$, it is possible to construct exactly one map $A_0 \to A_n$. Now, consider the process we have just described for constructing the set $\pre{\~C}(H_A, X)$. The input data is an object $A \in \~C$ and a presheaf $X \in \pre{\~C}$, and the resulting output is a set $\pre{\~C}(H_A, X)$. Now, given the same input data of an object $A$ and a presheaf $X$, there is another natural way to construct a set, namely by evaluating the presheaf $X$ at $A$, which produces the set $X(A)$. Thus, given the input data $(X, A)$, we have found two different ways to construct a set. Therefore, if we keep the informal principle described above in mind, we might conjecture that
\[\pre{\~C}(H_A, X) \cong X(A).\]
This turns out to be true, and this is the statement of the Yoneda lemma. Informally, Yoneda says that for any $A \in \~C$ and any presheaf $X$ on $\~C$, a natural transformation $H_A \to X$ is the same as an element of $X(A)$. Formally, it can be stated as follows:
\begin{thm}[Yoneda]
\label{thm:yoneda}
    Let $\~C$ be a locally small category, then 
    \[\pre{\~C}(H_A, X) \cong X(A)\]
    naturally in $A \in \~C$ and $X \in \pre{\~C}$.
\end{thm}
\begin{proof}

We need to define a bijection of sets $\pre{\~C}(H_A, X) \cong X(A)$ and show that this bijection is natural in $A$ and $X$. First, fix $A \in \~C$ and $X \in \pre{\~C}$. We will define maps 
\[\begin{tikzcd}
    {\pre{\~C}(H_A,X) } & {X(A)}
    \arrow["{(\hat{~})}", shift left, from=1-1, to=1-2]
    \arrow["{(\tilde{~})}", shift left, from=1-2, to=1-1]
\end{tikzcd}\]
and show that they are inverses of one another. First, we define the map $(\hat{~}): \pre{\~C}(H_A,X) \to X(A)$ given $\alpha: H_A \to X$, define $\hat{\alpha} \in X(A)$ by $\hat{\alpha} = \alpha_A(1_A)$. To define the map $(\tilde{~}): X(A) \to \pre{\~C}(H_A,X)$, let $x \in X$. We need to specify a natural tranformtaion $\tilde{x}: H_A \to X$. That is, for each $B \in \~C$, we need to produce a map $\tilde{x}_B: H_A(B) \to X(B)$. Given $f: B \to A$  in $H_A(B)$, define $\tilde{x}_B(f)$ by $\tilde{x}_B(f) = (X(f))(x) \in X(B)$. To show that $\tilde{x}$ defines a natural transformation, we need to show that for any map $g: B' \to B$ in $\~C$, the square
\[\begin{tikzcd}
    {\~C(B, A)} && {\~C(B',A)} \\
    {X(B)} && {X(B')}
    \arrow["{H_A(g)= -\circ g}", from=1-1, to=1-3]
    \arrow["{\tilde{x}_B}"', from=1-1, to=2-1]
    \arrow["{\tilde{x}_{B'}}", from=1-3, to=2-3]
    \arrow["Xg"', from=2-1, to=2-3]
\end{tikzcd}\]
commutes. Indeed, if $f: B \to A$ is an arrow in $\~C(B, A)$, then we have
\[\begin{tikzcd}
    {(f: B \to A)} & {(f \circ g: B' \to A)} \\
    {Xf(x)} & {\begin{tabular}{c} $(Xfg)(x)$ \\ \hline $(Xg(Xf(x)))$ \end{tabular}}
    \arrow[maps to, from=1-1, to=1-2]
    \arrow[maps to, from=1-1, to=2-1]
    \arrow[maps to, from=1-2, to=2-2]
    \arrow[maps to, from=2-1, to=2-2]
\end{tikzcd}\]
and $X(f \circ g) = (Xg) \circ (Xf)$ since $X$ is a contravariant functor, so the diagram commutes. We now need to show that $(\hat{~})$ and $(\tilde{~})$ are inverses. First, take $x \in X(A)$. Then 
\[\hat{\tilde{x}} = \tilde{x}_A(1_A) = (X1_A)(x) = 1_{X(A)}(x) = x\]
so $(\hat{\tilde{~}})$ is the identity on $X(A)$. Next, take $\alpha: H_A \to X$. We must show that $\tilde{\hat{\alpha}} = \alpha$. Two natural transformations are equal if and only if they are equal componentwise, so we need to verify that
\[(\tilde{\hat{\alpha}})_B(f) = \alpha_B(f)\]
for all $B \in \~C$ and $f: B \to A$ in $\~C$. We have
\[(\tilde{\hat{\alpha}})_B(f) = (Xf)(\hat{\alpha}) = (Xf)(\alpha_A(1_A)),\]
so we need to show that 
\begin{equation}
    \label{eq:3.2.1}
    (Xf)(\alpha_A(1_A)) = \alpha_B(f).
\end{equation}
Indeed, naturality of $\alpha$ gives us the commutative square
\end{proof}
\[\begin{tikzcd}[sep=scriptsize]
    {\~C(A,A)} && {\~C(B,A)} \\
    {X(A)} && {X(B)}
    \arrow["{H_A(f) =-\circ f}", from=1-1, to=1-3]
    \arrow["{\alpha_A}"', from=1-1, to=2-1]
    \arrow["{\alpha_B}", from=1-3, to=2-3]
    \arrow["Xf"', from=2-1, to=2-3]
\end{tikzcd}\]
and chasing $1_A$ around this square gives condition \ref{eq:3.2.1}. Finally, we need to prove naturality. In principle, we need to check naturality of both $(\tilde{~})$ and $(\hat{~})$, but by \label{lem:nat_isom_condition} it suffices to check naturality of $\hat{~}$ alone. We'll show that $(\hat{~})$  


\section{Some consequences of the Yoneda lemma}
\begin{cor}
Let $\~C$ be a locally small category. Let $C$ and $D$ be objects of $\~C$. Show that $C$ and $D$ are isomorphic if and only if $H_C$ and $H_D$ are naturally isomorphic functors (dually, $C$ and $D$ are isomorphic if and only if $H^C$ and $H^D$ are naturally isomorphic functors)
\end{cor}
\begin{proof}
    To prove ``if'', suppose $C \cong D$ in $\~C$. Then there exists $f: C \to D$ and $f^{-1}:D \to C$ with $f^{-1} \circ f = \id_C$ and $f \circ f^{-1} = \id_D$. To define a natural isomorphism $h_C \xRightarrow{\sim} h_D$, we will define natural transformations $\alpha: h_C \Rightarrow h_D$ and $\beta: h_D \Rightarrow h_C$ with $\beta \circ \alpha = \id_{h_C}$ and $\alpha \circ \beta = \id_{h_D}$. For $A \in \~C$, we define 
    \begin{align*}
        \alpha_A: \Hom_{\~C}(A,C) &\to \Hom_{\~C}(A,D) \\ 
        s &\mapsto f \circ s
    \end{align*} 
    and
    \begin{align*}
        \beta_A: \Hom_{\~C}(A,D) &\to \Hom_{\~C}(A,C) \\ 
        t &\mapsto f^{-1} \circ t
    \end{align*}
    Then for $g: B \to A$ in $\~C$, we have a commutative diagram
    \[\begin{tikzcd}
        {\Hom_{\~C}(A,C)} & {\Hom_{\~C}(B, C)} \\
        {\Hom_{\~C}(A,D)} & {\Hom_{\~C}(B,D)}
        \arrow["{- \circ g}", from=1-1, to=1-2]
        \arrow["{\alpha_C}"', from=1-1, to=2-1]
        \arrow["{\alpha_D}", from=1-2, to=2-2]
        \arrow["{- \circ g}"', from=2-1, to=2-2]
    \end{tikzcd}\]
    since both paths send $s \in \Hom_{\~C}(A,C)$ to $f \circ s \circ g \in \Hom_{\~C}(A,C)$. This shows that $\alpha$ is a natural transformation $H_C \Rightarrow H_D$, and using the analogous argument with $f^{-1}$ instead of $f$, we see that $\beta$ defines a natural transformation $H_D \Rightarrow H_C$. Then for any $A \in \~C$ and $t \in \Hom_{\~C}(B, D)$, we have
    \begin{align*}
        (\alpha_A \circ \beta_A)(t) &= \alpha_A(\beta_A(t)) \\ &= \alpha_A (f^{-1} \circ t) \\ &= f \circ f^{-1} \circ t  \\ &= t.
    \end{align*}
    and similarly, $(\beta_B\circ \alpha_B)(s) = s$ for any $s \in \Hom(A, C)$. Hence $\beta \circ \alpha = \id_{h_C}$ and $\alpha \circ \beta = \id_{h_D}$, which settles the forward implication.

    For the reverse implication, we will require the aid of Yoneda's lemma. Suppose $H_C \cong H_D$ in the functor category $\Fun \left(\~C^\op, \Set \right)$. Yoneda tells us that
    \[\widehat{C}(H_C, H_D) \cong H_D(C) = \Hom_{\~C}(C,D),\]
    so we have a natural isomorphism $\phi: H_C \Rightarrow H_D$. Under the bijection above, $\phi$ corresponds to a unique arrow $f = \phi_C(\id_C) \in \Hom_{\~C}(C, D)$. Being an isomorphism, $\phi$ has an inverse $\psi: H_D \Rightarrow H_C$, which will correspond to a unique morphism $g = \psi_D(\id_D) \in \Hom_{\~C}(D,C)$. We have $\psi \circ \phi = \id_{H_C}$, hence $(\psi_C \circ \phi_C)(\id_C) = \id_C$. But we also know that $\phi_C(\id_C) = f$ and $\psi_C(f) = g \circ f$ thanks to the commutativity of the diagram
    \[\begin{tikzcd}
    {H_D(D)} & {H_D(C)} \\
    {H_C(D)} & {H_C(C)}
    \arrow["{- \circ f}", from=1-1, to=1-2]
    \arrow["{\Psi_D}"', from=1-1, to=2-1]
    \arrow["{\Psi_C}", from=1-2, to=2-2]
    \arrow["{- \circ f}"', from=2-1, to=2-2]
\end{tikzcd}\]
Therefore $\id_C = \psi_C(\phi_C(\id_C)) = \psi_C(f) = g \circ f$. Repeating this logic for the composition $\phi \circ \psi$ at the object $D$, we get $\phi_D(g) = f \circ g = \text{id}_D$, so $C \cong D$ in $\~C$. 
\end{proof}

\chapter{Limits}

Limits are a fundamental concept in category theory, underlying numerous constructions throughout mathematics. At the heart of limits is the idea of a \emph{universal property}, which is essentially a statement of the form
\begin{center}
    \emph{``for all $X$ in a category there exists a unique morphism making $Y$ occur''}
\end{center}
these kinds of statements are ubiquitous in category theory, and many of them underlie important constructions in other disciplines. Indeed, we have already seen a very basic example of a universal property when we met terminal and initial objects (for the rest of this discussion, we'll focus on terminal objects without loss of generality). Recall that if $1$ is terminal in the category $\~C$, there is exactly one map $X \to 1$ for all objects $X \in \~C$. Another way we can phrase this is 
\begin{center}
    \emph{``for all $X \in \~C$ there exists a unique morphism $X \to 1$''}
\end{center}
so we can see that terminal objects have their own flavour of universal property (again, this applies by duality to the notion of initial objects). 

However, when we met terminal objects, we also saw more; proposition (\ref{prop:terminal_objects_are_isomorphic}) tells us that terminal objects are determined \emph{up to a unique isomorphism}, meaning that if $1$ and $1'$ are both terminal objects, then there is a unique isomorphism $1 \isom 1'$. The fact that we can establish such an isomorphism in the first place is because of the universal property associated with terminal objects. In general, every universal property determines some object in a category up to a unique isomorphism, which is an incredibly powerful property. In this section, we will meet more sophisticated examples of universal properties and see how they underlie various fundamental notions in category theory and beyond. Our end goal will be the definition of a \emph{limit}, but before we arrive at this definition we will introduce three important kinds of limits: products, equalizers, and pullbacks. With these objects defined, we will be better equipped to treat the more general definition limits. 


\section{Products}

Our first kind of limit is the \emph{product}. First, some basic terminology: a \emph{span} in a category $\~C$ consists of objects $X, Y, Z \in \~C$ and maps
\[\begin{tikzcd}[column sep=small]
            & {X} \arrow[dl, swap] \arrow[dr] & \\
            Y & & Z
        \end{tikzcd}\]
and a \emph{cospan} is a span in $\~C^\op$, i.e., a span with the arrows reversed:
\[\begin{tikzcd}[column sep=small]
    & X \\
    Y && Z
    \arrow[from=2-1, to=1-2]
    \arrow[from=2-3, to=1-2]
\end{tikzcd}\]
\begin{defn}
    Let $\~C$ be a category and $X, Y \in \~C$. A \emph{product} of $X$ and $Y$ in $\~C$ consists of an object $P$ in $\~C$ and a span
        \[\begin{tikzcd}[column sep=small]
            & {P} \arrow[dl, "\pi_X", swap] \arrow[dr, "\pi_Y"] & \\
            X & & Y
        \end{tikzcd}\]
    With the property that for every other span
        \[\begin{tikzcd}[column sep=small]
            & A \arrow[dl, "f", swap] \arrow[dr, "g"] & \\
            X & & Y
        \end{tikzcd}\]
    in $\~C$, there exists a unique map $\langle f, g \rangle: A \to P$ such that
    \begin{equation}
    \label{eq:product}
          \begin{tikzcd}
            & A \\
            \\
            & {P} \\
            X && Y
            \arrow["f"', curve={height=6pt}, from=1-2, to=4-1]
            \arrow["{\pi_X}", from=3-2, to=4-1]
            \arrow["g", curve={height=-6pt},from=1-2, to=4-3]
            \arrow["{\pi_Y}"', from=3-2, to=4-3]
            \arrow["{\langle f, g\rangle}"{description}, dashed, from=1-2, to=3-2]
        \end{tikzcd}
    \end{equation}
commutes. The maps $\pi_X$ and $\pi_Y$ are called \emph{projections}. We typically write $X \times Y$ for the product object $P$ in the above definition. A \emph{coproduct} in the category $\~C$ is a product in $\~C^\op$. Explicitly, this means a coproduct of $X$ and $Y$ as above consists of an object $C$ and maps 
\[\begin{tikzcd}[column sep=small]
    & C \\
    X && Y
    \arrow["{\pi_X}", from=2-1, to=1-2]
    \arrow["{\pi_Y}"', from=2-3, to=1-2]
\end{tikzcd}\]
with the property that for all objects and maps 
\[\begin{tikzcd}[column sep=small]
    & A \\
    X && Y
    \arrow["{f}", from=2-1, to=1-2]
    \arrow["{g}"', from=2-3, to=1-2]
\end{tikzcd}\]
there exists a unique map $\langle f, g \rangle  : C \to A$ such that
\[\begin{tikzcd}
    & A \\
    \\
    & C \\
    X && Y
    \arrow["{\pi_X}"', from=4-1, to=3-2]
    \arrow["{\pi_Y}", from=4-3, to=3-2]
    \arrow["f", curve={height=-6pt}, from=4-1, to=1-2]
    \arrow["g"', curve={height=6pt}, from=4-3, to=1-2]
    \arrow["{\langle f, g\rangle }"{description}, dashed, from=3-2, to=1-2]
\end{tikzcd}\]
commutes. In the case of coproducts, the maps $\pi_X$ and $\pi_Y$ are called \emph{coprojections}.
\end{defn}

\begin{ex}[Products and coproducts in $\Set$]
    In $\Set$, the cartesian product $X \times Y$ is a product. To see this, we first define projection maps 
    \[\begin{tikzcd}[column sep=small]
    & {X \times Y} \\
    X && Y
    \arrow["{\pi_X}"', from=1-2, to=2-1]
    \arrow["{\pi_Y}", from=1-2, to=2-3]
    \end{tikzcd} \quad \begin{tikzcd}
    & {(x, y)} \\
    x && y
    \arrow[maps to, from=1-2, to=2-1]
    \arrow[maps to, from=1-2, to=2-3]
\end{tikzcd}\]
Now let 
    \[\begin{tikzcd}[column sep=small]
    & A \\
    X && Y
    \arrow["f"', from=1-2, to=2-1]
    \arrow["g", from=1-2, to=2-3]
    \end{tikzcd}\]
be a span in $\Set$. We want a unique dotted arrow $\langle f,g \rangle$ making the diagram (\ref{eq:product}) commute. If we define
    \[\langle f, g \rangle (a) := (f(a), g(a)) \in X \times Y\]
for all $a \in A$, then (\ref{eq:product}) clearly commutes. To see this map is unique, suppose we had another map $\overline{\langle f, g \rangle}$ filling in the dotted arrow in (\ref{eq:product}). Then if $\overline{\langle f, g \rangle}(a) = (x, y) \in X \times Y$, we require that
    \[\pi_X \left((x, y) \right) = x = f(a) \quad \text{and} \quad \pi_Y \left((x, y) \right) = y = g(a)\]
and therefore
\[\overline{\langle f, g \rangle}(a) = \left(f(a), g(a)\right) = \langle f, g \rangle(a)\]
for all $a \in A$, so the map $\langle f, g \rangle$ is unique. Thus, we have produced projections $\pi_X$ and $\pi_Y$ and a unique map $\langle f, g \rangle$ making (\ref{eq:product}) commute, so we can conclude that $X \times Y$ is a valid product in $\Set$.

Recall that the union operation of sets $X \cup Y$ does not distinguish between elements from $X$ and elements in $Y$. For instance, if $X = \{a, b\}$ and $Y = \{a, c, d\}$ then $X \cup Y = \{a, b, c, d\}$ -- the union operation ``cannot tell'' that $a$ was common to both $X$ and $Y$. With this in mind, we can now introduce the coproduct in $\Set$, which is given by the \emph{disjoint union of sets} $X \sqcup Y$. This is like the union operation $X \cup Y$, but where we distinguish between elements that come from $X$ and elements that come from $Y$. There are several equivalent ways to define the disjoint union (indeed, one is as the coproduct in $\Set$), but perhaps the simplest is via the cartesian product:
\[X \sqcup Y = \left(  \{0\}  \times X \right) \cup \left( \{1\} \times Y\right)\]
So, if we again take $X = \{a, b\}$ and $Y = \{a, c, d\}$, we can compute their disjoint union as
\[X \sqcup Y = \{(0, a), (0, b), (1, a), (1, c), (1, d)\}\]
Thus, $X \sqcup Y$ takes the two sets $X$ and $Y$, makes them disjoint from one another, and then takes the union of these two sets, hence the name ``disjoint union''. To show that this defines a coproduct, we define projection maps 
\[\begin{tikzcd}[column sep=small]
    & {X \sqcup Y} \\
    X && Y
    \arrow["{\pi_X}", to=1-2, from=2-1]
    \arrow["{\pi_Y}"', to=1-2, from=2-3]
    \end{tikzcd}\]
by $\pi_X(x) = (0, x)$ and $\pi_Y(y) = (1, y)$. It is now straightforward to check that $\sqcup$ satisfies the universal property of the coproduct. 
\end{ex}

It is important to note that products do not exist in every category. However, we will now see that when products do exist, they are determined up to a unique isomorphism via the universal property.

\begin{prop}
    If $P$ and $Q$ are two products (resp. coproducts) in a category $\~C$, then $P \cong Q$, and moreover, this isomorphism is unique.
\end{prop}
\begin{proof}
    Suppose $P$ and $Q$ are both products in $\~C$, with projections
    \[\begin{tikzcd}[column sep=small]
    & P \\
    X && Y
    \arrow["{\pi_X}"', from=1-2, to=2-1]
    \arrow["{\pi_Y}", from=1-2, to=2-3]
    \end{tikzcd} \quad \text{and} \quad \begin{tikzcd}
    & Q \\
    X && Y
    \arrow["{\rho_X}"', from=1-2, to=2-1]
    \arrow["{\rho_Y}", from=1-2, to=2-3]
    \end{tikzcd}\]
    Then we can construct the diagram 
    \[\begin{tikzcd}
        & P \\
        \\
        & Q \\
        X && Y
        \arrow["{\pi_X}"', from=1-2, to=4-1]
        \arrow["{\rho_X}", from=3-2, to=4-1]
        \arrow["{\pi_Y}", from=1-2, to=4-3]
        \arrow["{\rho_Y}"', from=3-2, to=4-3]
    \end{tikzcd}\]
    and by the universal property of products, there is a unique map $\langle \pi_X, \pi_Y \rangle: P \to Q$ making
    \[\begin{tikzcd}
        & P \\
        \\
        & Q \\
        X && Y
        \arrow["{\pi_X}"', curve={height=12pt}, from=1-2, to=4-1]
        \arrow["{\rho_X}", from=3-2, to=4-1]
        \arrow["{\pi_Y}", curve={height=-12pt}, from=1-2, to=4-3]
        \arrow["{\rho_Y}"', from=3-2, to=4-3]
        \arrow["{\langle \pi_X, \pi_Y \rangle}"{description}, dashed, from=1-2, to=3-2]
    \end{tikzcd}\]
    commute. However we can also form the diagram
    \[\begin{tikzcd}
        & Q \\
        \\
        & P \\
        X && Y
        \arrow["{\rho_X}"', from=1-2, to=4-1]
        \arrow["{\pi_X}", from=3-2, to=4-1]
        \arrow["{\rho_Y}",  from=1-2, to=4-3]
        \arrow["{\pi_Y}"', from=3-2, to=4-3]
    \end{tikzcd}\]
and again, we know a unique map $\langle \rho_X, \rho_Y \rangle: Q \to P$ makes
    \[\begin{tikzcd}
        & Q \\
        \\
        & P \\
        X && Y
        \arrow["{\rho_X}"', curve={height=12pt}, from=1-2, to=4-1]
        \arrow["{\pi_X}", from=3-2, to=4-1]
        \arrow["{\rho_Y}", curve={height=-12pt}, from=1-2, to=4-3]
        \arrow["{\pi_Y}"', from=3-2, to=4-3]
        \arrow["{\langle \rho_X, \rho_Y \rangle}"{description}, dashed, from=1-2, to=3-2]
    \end{tikzcd}\]
commute. But now we have a commutative diagram
\[\begin{tikzcd}
    & P \\
    \\
    & Q \\
    \\
    & P \\
    X && Y
    \arrow["{\langle \pi_X, \pi_Y \rangle}"{description}, dashed, from=1-2, to=3-2]
    \arrow["{\langle \rho_X, \rho_Y \rangle}"{description}, dashed, from=3-2, to=5-2]
    \arrow["{\pi_X}", from=5-2, to=6-1]
    \arrow["{\pi_Y}"', from=5-2, to=6-3]
    \arrow["{\rho_X}", curve={height=12pt}, from=3-2, to=6-1]
    \arrow["{\rho_Y}"', curve={height=-12pt}, from=3-2, to=6-3]
    \arrow["{\pi_X}"', curve={height=24pt}, from=1-2, to=6-1]
    \arrow["{\pi_Y}", curve={height=-24pt}, from=1-2, to=6-3]
\end{tikzcd}\]
And we see that the composite arrow $\langle \rho_X, \rho_Y \rangle \langle \pi_X, \pi_Y \rangle : P \to P$ makes the reduced diagram
\[\begin{tikzcd}
    & P \\
    \\
    \\
    & P \\
    X && Y
    \arrow["{\pi_X}", from=4-2, to=5-1]
    \arrow["{\pi_Y}"', from=4-2, to=5-3]
    \arrow["{\pi_X}"', curve={height=30pt}, from=1-2, to=5-1]
    \arrow["{\pi_Y}", curve={height=-30pt}, from=1-2, to=5-3]
    \arrow["{\langle \rho_X, \rho_Y \rangle \langle \pi_X, \pi_Y \rangle}"{description}, dashed, from=1-2, to=4-2]
\end{tikzcd}\]
commute. However, we can think of another map making this diagram commute, namely $1_P: P \to P$. Since $P$ is a product, the map $\langle \rho_X, \rho_Y \rangle \langle \pi_X, \pi_Y \rangle$ is unique and thus
\[\langle \rho_X, \rho_Y \rangle \langle \pi_X, \pi_Y \rangle = 1_P\]

And that $\langle \rho_X, \rho_Y \rangle \langle \pi_X, \pi_Y\rangle = 1_Q$ follows, \emph{mutatis mutandis}, by interchanging $P$ and $Q$ in the above argument. Hence, there is a unique isomorphism $\langle \pi_X, \pi_Y \rangle: P \isom Q$, as desired. Dualizing the above argument, we see that when $P$ and $Q$ are coproducts there is a unique isomorphism $Q \isom P$, which completes the proof. 
\end{proof}

Therefore, we can speak of \emph{the} product in a category $\~C$, since any two products in $\~C$ will always be isomorphic. 

\begin{prop}
    Let $\~C$ be a category with a terminal object $1$. Then $X$ is a product of $X$ and $1$ in $\~C$. In symbols, $X \times 1 = X$.
\end{prop}

\begin{ex}
In the category $\Vect_k$ of vector spaces, the product is given by the direct sum $\oplus$. If $V$ and $W$ are vector spaces over the same field $k$, then $V \oplus W$ has elements $(v, w)$ (with $v \in V$ and $w \in W$). We therefore have linear projection maps
\[\begin{tikzcd}[column sep=small]
    & V \oplus W \\
    V && W
    \arrow["{\pi_V}"', from=1-2, to=2-1]
    \arrow["{\pi_W}", from=1-2, to=2-3]
    \end{tikzcd} \quad \begin{tikzcd}[column sep=small]
    & {(v, w)} \\
    v && w
    \arrow["{\pi_V}"', maps to, from=1-2, to=2-1]
    \arrow["{\pi_W}", maps to, from=1-2, to=2-3]
\end{tikzcd}\]
If we now take a vector space $X$ over $k$ with maps
\[\begin{tikzcd}[column sep=small]
    & X \\
    V && W
    \arrow["{f}"', from=1-2, to=2-1]
    \arrow["{g}", from=1-2, to=2-3]
    \end{tikzcd}\]
then we can indeed verify that $V \oplus W$ satisfies the universal property of the product.
\end{ex}

\begin{ex}[Posetal categories]
\hfill
\begin{enumerate}[label=(\roman*)]
    \item In the posetal category $(\R, \leq)$, the product of two real numbers $x$ and $y$ is their \emph{minimum} $\min\{x, y\}$. If $x \leq y$ in $(\R, \leq)$, we will draw an arrow $x \to y$ (we won't give arrows in posetal categories names, since this could lead to confusion). Now, $\min\{x, y\}$ satisfies $\min\{x, y\} \leq x$ and $\min\{x, y\} \leq y$, giving us projection maps:
\[\begin{tikzcd}[column sep=small]
    & {\min\{x, y\}} \\
    x && y
    \arrow[from=1-2, to=2-1]
    \arrow[from=1-2, to=2-3]
\end{tikzcd}\]
Furthermore, $\min\{x, y\}$ has the property that if $a \leq x$ and $a \leq y$, then $a \leq \min\{x, y\}$. This means that we have the diagram
\[
  \begin{tikzcd}
    & a \\
    \\
    & {\min\{x, y\}} \\
    x && y
    \arrow[from=1-2, to=4-1]
    \arrow[from=3-2, to=4-1]
    \arrow[from=1-2, to=4-3]
    \arrow[from=3-2, to=4-3]
    \arrow[from=1-2, to=3-2]
\end{tikzcd}
\]
and since all arrows are unique in $(\R, \leq)$, this commutes. We can therefore see that $\min\{x, y\}$ is the product of $x$ and $y$ in $(\R, \leq)$.

On the other hand, the coproduct of two real numbers $x$ and $y$ in $(\R, \leq)$ is their \emph{maximum} $\max\{x, y\}$, which satisfies $x \leq \max\{x, y\}$ and $y \leq \max\{x, y\}$ and has the additional property that whenever $x \leq a$ and $y \leq a$, it follows that $\max\{x, y\} \leq a$. This tells us the coproduct diagram
\[\begin{tikzcd}
    & a \\
    \\
    & {\max\{x, y\}} \\
    x && y
    \arrow[from=4-1, to=3-2]
    \arrow[from=4-3, to=3-2]
    \arrow[from=3-2, to=1-2]
    \arrow[from=4-1, to=1-2]
    \arrow[from=4-3, to=1-2]
\end{tikzcd}\]
commutes in $(\R, \leq)$.

\item Now fix a set $S$, and consider the poset $(\~P(S), \subseteq)$, regarded as a 
category. If $X, Y \in \~P(S)$, then the intersection $X \cap Y$ satisfies $X \cap Y \subseteq X$ and $X \cap Y \subseteq Y$, and has the additional property that if $A \subseteq X$ and $A \subseteq Y$, then $A \subseteq X \cap Y$. If we draw an arrow $A \to B$ when $A \subseteq B$ in $(\~P(S), \subseteq)$, we get the diagram 
\[
  \begin{tikzcd}
    & A \\
    \\
    & {X \cap Y} \\
    X && Y
    \arrow[from=1-2, to=4-1]
    \arrow[from=3-2, to=4-1]
    \arrow[from=1-2, to=4-3]
    \arrow[from=3-2, to=4-3]
    \arrow[from=1-2, to=3-2]
\end{tikzcd}
\]
and again, since every arrow in a posetal category is unique, this diagram commutes, so the product of $X$ and $Y$ in $(\~P(S), \subseteq)$ is $X \cap Y$. 

As one might guess, the coproduct in of $X$ and $Y$ in $(\~P(S), \subseteq)$ is their union $X \cup Y$. This is because $X \cup Y$ satisfies $X \subseteq X \cup Y$ and $Y \subseteq X \cup Y$, with the additional property that if $X \subseteq A$ and $Y \subseteq A$ then $X \cup Y \subseteq A$. In other words, the following coproduct diagram commutes in $(\~P(S), \subseteq)$:
\[\begin{tikzcd}
    & A \\
    \\
    & {X \cup Y} \\
    X && Y
    \arrow[from=4-1, to=3-2]
    \arrow[from=4-3, to=3-2]
    \arrow[from=3-2, to=1-2]
    \arrow[from=4-1, to=1-2]
    \arrow[from=4-3, to=1-2]
\end{tikzcd}\]

\item Let $x, y \in \N$. The greatest common divisor $\gcd(x, y)$ of $x$ and $y$ clearly has the property that $\gcd(x, y) ~|~ x$ and $\gcd(x, y) ~|~ y$. Furthermore, if $a~|~x$ and $a~|~y$, then $a ~|~\gcd(x, y)$. From this, we deduce that the product diagram
\[
  \begin{tikzcd}
    & a \\
    \\
    & {\gcd(x, y)} \\
    x && y
    \arrow[from=1-2, to=4-1]
    \arrow[from=3-2, to=4-1]
    \arrow[from=1-2, to=4-3]
    \arrow[from=3-2, to=4-3]
    \arrow[from=1-2, to=3-2]
\end{tikzcd}
\]
commutes in the posetal category $(\N, ~|~)$, so the product of two $x$ and $y$ in this category is their greatest common divisor.  

The coproduct in $(\N, ~|~)$ is given by the \emph{least common multiple} $\lcm(x, y)$ of $x$ and $y$, which is defined as the smallest positive integer with the property that $x ~|~ \lcm(x, y)$ and $y ~|~ \lcm(x, y)$. Furthermore, $\lcm(x, y)$ has the property that whenever $x ~|~ a$ and $y ~|~ a$ we have $\lcm(x, y) ~|~ a$. Hence, the diagram 
\[\begin{tikzcd}
    & a \\
    \\
    & {\lcm(x, y)} \\
    x && y
    \arrow[from=4-1, to=3-2]
    \arrow[from=4-3, to=3-2]
    \arrow[from=3-2, to=1-2]
    \arrow[from=4-1, to=1-2]
    \arrow[from=4-3, to=1-2]
\end{tikzcd}\]
commutes, and we see that $\lcm(x, y)$ satisfies the definition of the coproduct in $(\N, ~|~)$.
\end{enumerate}
\end{ex}
Let's now describe the our observations in the above example in a more general setting. Let $(P, \leq)$ be  a poset and take $x, y \in P$. A \emph{lower bound} for $x$ and $y$ is an element $a \in P$ such that $a \leq x$ and $a \leq y$. A \emph{greatest lower bound} or \emph{meet} of $x$ and $y$ is a lower bound $a$ for $x$ and $y$ with the further property that whenever $b$ is another lower bound for $x$ and $y$ we have $b \leq a$. When any poset $P$ is regarded as a category, the product of two objects $x$ and $y$ in this category (assuming it exists) will always be their meet. We typically use the notation $x \wedge y$ for the meet of $x$ and $y$ instead of the more traditional product notation $x \times y$. So in the previous example
\[x \wedge y = \min\{x, y\} \quad X \wedge Y = X \cap Y  \quad x \wedge y =  \gcd(x, y) \]
with the second example being the origin of the notation.

On the other hand, an \emph{upper bound} for $x$ and $y$ is an element $a \in P$ such that $x \leq a$ and $y \leq a$. A \emph{least upper bound} or \emph{join} of $x$ and $y$ is a lower bound $a$ for $x$ and $y$ with the further property that whenever $b$ is another lower bound for $x$ and $y$ we have $a \leq b$. When the poset $P$ is regarded as a category, the coproduct of two objects $x$ and $y$ in this category (when it exists) will always be their join. We denote the join of $x$ and $y$ by $x \vee y$. We can again rewrite the products in our previous example using this join notation:
\[x \vee y = \max\{x, y\} \quad X \vee Y = X \cup Y \quad x \vee y = \lcm(x, y)\]

We have so far only discussed \emph{binary} products and coproducts, which are products/coproducts of \emph{two} elements, like $X \times Y$. However, we can quite easily generalize this notion to products/coproducts $X \times Y \times Z$ of three elements, or even infinitely many elements. We just generalize our previous definition slightly:

\begin{defn}
\label{defn:gen_prod}
If $\{X_i\}_{i \in I}$ is an indexed family of objects of $\~C$, a \emph{product} of all the $X_i$'s together consists of an object $P$ in $\~C$ and maps
\[\left(P \To{\pi_i} X_i \right)_{i \in I}\]
    With the property that for all objects and maps
\[\left(A \To{f_i} X_i \right)_{i \in I}\]
in $\~C$, there exists a unique map $(f_i)_{i \in I}: A \to P$ such that $\pi_i \circ (f_i)_{i \in I} = f_i$ for all $i \in I$. Reversing the arrows in this definition, we get the definition of a coproduct of $\{X_i\}_{i \in I}$.
\end{defn}
We typically denote the product $P$ in the above definition by $\prod_{i \in I} X_i$, and for coproducts we use the notation $\coprod_{i \in I} X_i$. The map $(f_i)_{i \in I}$ is called the \emph{product of the maps $f_i: A \to X_i$}, and is defined by $(f_i)_{i \in I}(a) = f_i(a)$ for each $i \in I$ (thus, we can think of $(f_i)_{i \in I}$ as a function with components $f_i$ indexed by $i \in I$. This map is somewhat ubiquitous in mathematics: whenever we have a collection of maps $f_i: A \to X_i$ for all $i$ in some indexing set $I$, it is likely $(f_i)_{i \in I}$ will pop up at some point.  

\begin{ex}
In $\Set$, the product $\prod_{i \in I} X_i$ of a family of sets $\{X_i\}_{i \in I}$ is just the Cartesian product of all the $X_i$'s, and the coproduct of such a family is the disjoint union $\bigsqcup_{i \in I} X_i$, whose definition is a natural extension of the disjoint union of two sets:
\[\bigsqcup_{i \in I} X_i = \bigcup_{i \in I} \{i\} \times X_i\]
\end{ex}
\begin{ex}
In a posetal category $(P, \leq)$, a \emph{lower bound} for a family of elements $\{x_i\}_{i \in I}$ is an element $a \in P$ such that $a \leq x_i$ for all $i$. A \emph{greatest lower bound} of \emph{meet} of the family $\{x_i\}_{i \in I}$ is a lower bound for the family that is greater than any other lower bound. We denote the meet of $\{x_i\}_{i \in I}$ by $\bigwedge_{i \in I} x_i$. For instance, in the posetal category $(\R, \leq)$, the meet of a family $\{x_i\}_{i \in I}$ of real numbers is given by the infimum $\inf\{x_i\}_{i \in I}$, and in $(\~P(S), \subseteq)$, the meet of a family of sets $\{X_i\}_{i \in I}$ is just the intersection $\bigcap_{i \in I} X_i$.

Conversely, an \emph{upper bound} for a family of elements $\{x_i\}_{i \in I}$ is an element $b \in P$ such that $x_i \leq b$ for all $i$. A \emph{least upper bound} of \emph{join} of the family $\{x_i\}_{i \in I}$ is an upper bound for the family that is less than any other. We denote the join of $\{x_i\}_{i \in I}$ by $\bigvee_{i \in I} x_i$. So in $(\R, \leq)$, the join of $\{x_i\}_{i \in I}$ is given by the supremum $\sup\{x_i\}_{i \in I}$, and in $(\~P(S), \subseteq)$, the join of a family of sets $\{X_i\}_{i \in I}$ is just the union $\bigcup_{i \in I} X_i$ of the entire family.
\end{ex}

\begin{ex}
 Given an $I$-indexed family $\{M_i\}_{i \in I}$ of $(R,S)$-bimodules, recall that the product of this family is defined by
    \[\prod_{i\in I} M_i := \{(m_i)_{i \in I} : m_i \in M_i \text{ for all } i \in I\},\]
    with the obvious $(R,S)$-bimodule structure (addition and scalar multiplication defined component-wise). To show this object satisfies the universal property of the product in $_R\Mod_S$, we first define projections $\pi_i: \prod_{i \in I}M_i \to M_i$ by $\pi_i(m_i)_{i \in I} = m_i$ for all $i \in I$. As each $M_i$ is an $(R,S)$-bimodule, it is clear that the maps $\pi_i$ are homomorphisms. Now suppose we are given $N \in {}_R\Mod_S$ and maps $\{f_i: N \to M_i\}_{i \in I}$. Defining $f: N \to \prod_{i \in I}M_i$ by $f(n) := (f_i(n))_{i \in I}$, we see that the diagram
    \[
        \begin{tikzcd}
            N \\
            {\prod_{i \in I}M_i} & {M_i}
            \arrow["f"', from=1-1, to=2-1]
            \arrow["{f_i}", from=1-1, to=2-2]
            \arrow["{\pi_i}"', from=2-1, to=2-2]
        \end{tikzcd}
    \]
    commutes for every $i \in I$. Moreover, $f$ is the \emph{unique} map with this property, for if some other map $f'$ makes the diagram above commute, then for each $n \in N$ and $i \in I$, we must have $\pi_i(f'(n)) = f_i(n)$, hence $f = f'$. From this we see that $\prod_{i \in I}M_i$ satisfies the universal property of the product in $_R\Mod_S$.

    We now turn our attention to the coproduct in $_R\Mod_S$. By definition,
    \[\bigoplus_{i \in I}M_i := \left\{(m_i)_{i\in I} \in \prod_{i \in I}M_i : m_i = 0 \text{ for all but finitely many } i \in I \right\}.\]
    To show that $\bigoplus_{i \in I}M_i$ satisfies the universal property of the coproduct in $_R\Mod_S$, we first define coprojections $\iota_i: M_i \to \bigoplus_{i \in I}M_i$ by 
    \[\iota_i (x) := (m_i)_{i \in I} \text{ where } m_i := \begin{cases}x & i = j \\ 0 & i \ne j \end{cases}\, \text{ for all } x \in M_i.\]
    Now suppose we have $N \in {}_R\Mod_S$ and an $I$-indexed family of maps $\{g_i: M_i \to N\}_{i \in I}$. Given $m = (m_i)_{i \in I} \in \bigoplus_{i \in I} M_i$, define $g: \bigoplus_{i \in I} M_i \to N$ by $g(m) = \sum_{i \in I} g_i(m_i)$ (note this sum is well-defined since all but a finite number of the $m_i$ are zero). Then for each $x \in M_i$, we have $g(\iota_i(x)) = g_i(x) + \sum_{j \ne i} g_j(0) = g_i(x)$, so the diagram
    \[\begin{tikzcd}
    {M_i} & {\bigoplus_{i \in I}M_i} \\
    & N
    \arrow["{\iota_i}", from=1-1, to=1-2]
    \arrow["{g_i}"', from=1-1, to=2-2]
    \arrow["g", from=1-2, to=2-2]\end{tikzcd}\]
    commutes for all $i \in I$. If some other map $g'$ makes this diagram commute, then we'd have $g'(\iota_i(x)) = g_i(x)$ for all $x \in M_i$, so $g = g'$. Therefore $\bigoplus_{i \in I}M_i$ satisfies the universal property of the coproduct in ${}_R \Mod_S$.
\end{ex}


\section{Equalizers}

Before defining equalizers (our second kind of limit) we need to introduce a piece of terminology: a \emph{fork} in a category $\~C$ consists of objects and maps
\begin{equation}
\label{eq:fork}
    \begin{tikzcd}
        A & X & Y
        \arrow["f", from=1-1, to=1-2]
        \arrow["t"', shift right, from=1-2, to=1-3]
        \arrow["s", shift left, from=1-2, to=1-3]
    \end{tikzcd}
\end{equation}
such that $sf = tf$. 

\begin{defn}
    Let $\~C$ be a category and let $\begin{tikzcd}[column sep = small]
        X & Y
        \arrow["t"', shift right, from=1-1, to=1-2]
        \arrow["s", shift left, from=1-1, to=1-2]
    \end{tikzcd}$ be objects and maps in $\~C$. An \emph{equalizer} of $s$ and $t$ is an object $E$ together with a map $E \To{i} A$ such that 
\[\begin{tikzcd}
    E & X & Y
    \arrow["i", from=1-1, to=1-2]
    \arrow["t"', shift right, from=1-2, to=1-3]
    \arrow["s", shift left, from=1-2, to=1-3]
\end{tikzcd}\]
is a fork, with the property that for any fork of the form (\ref{eq:fork}) there exists a unique map $\bar{f}: A \to E$ such that the diagram
\[\begin{tikzcd}
    & E \\
    A & X
    \arrow["f"', from=2-1, to=2-2]
    \arrow["{\bar{f}}", dashed, from=2-1, to=1-2]
    \arrow["i", from=1-2, to=2-2]
\end{tikzcd}\]
commutes.
\end{defn}
Just like with products, an equalizer of $\begin{tikzcd}[column sep = small]
        X & Y
        \arrow["t"', shift right, from=1-1, to=1-2]
        \arrow["s", shift left, from=1-1, to=1-2]
    \end{tikzcd}$ may not always exist, but when it does, it is determined up to a unique isomorphism. 
\begin{ex}
In $\Set$, the equalizer of two functions $\begin{tikzcd}[column sep = small]
        X & Y
        \arrow["t"', shift right, from=1-1, to=1-2]
        \arrow["s", shift left, from=1-1, to=1-2]
    \end{tikzcd}$ is given by the set
\[E := \{x \in X : s(x) = t(x)\} \subseteq X\]
where we take $i: E \mono X$ to be the inclusion of $E$ into $X$. Then $si  = ti$, so this is a fork. Now, take a set $A$ and suppose $\begin{tikzcd}
            A & X & Y
            \arrow["f", from=1-1, to=1-2]
            \arrow["t"', shift right, from=1-2, to=1-3]
            \arrow["s", shift left, from=1-2, to=1-3]
        \end{tikzcd}$ is a fork. Then $f(a) \in X$ and $s(f(a)) = t(f(a))$, so $f(a) \in E$ for all $a \in A$. We may now define a map $\bar{f}: A \to E$ by $\bar{f}(a) = f(a)$ for all $a \in A$. This gives a commutative triangle 
        \[\begin{tikzcd}
            & E \\
            A & X
            \arrow["f"', from=2-1, to=2-2]
            \arrow["{\bar{f}}", dashed, from=2-1, to=1-2]
            \arrow["i", from=1-2, to=2-2, hook]
        \end{tikzcd}\]   
        and it is clear that the map $\bar{f}$ is the \emph{unique} map making this triangle commute. This shows that $E$ is indeed the equalizer of $\begin{tikzcd}[column sep = small]
        X & Y
        \arrow["t"', shift right, from=1-1, to=1-2]
        \arrow["s", shift left, from=1-1, to=1-2]
    \end{tikzcd}$.
\end{ex}


\begin{prop}[Equalizers are monics, coequalizers are epis]
\label{prop:equalizer_monic}
    Let $\~C$ be a category.
    \begin{enumerate}[label=(\roman*)]
        \item If \begin{equation}
            \label{eq:equalizer_monic}
                \begin{tikzcd}
                    E & X & Y
                    \arrow["i", from=1-1, to=1-2]
                    \arrow["t"', shift right, from=1-2, to=1-3]
                    \arrow["s", shift left, from=1-2, to=1-3]
                \end{tikzcd}
            \end{equation}
        is an equalizer diagram in $\~C$, then $i$ is monic.
        \item If \begin{equation}\begin{tikzcd}
            X & Y & C
            \arrow["s", shift left, from=1-1, to=1-2]
            \arrow["t"', shift right, from=1-1, to=1-2]
            \arrow["j", from=1-2, to=1-3]
        \end{tikzcd}\end{equation}
        is a coequalizer diagram (i.e., an equalizer in $\~C^\op$), then $j$ is epic.
\end{enumerate}
\end{prop}
\begin{proof}
    We will prove (i) and then conclude (ii) by duality. Suppose we have parallel arrows \begin{tikzcd}[column sep = small] A \arrow[r, shift left]{}{a} \arrow[r, shift right, swap]{}{a'} & E\end{tikzcd} such that $i a = ia'$. Now, we have 
    \begin{align*}
        sia &= sia' \\
        &= (si)a' \\
        &= (ti)a' & (\text{since } (\ref{eq:equalizer_monic}) \text{ is an equalizer}) \\
        &= t(ia')\\
        &= tia
    \end{align*}
    So the diagram
    \[\begin{tikzcd}
            A & X & Y
            \arrow["ia", from=1-1, to=1-2]
            \arrow["t"', shift right, from=1-2, to=1-3]
            \arrow["s", shift left, from=1-2, to=1-3]
        \end{tikzcd}\]
    is a fork, and since (\ref{eq:equalizer_monic}) is an equalizer diagram, there is a unique map $u$ making

    \[\begin{tikzcd}
        & E \\
        A & X
        \arrow["i", from=1-2, to=2-2]
        \arrow["{ u}", dashed, from=2-1, to=1-2]
        \arrow["ia"', from=2-1, to=2-2]
    \end{tikzcd}\]
    Then we clearly have $u = a$. But we also know that $ia = ia'$, so $u = a' = a$, and hence $i$ is monic. Reversing all the arrows, we get the result of (ii), and the proof is complete.
\end{proof}

Because of this proposition, it is conventional to draw equalizers like 
\[\begin{tikzcd}
            E & X & Y
            \arrow["i", hook, from=1-1, to=1-2]
            \arrow["t"', shift right, from=1-2, to=1-3]
            \arrow["s", shift left, from=1-2, to=1-3]
        \end{tikzcd}\]

\begin{ex}
As we have seen, an equalizer can be used to describe the set of solutions to an equation. When combined with products, we can use them to describe the set of solutions to a whole family of equations. Take a set $I$ and a family of maps
\[\left( \begin{tikzcd} X \arrow[r, shift left]{}{s_i} \arrow[r, shift right, swap]{}{t_i} & Y_i \end{tikzcd} \right)_{i \in I}\] 
in $\Set$. So we have a pair of maps $\begin{tikzcd}[column sep = small] X \arrow[r, shift left]{}{s_i} \arrow[r, shift right, swap]{}{t_i} & Y_i\end{tikzcd}$, one for each $i \in I$. The solution set 
\[\{x \in X : s_i(x) = t_i(x) \text{ for all } i \in I \}\]
is the equalizer of the product maps 
\[\begin{tikzcd} X \arrow[r, shift left]{}{(s_i)_{i \in I}} \arrow[r, shift right, swap]{}{(t_i)_{i \in I}} & \prod_i Y_i\end{tikzcd}\]
\end{ex}

\begin{ex}[Kernels are equalizers]
    The next example is one of the most important cases of equalizers. Let $\theta: G \to H$ be a group homomorphism. Associated to $\theta$ is a diagram

    We will return to the correspondence between equalizers and kernels in Chapter \ref{chap:abelian_cats}, when we discuss abelian categories.
\end{ex}

\section{Pullbacks}

We will now introduce pullbacks, which are similar to products, but with a little more information involved.

\begin{defn}
    Let $\~C$ be a category, and suppose we have objects and maps of $\~C$ in the configuration
    \begin{equation}
    \label{eq:setup_pullback}
        \begin{tikzcd}
                & Y \\
                X & Z
                \arrow["s"', from=2-1, to=2-2]
                \arrow["t", from=1-2, to=2-2]
        \end{tikzcd}
    \end{equation}
    A \emph{pullback} of this diagram is an object $P \in \~C$, together with maps $\pi_X: P \to X$ and $\pi_Y: P \to Y$ (as usual, these are called \emph{projections}) such that
    \begin{equation}
    \label{eq:pullbacksquare}
        \begin{tikzcd}
                P & Y \\
                X & Z
                \arrow["t", from=1-2, to=2-2]
                \arrow["{\pi_X}"', from=1-1, to=2-1]
                \arrow["s"', from=2-1, to=2-2]
                \arrow["{\pi_Y}", from=1-1, to=1-2]
        \end{tikzcd}
    \end{equation} 
    commutes, and with the property that for any commutative square
    \[
        \begin{tikzcd}
                A & Y \\
                X & Z
                \arrow["{g}", from=1-1, to=1-2]
                \arrow["s"', from=2-1, to=2-2]
                \arrow["t", from=1-2, to=2-2]
                \arrow["{f}"', from=1-1, to=2-1]
        \end{tikzcd}
    \]
    in $\~C$, there is a unique map $\langle f, g \rangle: A \to P$ such that 

\begin{equation}
\label{eq:pullback}
    \begin{tikzcd}
            A \\
            & P & Y \\
            & X & Z
            \arrow["{\langle f, g \rangle}"{description}, dashed, from=1-1, to=2-2]
            \arrow["{\pi_X}", from=2-2, to=3-2, swap]
            \arrow["{\pi_Y}", from=2-2, to=2-3]
            \arrow["s"', from=3-2, to=3-3]
            \arrow["t", from=2-3, to=3-3]
            \arrow["{f}"', bend right, from=1-1, to=3-2]
            \arrow["{g}", bend left, from=1-1, to=2-3]
    \end{tikzcd}
\end{equation}
    commutes. (Note that since the commutativity of the square is already given, so commutativity of (\ref{eq:pullback}) means only that $\pi_X \circ \langle f, g \rangle = f$ and $\pi_Y \circ \langle f, g \rangle = g$). 
\end{defn}

The square (\ref{eq:pullbacksquare}) is called a pullback square. Another name for a pullback is a \emph{fibred} product. The following example explains this alternative name.

\begin{ex}
    Indeed, if $1$ is a terminal object in a category $\~C$, a product $P$ of $X$ and $Y$ is precisely a pullback of the diagram
    \begin{equation}
    \label{ex:terminal_setup}
        \begin{tikzcd}
            & Y \\
            X & 1
            \arrow["s"', from=2-1, to=2-2]
            \arrow["t", from=1-2, to=2-2]
        \end{tikzcd}
    \end{equation}
    To see this, take objects and maps 
    \[\begin{tikzcd}[column sep=small]
            & A \arrow[dl, "f", swap] \arrow[dr, "g"] & \\
            X & & Y
    \end{tikzcd}\]
    and note that when $P$ is a product, there are projections $\pi_X, \pi_Y$ and a unique map $\langle f, g \rangle$ making 
    \begin{equation}
    \label{eq:example_terminal_products}
        \begin{tikzcd}
        A \arrow[rd, "{\langle f, g \rangle}" description, dashed] \arrow[rdd, "f", bend right, swap] \arrow[rrd, "g", bend left] &                                      &   \\
                                                                                                                  & P \arrow[r, "\pi_Y"] \arrow[d, "\pi_X"'] & Y \\
                                                                                                                  & X                                    &  
        \end{tikzcd}
    \end{equation}
    commute. Since $1$ is terminal, we easily check that the diagram
    \[
        \begin{tikzcd}
            A \\
            & P & Y \\
            & X & 1
            \arrow["{\langle f, g \rangle}"{description}, dashed, from=1-1, to=2-2]
            \arrow["{\pi_X}"', from=2-2, to=3-2]
            \arrow["{\pi_Y}", from=2-2, to=2-3]
            \arrow["s"', from=3-2, to=3-3]
            \arrow["t", from=2-3, to=3-3]
            \arrow["{f}"', bend right, from=1-1, to=3-2]
            \arrow["{g}", bend left, from=1-1, to=2-3]
        \end{tikzcd}
    \]
    commutes. Importantly, however, we know that the diagram (\ref{eq:example_terminal_products}) commutes, which can only happen since $P$ is a product of $X$ and $Y$. The rest of the diagram above will always commute because $1$ is terminal, but (\ref{eq:example_terminal_products}) will not commute unless $P$ is a product. Therefore, any pullback of (\ref{ex:terminal_setup}) is a product of $X$ and $Y$.
\end{ex}

\begin{ex}
    The pullback of the diagram (\ref{eq:setup_pullback}) in $\Set$ is 
    \[P = \{(x, y) \in X \times Y : s(x) = t(y)\}\]
    To see this, suppose we are given the diagram (\ref{eq:setup_pullback}) for sets $X, Y, Z$ and functions $s, t$. Let the projection $\pi_X: P \to X$ be given by $\pi_X(x, y) = x$ for all $(x, y) \in P$. We define our second projection $\pi_Y: P \to Y$ similarly by letting $\pi_Y(x, y) = y$ for all $(x, y) \in P$. With $\pi_X$ and $\pi_Y$ defined in this way, it follows that the square
    \[
        \begin{tikzcd}
           P & Y \\
           X & Z
           \arrow["t", from=1-2, to=2-2]
           \arrow["{\pi_X}"', from=1-1, to=2-1]
           \arrow["s"', from=2-1, to=2-2]
           \arrow["{\pi_Y}", from=1-1, to=1-2]
        \end{tikzcd}
    \]
    commutes.
    Now, take another commutative square
    \[
        \begin{tikzcd}
        A & Y \\
        X & Z
        \arrow["{g}", from=1-1, to=1-2]
        \arrow["s"', from=2-1, to=2-2]
        \arrow["t", from=1-2, to=2-2]
        \arrow["{f}"', from=1-1, to=2-1]
        \end{tikzcd}
    \]
    in $\Set$. Define a map $\langle f, g \rangle : A \to P$ by 
    \[\langle f, g \rangle (a) = (f(a), g(a))\]
    With $\langle f, g \rangle $ defined in this way, it is easy to see that the diagram
    \[
        \begin{tikzcd}
            A \\
            & P & Y \\
            & X & Z
            \arrow["{\langle f, g \rangle}"{description}, dashed, from=1-1, to=2-2]
            \arrow["{\pi_X}"', from=2-2, to=3-2]
            \arrow["{\pi_Y}", from=2-2, to=2-3]
            \arrow["s"', from=3-2, to=3-3]
            \arrow["t", from=2-3, to=3-3]
            \arrow["{f}"', bend right, from=1-1, to=3-2]
            \arrow["{g}", bend left, from=1-1, to=2-3]
        \end{tikzcd}
    \]
    commutes. We now need to show that $\langle f, g \rangle $ is a unique map. To that end, suppose we had some other arrow $\overline{\langle f, g \rangle}: A \to P$ making the above diagram commute. Then for any $a \in A$, $\overline{\langle f, g \rangle}(a) = (x, y) \in P$, so we have 
   \[ \pi_X \overline{\langle f, g \rangle}(a) = x \quad \text{and} \quad  \pi_Y \overline{\langle f, g \rangle}(a) = y\] 
    But since $\overline{\langle f, g \rangle}$ makes the diagram above commute, we deduce that $x = f(a)$ and $y = g(a)$, and hence
    \[\overline{\langle f, g \rangle} (a) = \left( f(a), g(a) \right) = \langle f, g \rangle\]
    so the map $\langle f, g \rangle$ is unique, and therefore $P$ satisfies the definition of a pullback in $\Set$.
\end{ex}

\begin{prop}
    If $s$ is a monomorphism and the square
        \[\begin{tikzcd}
            P & Y \\
            X & Z
            \arrow["{\pi_Y}", from=1-1, to=1-2]
            \arrow["s"', from=2-1, to=2-2, hook]
            \arrow["{\pi_X}"', from=1-1, to=2-1]
            \arrow["t", from=1-2, to=2-2]
        \end{tikzcd}\]
    is a pullback, then $\pi_Y$ is also a monomorphism.
\end{prop}

\begin{proof}
Let $A \in \~C$ and suppose we have maps \begin{tikzcd}[column sep = small]A \arrow[r, shift left]{}{a'} \arrow[r, shift right, swap]{}{a} & P\end{tikzcd} such that $\pi_Y a = \pi_Y a' $. Applying $t$ to both sides, we get
\[t\pi_Y a = t\pi_Y a' \]
Now, since the square
        \[\begin{tikzcd}
            P & Y \\
            X & Z
            \arrow["{\pi_Y}", from=1-1, to=1-2]
            \arrow["s"', from=2-1, to=2-2, hook]
            \arrow["{\pi_X}"', from=1-1, to=2-1]
            \arrow["t", from=1-2, to=2-2]
        \end{tikzcd}\]
commutes, we have $t\pi_Y = s\pi_X$, so we can write $s\pi_X a = s\pi_X a'$, and $s$ being monic, we deduce that
\[\pi_X a= \pi_X a'\]
Now, this means that the diagram
    \[
        \begin{tikzcd}
            A \\
            & P & Y \\
            & X & Z
            \arrow["{a}"{description}, dashed, from=1-1, to=2-2]
            \arrow["{\pi_X}"', from=2-2, to=3-2]
            \arrow["{\pi_Y}", from=2-2, to=2-3]
            \arrow["s"', from=3-2, to=3-3, hook]
            \arrow["t", from=2-3, to=3-3]
            \arrow["{\pi_X a'}"', bend right, from=1-1, to=3-2]
            \arrow["{\pi_Y a'} ", bend left, from=1-1, to=2-3]
        \end{tikzcd}
    \]
    commutes. This means that both $a$ and $a'$ should also fill in the dotted arrow $A \to P$. However, by universality of the pullback, this arrow is unique, so we conclude that $a = a'$, and therefore $\pi_Y$ is monic. 
\end{proof}

The fact that monomorphisms satisfy the above property is often phrased as ``monics are \emph{stable under pullbacks}''.

\begin{lem}[Pullback lemma]

 In a category $\~C$, consider a diagram
\[\begin{tikzcd}
    A & B & C \\
    {X} & {Y} & {Z}
    \arrow["g", from=1-2, to=2-2]
    \arrow["v", from=1-2, to=1-3]
    \arrow["{h}", from=1-3, to=2-3]
    \arrow["{q}"', from=2-2, to=2-3]
    \arrow["u", from=1-1, to=1-2]
    \arrow["{p}"', from=2-1, to=2-2]
    \arrow["f", from=1-1, to=2-1]
\end{tikzcd}\]
Assume that the right-hand square is a pullback. Then the left-hand square is a pullback if and only if the outer rectangle is a pullback.
\end{lem}

\begin{proof}
To prove ``if'', suppose that the two inner squares are pullbacks. To show that the outer diagram is a pullback, suppose we have a commuting diagram 
    \[\begin{tikzcd}
    P \\
    &&& C \\
    & {X} & {Y} & {Z}
    \arrow["{h}", from=2-4, to=3-4]
    \arrow["{q}"', from=3-3, to=3-4]
    \arrow["{p}"', from=3-2, to=3-3]
    \arrow["{r'}"', bend right, from=1-1, to=3-2]
    \arrow["r", bend left, from=1-1, to=2-4]
    \end{tikzcd}\]
Because the right-hand square is a pullback and the rectangle above commutes, we can find a unique map $w: P \to B$ making the following diagram commute:
\begin{equation}
    \label{eq:pullback_lemma_rightsq}
    \begin{tikzcd}
        P \\
        && B & C \\
        & {X} & {Y} & {Z}
        \arrow["{h}", from=2-4, to=3-4]
        \arrow["{q}"', from=3-3, to=3-4]
        \arrow["{p}"', from=3-2, to=3-3]
        \arrow["{r'}"', bend right, from=1-1, to=3-2]
        \arrow["r", bend left, from=1-1, to=2-4]
        \arrow["g"', from=2-3, to=3-3]
        \arrow["v", from=2-3, to=2-4]
        \arrow["w"{description}, dashed, from=1-1, to=2-3]
    \end{tikzcd}
\end{equation}
But now we have a commutative diagram 
\[\begin{tikzcd}
    P \\
    && B \\
    & {X} & {Y}
    \arrow["{p}"', from=3-2, to=3-3]
    \arrow["g", from=2-3, to=3-3]
    \arrow["{r'}"', bend right, from=1-1, to=3-2]
    \arrow["{w}", bend left, dashed, from=1-1, to=2-3]
\end{tikzcd}\]
and since the left-hand square is a pullback, we must have another unique map $y: P \to A$ such that $uy = w$ and $fy = r'$:
\begin{equation}
    \label{eq:pullbacklem_final}
    \begin{tikzcd}
        P \\
        & A & B & C \\
        & {X} & {Y} & {Z}
        \arrow["f"', from=2-2, to=3-2]
        \arrow["g"', from=2-3, to=3-3]
        \arrow["{h}", from=2-4, to=3-4]
        \arrow["u", from=2-2, to=2-3]
        \arrow["v", from=2-3, to=2-4]
        \arrow["{p}"', from=3-2, to=3-3]
        \arrow["{q}"', from=3-3, to=3-4]
        \arrow["{y}"{description}, dashed, from=1-1, to=2-2]
        \arrow["{r'}"', curve={height=18pt}, from=1-1, to=3-2]
        \arrow["w"{description}, curve={height=-6pt}, dashed, from=1-1, to=2-3]
        \arrow["r", curve={height=-18pt}, from=1-1, to=2-4]
    \end{tikzcd}
\end{equation}
We claim $y$ is the morphism we're looking for. We already have $fy = r'$ and $vuy = vw$. Diagram (\ref{eq:pullback_lemma_rightsq}) tells us $vw = r$, so
$vuy = vw = r$ as required. To show that $y$ is unique, suppose there is another map $y': P \to A$ such that $(vu)y' = r$ and $fy' = r'$. Let $w' = uy'$, then $vw' = r$, and $gw' = guy' = pfy' = pr'$, so by the uniqueness property of the right-hand square, we have $w = w'$. But now $w = uy'$ and $fy' = r'$, so by the uniqueness property of the left-hand square we obtain $y = y'$ as desired.

To prove ``only if'', suppose the outer rectangle is a pullback, and consider a commutative diagram
\[\begin{tikzcd}
    P \\
    && Y \\
    & A & B
    \arrow["r", curve={height=-6pt}, from=1-1, to=2-3]
    \arrow["{r'}"', curve={height=6pt}, from=1-1, to=3-2]
    \arrow["g", from=2-3, to=3-3]
    \arrow["p"', from=3-2, to=3-3]
\end{tikzcd}\]
We can expand this to the following commuting diagram:
\[\begin{tikzcd}
    P \\
    && Y & Z \\
    & A & B & C
    \arrow["r", from=1-1, to=2-3]
    \arrow["{v r}", curve={height=-12pt}, from=1-1, to=2-4]
    \arrow["{r'}"', curve={height=6pt}, from=1-1, to=3-2]
    \arrow["v", from=2-3, to=2-4]
    \arrow["g", from=2-3, to=3-3]
    \arrow["h", from=2-4, to=3-4]
    \arrow["p"', from=3-2, to=3-3]
    \arrow["q"', from=3-3, to=3-4]
\end{tikzcd}\]
Now the outer rectangle is a pullback, so we can find a unique arrow $w: P \to X$ fitting into the following commutative diagram:
\[\begin{tikzcd}
    P \\
    & X & Y & Z \\
    & A & B & C
    \arrow["w", dashed, from=1-1, to=2-2]
    \arrow["{v r}", curve={height=-12pt}, from=1-1, to=2-4]
    \arrow["{r'}"', curve={height=6pt}, from=1-1, to=3-2]
    \arrow["u", from=2-2, to=2-3]
    \arrow["f", from=2-2, to=3-2]
    \arrow["v", from=2-3, to=2-4]
    \arrow["g", from=2-3, to=3-3]
    \arrow["h", from=2-4, to=3-4]
    \arrow["p"', from=3-2, to=3-3]
    \arrow["q"', from=3-3, to=3-4]
\end{tikzcd}\]
We already know that $w$ is unique and $f w = r'$, so to show that the left-hand square is a pullback, we just need to check that $u w = r$, so that the diagram
\[\begin{tikzcd}
    P \\
    & X & Y \\
    & A & B
    \arrow["w"{description}, dashed, from=1-1, to=2-2]
    \arrow["r", curve={height=-6pt}, from=1-1, to=2-3]
    \arrow["{r'}"', curve={height=6pt}, from=1-1, to=3-2]
    \arrow["u", from=2-2, to=2-3]
    \arrow["f"', from=2-2, to=3-2]
    \arrow["g", from=2-3, to=3-3]
    \arrow["p"', from=3-2, to=3-3]
\end{tikzcd}\]
commutes. Indeed, observe that $uw$ and $r$ both fill in the dotted arrow in the following diagram:
\[\begin{tikzcd}
    P \\
    & Y & Z \\
    & B & C
    \arrow[dashed, from=1-1, to=2-2]
    \arrow["{v r}", curve={height=-6pt}, from=1-1, to=2-3]
    \arrow["{p r'}"', curve={height=6pt}, from=1-1, to=3-2]
    \arrow["v", from=2-2, to=2-3]
    \arrow["g", from=2-2, to=3-2]
    \arrow["h", from=2-3, to=3-3]
    \arrow["q"', from=3-2, to=3-3]
\end{tikzcd}\]
Since the right-hand square is a pullback, there can only be one arrow that completes this diagram, so $uw = r$ and we see that the left-hand square must also be a pullback. This proves the claim. 
\end{proof}

\begin{prop}
\label{prop:pushout_tensorprod}
    The tensor product $A \otimes_R B$ is a pullback in $\cAlg_R^\op$, that is, a pushout in $\cAlg_R$.
\end{prop}
\section{General limits and colimits}

So far we've examined three constructions and their duals: products/coproducts, equalizers/coequalizers, and pullbacks/pushouts. Let's think about what these constructions have in common. First, let's note that they all begin with some kind of diagram, i.e., some objects in a category $\~C$ and some arrows in $\~C$ between them. In the definition of a product, our starting diagram is just the two objects $X$ and $Y$:
\begin{equation}
\label{diagramshape_product}
\begin{tikzcd}
    X & Y
\end{tikzcd}
\end{equation}
for equalizers, we start with two objects and two arrows:
\begin{equation}
\label{diagramshape_equa}
\begin{tikzcd}
    X & Y
    \arrow["s", shift left, from=1-1, to=1-2]
    \arrow["t"', shift right, from=1-1, to=1-2]
\end{tikzcd}
\end{equation}
and for pullbacks, we start with the familiar diagram
\begin{equation}
\label{diagramshape_pullback}
\begin{tikzcd}
    & Y \\
    X & Z
    \arrow["s"', from=2-1, to=2-2]
    \arrow["t", from=1-2, to=2-2]
\end{tikzcd}
\end{equation}
We would like to formulate a definition for the limit of any given diagram. To do this, we need to formalize our notion of what a ``diagram in a category'' is. We have already seen that a functor from the one-object category $\mathbf{1} = \boxed{\ast}$ to a category $\~C$ is simply a choice of an object $X \in \~C$, and similarly, a choice of map $X \to Y$ in $\~C$ is simply a functor $\mathbf{I} \to \~C$ where $\mathbf{I} = \boxed{\cdot \to \cdot}$ is the interval category. In both these situations, we call the small category in the codomain our \emph{index category}. It is easy to come up with more complicated examples of index categories. For instance, consider the small categories
\begin{equation} \mathbf{T} = 
\begin{boxedtikzcd}
\cdot & \cdot
\end{boxedtikzcd}~, \quad \mathbf{E} = \begin{boxedtikzcd}
    \cdot & \cdot
    \arrow[shift left, from=1-1, to=1-2]
    \arrow[shift right, from=1-1, to=1-2]
\end{boxedtikzcd}~, \quad \text{and} \quad \mathbf{P} = \begin{boxedtikzcd}    & \cdot \\
    \cdot & \cdot
    \arrow[from=2-1, to=2-2]
    \arrow[from=1-2, to=2-2]
\end{boxedtikzcd}
\end{equation}
then functors $\mathbf{T} \to \~C$, $\mathbf{E} \to \~C$, and $\mathbf{P} \to \~C$ give the data of (\ref{diagramshape_product}), (\ref{diagramshape_equa}), and (\ref{diagramshape_pullback}) respectively. 
\begin{defn}
Let $\~C$ be a category and let $\mathbf{J}$ be a small category. A functor $\mathbf{J} \to \~C$ is called a \emph{diagram of shape $\mathbf{J}$} in $\~C$. 
\end{defn} 
So (\ref{diagramshape_product}), (\ref{diagramshape_equa}), and (\ref{diagramshape_pullback}) are diagrams of shape $\mathbf{T}$, $\mathbf{E}$, and $\mathbf{P}$. 
\begin{defn}
Let $\~C$ be a category, $\mathbf{J}$ a small category, and $D: \mathbf{J} \to \~C$ a diagram in $\~C$.
\begin{enumerate}
\item A \emph{cone} on the diagram $D$ is an object $A \in \~C$ (the \emph{vertex} of the cone) together with a family
\begin{equation}
\label{cone}
\left(A \To{f_J} D(J) \right)_{J \in \mathbf{J}}
\end{equation}
of maps in $\~C$ such that for each map $J \To{u} K$ in  $\mathbf{J}$, the triangle
\[\begin{tikzcd}
    & A \\
    {D(J)} && {D(K)}
    \arrow["{f_J}"', from=1-2, to=2-1]
    \arrow["{f_K}", from=1-2, to=2-3]
    \arrow["Du"', from=2-1, to=2-3]
\end{tikzcd}\]
commutes. 
\item Let $D: \mathbf{J} \to \~C$ be a diagram in $\~C$, and write $D^\op$ for the corresponding functor $\mathbf{J}^\op \to \~C^\op$. A \emph{cocone} on $D$ is, naturally, a cone on $D^\op$. Explicitly, this means that a cocone is an object $A$ (the \emph{vertex} of the cocone) together with a family
\begin{equation}
\label{cocone}
\left(D(J) \To{f_J} A \right)_{J \in \mathbf{J}}
\end{equation}
of maps in $\~C$ such that for all maps $J \To{u} K$ in $\mathbf{J}$, the diagram
\[\begin{tikzcd}
    & A \\
    {D(J)} && {D(K)}
    \arrow["{f_J}", from=2-1, to=1-2]
    \arrow["{f_K}"', from=2-3, to=1-2]
    \arrow["Du"', from=2-1, to=2-3]
\end{tikzcd}\]
commutes.
\end{enumerate}
\end{defn}
As one might have guessed by the similarity between the definition the definition of a cone and the definition of the generalized product (\ref{defn:gen_prod}), cones and cocones will play an important role in defining limits and colimits.
\begin{defn}
Let $\~C$ be a category, $\mathbf{J}$ a small category, and $D: \mathbf{J} \to \~C$ a diagram in $\~C$.
\begin{enumerate}[label=(\roman*)]
\item A \emph{limit} of the diagram $D$ is a cone 
\[
\left(L \To{\pi_J} D(J)\right)_{J \in \mathbf{J}}
\]
(sometimes called the \emph{limit cone} for emphasis) with the property that for any other cone (\ref{cone}) on $D$, there exists a unique map $\hat{f_J}: A \to L$ such that $\hat{f}_J \pi_J = f_J$ for all $J \in \mathbf{J}$. In this situation we say that the cone (\ref{cone}) \emph{factors through} the cone $\left(L \To{\pi_J} D(J)\right)_{J \in \mathbf{J}}$. The maps $\pi_J$ are called the \emph{projections} of the limit. 

\item The \emph{colimit} of $D$ is the limit of $D^\op$ (where $D^\op$ again denotes the functor $\mathbf{J}^\op \to \~C^\op$). Explicitly, this means a colimit of $D$ is a cocone
\[\left(D(J) \To{\pi_J} C \right)_{J \in \mathbf{J}}\]
(the ``colimit cocone'') with the property that for any cocone (\ref{cocone}) on $D$, there exists a unique map $\hat{f_J}: C \to A$ such that $\hat{f_J} \pi_J = f_J$ for all $J \in \mathbf{J}$. We call the maps $\pi_J$ \emph{coprojections}.
\end{enumerate}
We write $\lim_{\mathbf{J}} D$ (or just $\lim D$, if $\mathbf{J}$ is unambiguous) for the vertex of the cone $\left(L \To{\pi_J} D(J)\right)_{J \in \mathbf{J}}$ and $C = \colim_{\mathbf{J}} D$ (or just $\colim D$) for the vertex of the cocone $\left(D(J) \To{\pi_J} C \right)_{J \in \mathbf{J}}$. By a slight abuse of notation, we will sometimes refer to $L$ (resp. $C$) as the limit (resp. colimit) instead of the entire cone (resp. cocone).
\end{defn}

The definition of a limit tells us that for each pair of objects $J, K \in \mathbf{J}$ and every map $J \To{u} K$, we have a commutative diagram
\[\begin{tikzcd}
    & A \\
    \\
    & {\lim_{\mathbf{J}} D} \\
    {D(J)} && {D(K)}
    \arrow["{\langle f_J, f_K  \rangle}"{description}, dashed, from=1-2, to=3-2]
    \arrow["{f_J}"', curve={height=12pt}, from=1-2, to=4-1]
    \arrow["{f_K}", curve={height=-12pt}, from=1-2, to=4-3]
    \arrow["{\pi_J}", from=3-2, to=4-1]
    \arrow["{\pi_K}"', from=3-2, to=4-3]
    \arrow["Du"', from=4-1, to=4-3]
\end{tikzcd}\]
Similarly, the definition of a colimit tells us that for all such $J, K$, and $J \To{u} K$, we have a commutative diagram
\[\begin{tikzcd}
    {D(J)} && {D(K)} \\
    & {\colim_{\mathbf{J}} D} \\
    \\
    & A
    \arrow["Du", from=1-1, to=1-3]
    \arrow["{\pi_J}", from=1-1, to=2-2]
    \arrow["{f_J}"', curve={height=12pt}, from=1-1, to=4-2]
    \arrow["{\pi_K}"', from=1-3, to=2-2]
    \arrow["{f_K}", curve={height=-12pt}, from=1-3, to=4-2]
    \arrow["{\langle f_J, f_K  \rangle}"{description}, dashed, from=2-2, to=4-2]
\end{tikzcd}\]

\begin{ex}
Let $X_0$ be a set, and let $\mathbf{J} = \left( \~P(X_0), \, \subseteq \right)$. We can visualize $\mathbf{J}$ as follows:
\[
\cdots \subseteq X_2 \subseteq X_1 \subseteq X_0
\]
A diagram $D: \mathbf{J} \to \~C$ consists of a sequence
\[\cdots \To{s_3} X_2 \To{s_2} X_1 \To{s_1} X_0\] 
of maps in $\~C$. When $\~C$ is the category of sets $\Set$, the most straightforward diagram we can construct is the diagram $D: \mathbf{J} \to \Set$ which simply assigns each set $X_j$ to itself, and each inclusion of sets $X_{j + 1} \subseteq X_{j}$ to the inclusion map $s_{j + 1}: X_{j + 1} \mono X_{j}$. We claim that the limit of this diagram is $L = \bigcap_{j \in \N} X_j$. To verify this claim, we need to produce an appropriate limit cone and show it is universal among all other cones. A cone on $D$ with vertex $L$ consists of a objects and maps 
\[\left(L \To{\pi_j} X_j\right)_{j \in \N}\]
Now, since $L\subseteq X_j$ for all $j$ and every $X_j$ is contained in $X_0$, we can simply take each projection $\pi_j$ to be the same inclusion map $\pi: L \mono X_0$. Now take another cone
\[\left(A \To{f_j} X_j \right)_{j \in \N}\] 
on $D$. Since each $X_j$ is contained in $X_0$, we can again take each map $f_j: A \to X_j$ in this cone to be the same map $f: A \to X_0$. Verifying that we have a limit now comes down to showing there is a unique map $\langle f_j \rangle_{j \in \N}$ making the triangle 
\[\begin{tikzcd}
    A \\
    {\bigcap_{j \in \N}X_0} & {X_0}
    \arrow["{\langle f_j \rangle_{j \in \N}}"', dashed, from=1-1, to=2-1]
    \arrow["\pi"', hook, from=2-1, to=2-2]
    \arrow["f", from=1-1, to=2-2]
\end{tikzcd}\]
commute for each $j$, and it is easy to check that $f: A \to L \subseteq X_0$ is one map with this property. If $g$ is another such map, then $(\pi g) (a) = g(a) = f(a)$ for all $a \in A$, so the map $f$ is unique. Thus, we see that $L$ satisfies the universal property of the limit, so 
\[\lim D = \bigcap_{j \in \N} X_j\]
as claimed. In this and similar contexts, limits are sometimes referred to as \emph{inverse limits} (especially in older texts).
\end{ex}

\begin{ex}
\label{ex:directlim}
The dual concept of the inverse limit is the \emph{direct limit}, which refers to the colimit of a diagram $D$ of the following form
\begin{equation}
D: X_0 \To{s_1} X_1 \To{s_2} X_2 \To{s_3} \cdots
\end{equation}
In this case, the indexing category $\mathbf{J}$ is a posetal category whose underlying poset has the form
\[X_0 \subseteq X_1 \subseteq X_2 \subseteq \cdots\]
When the $X_i$'s are sets, a similar argument to the previous example shows that the colimit of the diagram $D$ is given by    
\[\colim D = \bigcup_{j \in \N} X_j\] 
We will not use the terms ``inverse limit'' and ``direct limit'' in this book, but they are worth knowing about, as many other texts make use of this terminology.
\end{ex}

\begin{ex}
    Let $n \geq 1$ be an integer. In this example we explain how the abelian group $\Q$ can be thought of as the colimit $\colim \frac{1}{n} \Z$. Note that $\frac{1}{n}\Z$ is a subgroup of $\Q$ for any $n \geq 1$. Given two subgroups $\frac{1}{m}\Z$ and $\frac{1}{n}\Z$, there is an inclusion $\frac{1}{m}\Z \subseteq \frac{1}{n} \Z$ if and only if $m$ divides $n$. We now need to describe the diagram for our colimit. Our indexing category in this case is the the posetal category $\mathbf{J} = (\N, ~|~)$ of natural numbers ordered by \emph{divisibility}. If $m ~|~ n$ in $\mathbf{J}$, then $mx = n$ for some $x \in \N$ and we can define an inclusion $s_n: \frac{1}{m}\Z \mono \frac{1}{n}\Z$ of abelian groups by
    \begin{align*}
        s_n: \frac{1}{m} \Z &\mono \frac{1}{n} \Z \\
        \frac{a}{m} &\mapsto \frac{ax}{mx} = \frac{ax}{n}
    \end{align*}
    Therefore, we get a diagram $D: \mathbf{J} \to \Ab$, where $D(n) = \frac{1}{n}\Z$ for any $n \geq 1$ and there is a homomorphism $s_n: \frac{1}{m} \Z \mono \frac{1}{n} \Z$ whenever $m ~|~ n$ in $\mathbf{J}$. Note that unlike in the previous two examples, this diagram does not have a nice representation as a single chain of arrows in $\Ab$. However, we can still compute this colimit, even if it is difficult to conceptualize pictorially. To determine the colimit, we need to find a cocone on $D$ through which all other cocones factor uniquely. For any $n \geq 1$, let $\pi_n$ be the map that sends $\frac{a}{n} \in \frac{1}{n} \Z$ to itself, treated as an element of $\Q$. Then clearly $\left(\frac{1}{n}\Z \To{\pi_n} \Q \right)_{n \geq 1}$ is a cocone on $D$. Let $\left(\frac{1}{n}\Z \To{f_n} A \right)_{n \geq 1}$ be another cocone on $D$. We now need to show that whenever $m ~|~ n$ (i.e., $m \to n$ in $\mathbf{J}$) there is a unique map $\langle f_m, f_n \rangle$ rendering the following diagram commutative:
    \[\begin{tikzcd}
    & A \\
    \\
    & \Q \\
    {\frac{1}{m}\Z} && {\frac{1}{n}\Z}
    \arrow["s_n"', hook, from=4-1, to=4-3]
    \arrow["{f_m}", curve={height=-12pt}, from=4-1, to=1-2]
    \arrow["{f_n}"', curve={height=12pt}, from=4-3, to=1-2]
    \arrow["{\langle f_m, f_n \rangle}"{description}, dashed, from=3-2, to=1-2]
    \arrow["{\pi_m}"', from=4-1, to=3-2]
    \arrow["{\pi_n}", from=4-3, to=3-2]
\end{tikzcd}\]


    % We therefore get a chain of inclusions
    % \[\Z \subseteq \frac{1}{m} \Z \subseteq \frac{1}{n} \Z \subseteq \cdots\]
    % and since this is a poset, we can take it to be our indexing category $\mathbf{J}$. A diagram $D: \mathbf{J} \to \Ab$
    % has the form
    % \[\Z \mono \frac{1}{m} \Z \mono \frac{1}{n} \Z \mono \cdots\]
    % where each homomorphism $\frac{1}{m} \Z \mono \frac{1}{n}\Z$ is the obvious inclusion. Following Example \ref{ex:directlim} above, we deduce that 
    % \[\colim \frac{1}{n} \Z = \bigcup\]
\end{ex}

\begin{prop}
    Let $D: \mathbf{J} \to \~C$ be a diagram and suppose $\mathbf{J}$ has a terminal object $1_{\mathbf{J}}$. Then
    \[\colim_{\mathbf{J}}D = D(1_{\mathbf{J}}).\]
\end{prop}

\begin{proof}
    For each $j \in \mathbf{J}$, we have a unique map $j \to 1_{\mathbf{J}}$, so we can define a cocone on $D$ by
    \[\left(F(j) \To{\alpha_j} F(1_{\mathbf{J}})\right)_{j \in \mathbf{J}}.\]
    Now take another cocone $\left(F(j) \To{\beta_j} C\right)_{j \in \mathbf{J}}$ on $D$. Then the map $\beta_{1_{\mathbf{J}}}: F(1_{\mathbf{J}}) \to C$ makes the diagram
\[\begin{tikzcd}[column sep = scriptsize]
    {F(j)} && {F(k)} \\
    & {F(1_{\mathbf{J}})} \\
    \\
    & C
    \arrow[from=1-1, to=1-3]
    \arrow["{\alpha_j}"', from=1-1, to=2-2]
    \arrow["{\beta_j}"', curve={height=6pt}, from=1-1, to=4-2]
    \arrow["{\alpha_k}", from=1-3, to=2-2]
    \arrow["{\beta_k}", curve={height=-6pt}, from=1-3, to=4-2]
    \arrow["{\beta_{1_{\mathbf{J}}}}"{description}, dashed, from=2-2, to=4-2]
\end{tikzcd}\]
commute. If another map $\vp$ makes this diagram commute, 
\end{proof}

\section{Functors and limits}
\newpage 
\section{Limits and colimits of presheaves}

\begin{rmk}
    Using the Yoneda embedding $H_\bullet: \~C \mono \Fun(\~C^\op, \Set)$ to pass from a category to its presheaf category can be regarded as freely adjoining colimits to $\~C$. For this reason, the Yoneda embedding is sometimes called the \emph{free cocompletion} of $\~C$.
\end{rmk}

\section{Adjoints and limits}

\section{Filtered (co)limits}

In this section we will discuss a special kind of (co)limit called a \emph{filtered (co)limit} that is often encountered in various mathematical contexts. These are (co)limits of diagrams $D: \mathbf{J} \to \~C$, but with an additional hypothesis on the indexing category $\mathbf{J}$. This is explained in the following definition:
\begin{defn}
\label{defn:filtered_limits}
    A nonempty category $\~J$ is said to be \emph{filtered} when
    \begin{enumerate}[label=(\roman*)]
        \item for every pair of objects $J$ and $J'$ in $\~J$, there exists an object $K$ and arrows $J \to K$ and $J' \to K$ in $\~J$, as in the diagram
        \[\begin{tikzcd}
            J && {J'} \\
            & K
            \arrow["f"', dashed, from=1-1, to=2-2]
            \arrow["{f'}", dashed, from=1-3, to=2-2]
        \end{tikzcd}\]
        \item for every two parallel arrows $\begin{tikzcd}
            J & {J'}
            \arrow["u", shift left, from=1-1, to=1-2]
            \arrow["v"', shift right, from=1-1, to=1-2]
        \end{tikzcd}$ in $\~J$ there exists an object $K$ and an arrow $w: J' \to K$ such that $wu = wv$, as in the commutative diamond
        \[\begin{tikzcd}
                & {J'} \\
                I && K \\
                & J'
                \arrow["w", dashed, from=1-2, to=2-3]
                \arrow["u", from=2-1, to=1-2]
                \arrow["v"', from=2-1, to=3-2]
                \arrow["w"', dashed, from=3-2, to=2-3]
            \end{tikzcd}\]
    \end{enumerate}
    Condition (i) states that the discrete 



    We say that $\~J$ is \emph{cofiltered} if $\~J^\op$ is filtered. If $\~C$ is any category and $\mathbf{J}$ is a small filtered category, then the limit of a diagram $D: \mathbf{J} \to \~C$ (called a \emph{filtered diagram}) is called a \emph{filtered} limit, and the colimit of $D$ is called a \emph{filtered colimit}. This should not be confused with \emph{cofiltered limits} (resp. colimits), which are limits (resp. colimits) of diagrams $D: \mathbf{J} \to \~C$ with $\mathbf{J}$ a small cofiltered category.
\end{defn}

One way to think about filtered categories is as categorical analogues of directed sets. Recall that a directed set is a poset $(X, \leq)$ in which every pair of elements has a least upper bound. In much the same way, a filtered category is a category where every finite diagram has a cocone. Indeed, if $(X, \leq)$ is a directed set, then the posetal category induced by $(X, \leq)$ is filtered. 



\begin{ex}
    index category of sets ordered by inclusion
\end{ex}

\begin{rmk}
We will mainly focus on filtered colimits in this text. In fact, if the category $\~C$ has a terminal object $1$, then any limit can be turned into a filtered limit in a trivial manner: Let $\mathbf{J}^+$ be the category obtained by adjoining a new terminal object, called $1$, to the category $\mathbf{J}$ (this means that if $\mathbf{J}$ already has a terminal object, we adjoin a second terminal object anyway). This means that $\ob({\mathbf{J}^+}) = \mathbf{J} \cup \{1\}$ and for each $J \in \mathbf{J}^+$ we adjoin a new arrow $J \to 1$ (the unique arrow in $\mathbf{J}^+(J, 1)$). Then $\mathbf{J}^+$ is a filtered category (check!) and so the limit of any diagram $\mathbf{J}^+ \to \~C$ is filtered. Therefore, in a category with a terminal object, any diagram $\mathbf{J} \to \~C$ can be trivially modified to create a filtered diagram $\mathbf{J}^+ \to \~C$. For this reason, filtered limits are not as special as filtered colimits, so we will typically focus on the latter.
\end{rmk}

\begin{ex}
One of the main reasons why filtered colimits are so useful is because the equivalence relation generated by a the colimit of a filtered diagram $D: \mathbf{J} \to \Set$ is stronger than the usual equivalence relation: If $D: \mathbf{J} \to \Set$ is a filtered diagram, then the colimit of $D$ is given by
\[\colim_{\mathbf{J}} D = \left(\bigsqcup_{J \in \mathbf{J}} D(J)  \right) \bigg/ \sim \]
Where $\sim$ is the equivalence relation defined by $(d_I, I) \sim (d_J, J)$ in $\colim_\mathbf{J}D$ if and only if there exists $K \in \mathbf{J}, \, \vp: I \to K,$ and $\vp': J \to K$ such that $D\vp(d_I) = D\vp'(d_J)$. The same principle is true for filtered diagrams in any set-like category, such as $\Ab$, $\cRing$, $\Mod_R$, etc... 
\end{ex}

\subsection*{Ind-objects}
We now turn our attention to a useful construction involving filtered colimits. In this section, we will write filtered diagrams in $\mathbf{I} \to \~C$ as $\{X_i\}_{\mathbf{I}}$ (this notation is due to Lurie).  Given a category $\~C$, there is a well-defined way of ``completing'' $\~C$ by adjoining a colimit for every filtered diagram in $\~C$. A standard example of a filtered diagram that one might encounter is an ascending chain of inclusions in $\~C$:
\[X_1 \mono X_2 \mono X_3 \mono \cdots\]
For a general category $\~C$, the colimit of such a diagram need not exist. However, we can add in the remaining colimits to get a new category called the \emph{ind-completion} (short for ``inductive completion'') of $\~C$.

Notice that given two filtered diagrams $\{X_i\}_{\mathbf{I}}$ and $\{Y_j\}_{\mathbf{J}}$, the assignment $(i,j) \mapsto \Hom_\~C(X_i, Y_j)$ defines a new filtered diagram $\mathbf{I}^\op \times \mathbf{J} \to \Set$\footnote{As a note on terminology, a $\Set$-valued bifunctor which is contravariant in the first argument and covariant in the second is sometimes called a \emph{profunctor}. The hom bifunctor in Definition \ref{defn:hom_bifun} is the archetypal example of a profunctor}.

\begin{defn}
    The \emph{ind-completion} (short for ``inductive completion'') of a category $\~C$ is the category $\Ind(\~C)$ whose objects are filtered diagrams in $\~C$ (called \emph{ind-objects}), and where
    \begin{align}\Hom_{\Ind(\~C)}\left(\{X_i\}_\mathbf{I}, \{Y_j\}_\mathbf{J}\right) &:= \lim_{i \in \mathbf{I}}\colim_{j \in \mathbf{J}} \Hom_\~C(X_i, Y_j) \\ 
    &=\lim_{i \in \mathbf{I}}\Hom_{\Ind(\~C)}\left(X_i, \colim_{j \in \mathbf{J}} Y_j\right).
    \end{align}
    Equivalently, $\Ind(\~C)$ is the full subcategory of $\Fun\left(\~C^\op, \Set\right)$ consisting of those presheaves $X$ which are small filtered colimits of representables, meaning that there is a filtered diagram $\{X_i\}_{i \in \mathbf{I}}$ with
    \[X \cong \colim_{i \in \mathbf{I}}H_\bullet(X_i)\]
    (where recall $H_\bullet$ denotes the Yoneda embedding). An ind-object $\{X_i\}_{\mathbf{I}}$ is called \emph{strict} if $\mathbf{I}$ can be chosen to be the posetal category $(\N, \le)$. 
\end{defn}

\begin{rmk}
Note the similarity between the second characterization of $\Ind(\~C)$ and the co-Yoneda lemma.
\end{rmk}

By identifying an object $X \in \~C$ with its singleton diagram $\ast \mapsto X$, we obtain an embedding $\~C \mono \Ind(\~C)$, which clearly factors through the Yoneda embedding:
\[\begin{tikzcd}
    & {\Ind(\~C)} \\
    {\~C} & {\Fun(\~C^\op,\Set)}
    \arrow[hook, from=1-2, to=2-2]
    \arrow[hook, from=2-1, to=1-2]
    \arrow["{H_\bullet}"', hook, from=2-1, to=2-2]
\end{tikzcd}\]
\part{Classical categorical structures}

\chapter{Monoidal categories}
% In the second half of this book, we will move past the basic definitions of category theory and explore more sophisticated categorical structures that one often encounters in mathematics. These are special kinds of categories constructed by placing additional hypotheses on the objects and arrows of an existing category.

\section{Definition and examples}

Roughly speaking, a category is \emph{monoidal} when it has a ``product'' operation, like the direct product $\times$, the direct sum $\oplus$, or the tensor product $\otimes$. In a general monoidal category, we write $\otimes$ for this product. 

\begin{defn}
    A \emph{monoidal category} is a category $\~C$ equipped with the following data:
    \begin{enumerate}
        \item A functor $\otimes: \~C \times \~C \to \~C$ (called the \emph{tensor product functor} and written $\otimes(X, Y) = X \otimes Y$).
        \item An object $1 \in \~C$ (called the \emph{tensor unit}) 
        \item A natural isomorphism $\alpha$ (called an \emph{associativity constraint}) with components 
        \[\alpha_{X, Y, Z}: X \otimes (Y \otimes Z) \isom (X \otimes Y) \otimes Z.\]
        \item A natural isomorphism $\rho: (-) \otimes 1 \isom (-)$ (called the \emph{right unitor}), with components
        \[\rho_X : X \otimes 1 \isom X.\]
        \item A natural isomorphism $\lambda: 1 \otimes (-) \isom (-)$ (called the \emph{left unitor}), with components
        \[\lambda_X: 1 \otimes X \isom X.\]
    \end{enumerate}
    This data must satisfy the following identities:
    \begin{enumerate}
        \item Pentagon axiom: For all $W, X, Y, Z \in \~C$, the following diagram commutes:
        \begin{equation}\label{eq:pentagon}
        \begin{tikzcd}[column sep=tiny,row sep=scriptsize]
            & {W \otimes(X \otimes(Y \otimes Z))} \\
            {W \otimes ((X \otimes Y) \otimes Z)} && {(W \otimes X) \otimes (Y \otimes Z)} \\
            \\
            {(W \otimes (X \otimes Y)) \otimes Z} && {((W \otimes X) \otimes Y)\otimes Z}
            \arrow["{1 \otimes \alpha}"', from=1-2, to=2-1]
            \arrow["\alpha", from=1-2, to=2-3]
            \arrow["\alpha"', from=2-1, to=4-1]
            \arrow["\alpha", from=2-3, to=4-3]
            \arrow["{\alpha \otimes 1}"', from=4-1, to=4-3]
        \end{tikzcd}
        \end{equation}
        \item 
    \end{enumerate}
\end{defn}




\begin{defn}
    A monoidal category is called \emph{strict} when every natural \emph{strict monoidal category} is a triple $(\~C, \otimes, I)$ consisting of a category $\~C$, a bifunctor $\otimes: \~C \times \~C \to \~C$, and an object $I \in \~C$ (called the \emph{monoidal unit}) such that for all objects $A, B, C$ and arrows $f, g, h$ in $\~C$, we have
    \begin{enumerate}[label=(\roman*)]
        \item $(A \otimes B) \otimes C = A \otimes (B \otimes C)$
        \item $I \otimes A = A$
        \item $A \otimes I = A$
        \item $(f \otimes g) \otimes h = f \otimes (g \otimes h)$
        \item $1_I \otimes f = f$
        \item $f \otimes 1_I = f$
    \end{enumerate}
\end{defn}

\chapter{Enrichment}
\label{chap:enrichment}

\chapter{Abelian categories}
\label{chap:abelian_cats}
These categories abstract special properties that hold in categories like $\Ab, \Mod_R,$ and $_{R}\Mod$. Namely, in each of these categories, the $\Hom$-sets are abelian groups and composition is bilinear. Moreover, in an abelian category, all finite limits and colimits exist and are particularly well-behaved. As a consequence of this, kernels and cokernels have incredibly special properties in abelian categories, and the study of these categories is incredibly elegant as a result.   

\section{Kernels and cokernels}
Recall that a \emph{zero} or \emph{null} object in a category $\~C$ is an object $0 \in \~C$ which is both initial and terminal. If $\~C$ has a zero object, then for all $A, B \in \~C$ the unique arrows $A \to 0$ and $0 \to B$ have a composite $0_{AB}: A \to B$ called the \emph{zero arrow} from $A$ to $B$. Since any two terminal/initial objects in the same category are isomorphic, the composition $A \to 0 \to B$ is independent of the choice of zero object. Moreover, any composite of arrows with at least one factor a zero will be equal to another zero arrow, since if $f: B \to C$ in $C$ then we have the composite
\[\begin{tikzcd}[column sep =small]
    A & 0 & B & C
    \arrow[from=1-1, to=1-2]
    \arrow[from=1-2, to=1-3]
    \arrow["{\exists!}"', bend right, dashed, from=1-2, to=1-4]
    \arrow["f", from=1-3, to=1-4]
\end{tikzcd}\]
but there can be only one map from $0$ to $C$, so the composite $0 \to B \To{f} C$ must equal the unique map $0 \to C$, so the above diagram reduces to $0_{AC}: A \to 0 \to C$, i.e., $0_{AB} f = 0_{AC}$, and a more general result holds for any sequence of arrows where at least one factor is a zero arrow. 

\begin{defn}
    Suppose $\~C$ is a category with a zero object. 
    \begin{enumerate}[label=(\roman*)]
    \item A \emph{kernel} of an arrow $f: A \to B$ is defined to be the equalizer of the diagram $\begin{tikzcd}[column sep = small]
        A & B
        \arrow["0"', shift right, from=1-1, to=1-2]
        \arrow["f", shift left, from=1-1, to=1-2]
    \end{tikzcd}$. That is, a kernel of $f: A \to B$ is an object $K \in \~C$ and a map $k: K \to A$ such that $k f = 0_{KA}$, and for every $g: K' \to A$ with $gf = 0$, there is a unique map $g': K' \to K$ such that the diagram
    \[\begin{tikzcd}[row sep = small]
        K \\
        & A & B \\
        K'
        \arrow["k"{description}, from=1-1, to=2-2]
        \arrow["0", shift left, from=1-1, to=2-3]
        \arrow["f"{description, pos=0.4}, from=2-2, to=2-3]
        \arrow["{\exists! \, g'}", dashed, from=3-1, to=1-1]
        \arrow["g"{description}, from=3-1, to=2-2]
        \arrow["0"', shift right, from=3-1, to=2-3]
    \end{tikzcd}\]
    commutes. 
    \item A \emph{cokernel} of an arrow $f: A \to B$ is a kernel of $f$ in $\~C^\op$. Explicitly, this means that a cokernel is an object $C \in \~C$ and a map $c: B \to C$ such that $fc = 0_{AC}$, and for every $h: C' \to A$ with $fh = 0$, there is a unique map $h': C \to C'$ such that
    \[\begin{tikzcd}[row sep = small]
        && C \\
        A & B \\
        && {C'}
        \arrow["{\exists ! \,h'}", dashed, from=1-3, to=3-3]
        \arrow["0", shift left, from=2-1, to=1-3]
        \arrow["f"{description}, from=2-1, to=2-2]
        \arrow["0"', shift right, from=2-1, to=3-3]
        \arrow["{c}"{description}, from=2-2, to=1-3]
        \arrow["h"{description}, from=2-2, to=3-3]
    \end{tikzcd}\]
    commutes.
\end{enumerate}
\end{defn}

So a category $\~C$ has kernels (resp. cokernels) if it has all equalizers (resp. coequalizers) and a zero object. Moreover, since a kernel (resp. cokernel) is an equalizer (resp. coequalizer), it is defined up to a unique isomorphism, so we may speak of \emph{the} kernel/cokernel of a map $f$. As is the case for all equalizers (Proposition \ref{prop:equalizer_monic}), the map $k: K \to A$ is monic and the map $c B \to C$ is epic. Thus, we can think of the kernel $k: K \to A$ as a subobject of $A$, i.e., an equivalence class of monics $K \mono A$. For this reason, we typically write $\ker f$ for the arrow $k: K \to A$. Similarly, we write $\coker f$ for the arrow $c: B \to C$.


\begin{defn}
    A \emph{preadditve category} or \emph{$\Ab$-category }is a category $\~A$ that is enriched over the category $\Ab$ of abelian groups. Explicitly, this means that $\~A$ satisfies the following two conditions:
    \begin{enumerate}[label=(\roman*)]
        \item for any pair of objects $A, B \in \~A$, the hom-set $\Hom(A, B)$ has the structure of an additive abelian group;
        \item composition of morphisms distributes over the addition of morphisms given by the abelian group structure, i.e., if $k \in \Hom(A, B)$, $g, h \in \Hom(B, C)$ and $f \in \Hom(C, D)$, then
        \[f  (g + h) = f  g + f  h \quad \text{and} \quad (g + h)  k = g  k + h  k.\]
    \end{enumerate}
\end{defn}

\begin{ex}
    In the category $\Grp$, the trivial group is the zero object, and for any two abelian groups $G$ and $H$, the zero map $0_{GH}: G \to H$ is the unique map which sends every element in $G$ to the identity of $H$. In this case, the kernel of a homomorphism $f: G \to H$ is the pair $(K, K \to G)$, where $K = \{x \in G : f(x) = 1_H\}$ is the usual kernel from group theory and the map $K\to G$ is the inclusion. We use the same definition of kernel in the category $\Ab$ of abelian groups. In $\Grp$, the cokernel object associated to $f: G \to H$ is the quotient $H/f(G)$ of $H$ by the image of $f$, $H \to C$ is the usual projection from a group to its quotient.
\end{ex}

\begin{ex}
    Recall that $\Set_*$ denotes the category of pointed sets, which we introduced in Exercise X . In $\Set_*$, the one-point set is a zero object, and given pointed sets $P$ and $Q$, the zero map $P \to Q$ is the function taking all of $P$ to the base point $*_Q$ in $Q$. If $f: P \to Q$ is a morphism of pointed sets, the kernel $K \to P$ is given by the subset 
    \[K = \{x \in P : f(x) = *_Q\}\]
    where the base point in $K$ is the same as the base point in $P$. We can use a similar method to define kernels in the category $\Top_\ast$ of pointed topological spaces. Interestingly, in $\Grp$, an epimorphism is determined (up to isomorphism) by its kernel, but this is not at all the case in $\Set_*$ or $\Top_*$.
\end{ex}

\begin{ex}
    In any preadditive category $\~A$, all equalizers are kernels. In such a category each $\~A(B, C)$ is an abelian group, so given a parallel pair of arrows $\begin{tikzcd}[column sep = small]
        B & C
        \arrow["g"', shift right, from=1-1, to=1-2]
        \arrow["f", shift left, from=1-1, to=1-2]
    \end{tikzcd}$, a third arrow $h: A \to B$ satisfies $hf = hg$ if and only if $(h(f - g))(a) = (f - g)(h(a)) = 0$ for all $a \in A$. Hence $hf = hg$ if and only if $h(a) \in \ker(f - g)$ for all $a \in A$, so an equalizer of $f$ and $g$ can either be described as the limit of $\begin{tikzcd}[column sep = small]
        B & C
        \arrow["g"', shift right, from=1-1, to=1-2]
        \arrow["f", shift left, from=1-1, to=1-2]
    \end{tikzcd}$ or the kernel of $f - g$. 
\end{ex}

If a category $\~C$ has a null object, kernels and cokernels, then we can factor any arrow $f: A \to B$ as $f = qk_{c_f}$, where $q$ is the unique arrow making the following diagram commute:
\[\begin{tikzcd}[row sep = small]
    {\ker c_f} \\
    & B & {\coker f} \\
    A
    \arrow["{k_{c_f}}", from=1-1, to=2-2]
    \arrow["{c_f}", from=2-2, to=2-3]
    \arrow["{\exists ! \, q}", dashed, from=3-1, to=1-1]
    \arrow["f"', from=3-1, to=2-2]
\end{tikzcd}\]


\begin{lem}
Let $\~C$ be a category with a null object, kernels and cokernels, and let $f: A \to B$ be an arrow with factorization $f = qk_{c_f}$ as above. Suppose further that $g: B \to C$ is an arrow with kernel $k_g: \ker g \to B$ such that $f = q' k_g$. Then the commutative square
\begin{equation}
\label{eq:kernel_square}
\begin{tikzcd}
    A & {\ker c_f} \\
    {\ker g} & B
    \arrow["q", from=1-1, to=1-2]
    \arrow["{q'}"', from=1-1, to=2-1]
    \arrow["t"', dashed, from=1-2, to=2-1]
    \arrow["{k_{c_f}}", from=1-2, to=2-2]
    \arrow["{k_g}"', from=2-1, to=2-2]
\end{tikzcd}
\end{equation}
has a unique diagonal arrow $t: \ker c_f \to \ker g$ such that $tk_g = k_{c_f}$ and $qt = q'$. Moreover, if $\~C$ has all equalizers and every monic in $\~C$ is a kernel, then $q$ is an epimorphism.
\end{lem}

    
\section{Additive categories}

\begin{prop}
    The following properties of an object $Z$ in a preadditive category $\~A$ are equivalent:
    \begin{enumerate}[label=(\roman*)]
        \item $Z$ is initial;
        \item $Z$ is terminal;
        \item the identity arrow $1_Z: Z \to Z$ is equal to zero map $0_{ZZ}: Z \to Z$: $1_Z = 0_{ZZ}$;
        \item the abelian group $\~A(Z, Z)$ is the zero group.
    \end{enumerate}
    In particular, any initial (or terminal) object in $\~A$ is a null object.
\end{prop}

\begin{proof}
    If $Z$ is initial, then there is only one map $Z \to Z$, which must be the identity. This gives the implications (i) $\implies$ (iii) and (i) $\implies$ (iv). Now, if $1_Z = 0_{ZZ}$, then any $f: A \to Z$ has $f = f1_Z = f0_{ZZ} = 0_{BZ}$, so there is only one arrow in $\~A(A, Z)$, namely the zero arrow $0_{AZ}: A \to Z$ so $Z$ is terminal. This gives (iii) $\implies$ (ii) and (iii) $\implies$ (iv). The rest follows by duality.
\end{proof}

When there is a zero object $0$ in the preadditive category $\~A$, the unique maps $A \to 0$ and $0 \to B$ are the zero elements of $\~A(A, 0)$ and $\~A(0, B)$ respectively. Hence the composite $A \to 0 \to B$ is both the zero morphism $A \to B$ and the zero element of the abelian group $\~A(A, B)$.

\begin{defn}
    Let $\~A$ be a preadditive category. A \emph{biproduct diagram} for the objects $A, B \in \~A$ is a diagram 
    \begin{equation}
    \label{eq:biprod_diagram}
        \begin{tikzcd}
            A & C & B
            \arrow["{i_1}"', shift right, from=1-1, to=1-2]
            \arrow["{p_1}"', shift right, from=1-2, to=1-1]
            \arrow["{p_2}", shift left, from=1-2, to=1-3]
            \arrow["{i_2}", shift left, from=1-3, to=1-2]
        \end{tikzcd}
    \end{equation}
with arrows $p_1, p_2, i_1, i_2$ such that 
\begin{equation}
\label{eq:biproduct_rules}
i_1 p_1 = 1_A, \quad i_2 p_2 = 1_B, \quad p_1 i_1 + p_2 i_2 = 1_C
\end{equation}
\end{defn}

\begin{thm}
    Two object $A$ and $B$ in a preadditive category $\~A$ have a product in $\~A$ if and only if they have a biproduct in $\~A$. Specifically, given a biproduct diagram (\ref{eq:biprod_diagram}), the object $C$ with the projections $p_1$ and $p_2$ is a product of $A$ and $B$, while dually, $C$ with $i_1$ and $i_2$ is a coproduct of $A$ and $B$. In particular, two objects $A, B \in \~A$ have a product if and only if they have a coproduct.
\end{thm}

\begin{proof}
    First assume that we have a biproduct diagram (\ref{eq:biprod_diagram}) satisfying the condition (\ref{eq:biproduct_rules}). Then 
    \begin{align*}
        i_2p_1 &= i_2(p_1 i_1 + p_2 i_2) p_1 \\
               &= (i_2 p_1 i_1 + i_2 p_2 i_2)p_1 \\
               &= i_2 p_1 i_1 p_1 + i_2 p_2 i_2 p_1 \\
               &= i_2 p_1 1_A + 1_B i_2 p_1 \\
               &= i_2 p_1 + i_2 p_1
    \end{align*}
    and subtracting $i_2 p_1$ from both sides yields $i_2 p_1 = 0$; the fact that $i_1 p_2 = 0$ follows symmetrically from the argument above. Our goal is to show that the diagram 
    \[\begin{tikzcd}
        & C \\
        A && B
        \arrow["{p_1}"', from=1-2, to=2-1]
        \arrow["{p_2}", from=1-2, to=2-3]
    \end{tikzcd}\]
    is a product, i.e., for any diagram
    \begin{equation}
    \label{eq:span}
    \begin{tikzcd}
    & D \\
    A && B
    \arrow["{f_1}"', from=1-2, to=2-1]
    \arrow["{f_2}", from=1-2, to=2-3]
    \end{tikzcd}
    \end{equation}
    there exists a unique arrow $h: D \to C$ rendering the diagram
    \begin{equation}
    \label{eq:abproduct}
    \begin{tikzcd}
        &&& A \\
        D && C \\
        &&& B
        \arrow["{f_1}", from=2-1, to=1-4]
        \arrow["h"{description}, dashed, from=2-1, to=2-3]
        \arrow["{f_2}"', from=2-1, to=3-4]
        \arrow["{p_1}"', from=2-3, to=1-4]
        \arrow["{p_2}", from=2-3, to=3-4]
    \end{tikzcd}\end{equation}
    commutative. So, consider a diagram of the form (\ref{eq:span}), 
    and set $h =  f_1 i_1 +  f_2 i_2: D \to C$. Then we have
    \begin{align*}
        hp_1 &=  (f_1i_1 + f_2 i_2)p_1 \\
                    &= f_1i_1p_1 + f_2 i_2p_1  \\
                    &= f_1 1_A + f_2 0 \\
                    &= f_1,
    \end{align*}
    and similarly, $hp_2 = f_2$, so $h$ makes the diagram (\ref{eq:abproduct}) commute. If another arrow $h': D \to C$ also makes (\ref{eq:abproduct}) commute, then we have
    \begin{align*}
        h' &= h'(p_1 i_1 + p_2 i_2) \\
        &= h'p_1 i_1 + h'p_2 i_2\\
        &= f_1i_1 + f_2i_2 \\
        &= h.
    \end{align*}
    Thus, $h: D \to C$ is the unique arrow making (\ref{eq:abproduct}) commute, so the objects $A$ and $B$ have a product in $\~A$. The assignment $h \mapsto (f_1, f_2)$ is an isomorphism 
    \[\~A(D, C) \cong \~A(D, A) \oplus \~A(D, B),\]
    where $\oplus$ denotes the direct sum of abelian groups.
\end{proof}

In categories like $\Ab$ and $\Mod_R$, the biproduct of $A$ and $B$ is often called a \emph{direct sum} and written $A \oplus B$. In these categories, (and indeed, in any preadditive category), the product and coproduct functors are naturally isomorphic (cf. \cite{maclane}, p. 196), and can therefore both be identified with the direct sum $\oplus$ (this explains the terminology ``biproduct''). Moreover, given objects $A_1, \cdots, A_n \in \~A$, we can construct the $n$-fold biproduct $\bigoplus_{j = 1}^n A_j$, which is characterized by the diagram 
\[\begin{tikzcd}
    {A_j} & {\bigoplus_{j = 1}^n A_j} & {A_k}
    \arrow["{i_j}", from=1-1, to=1-2]
    \arrow["{p_k}", from=1-2, to=1-3]
\end{tikzcd}, \quad j, k = 1, \hdots, n \]
and the equations
\[i_1 p_1 + \cdots + i_np_n = 1, \quad p_k i_j = \delta_{kj}\]
where $\delta_{jk}$ denotes the Kronecker delta. 

\begin{rmk}
Older texts often may also refer to this biproduct as an ``internal'' direct sum, because the construction of the biproduct relies entirely on the data $\begin{tikzcd}[sep = scriptsize]
            A & C & B,
            \arrow["{i_1}"', shift right, from=1-1, to=1-2]
            \arrow["{p_1}"', shift right, from=1-2, to=1-1]
            \arrow["{p_2}", shift left, from=1-2, to=1-3]
            \arrow["{i_2}", shift left, from=1-3, to=1-2]
\end{tikzcd}$and makes no assumptions about other objects or arrows in the underlying category $\~A$. This is in contrast to the usual product or ``external'' direct sum (again, an older term that we will not adopt here), whose definition depends on objects and arrows in the whole category. 
\end{rmk}
 

\begin{defn}\label{def:additive category}
    An \emph{additive category} is a preadditive category $\~A$ satisfying the following additional conditions:
    \begin{enumerate}[label=(\roman*)]
        \item $\~A$ has a zero object $0$ such that $\Hom(0, 0)$ is the trivial group;
        \item  there exists a biproduct diagram (\ref{eq:biprod_diagram}) for any pair of objects $A, B \in \~A$.
    \end{enumerate}
\end{defn}
Having defined (pre)additive categories, we may now define \emph{additive functors}, which are precisley those functors that preserve the additive structure:
\begin{defn}
\label{defn:adfunctor1}
    If $\~A$ and $\~B$ are preadditive categories, an \emph{additive functor} is a functor $T: \~A \to \~B$ such that for all $f, f':A \to B$ in $\~A$ we have
    \[T(f + f') = T(f) + T(f').\]
    That is, $T$ is additive if the induced maps $\Hom(A, B) \to \Hom(T(A), T(B))$ are group homomorphisms.
\end{defn}

If $\~A$ is an additive category, then all of $\~A$'s additive structure can be phrased in terms of biproducts, so we can reformulate the definition of an additive functor in terms of biproduct diagrams. This is the main content of the following proposition:
\begin{prop}
    If $\~A$ and $\~B$ are preadditive categories and $\~A$ has a biproduct for each pair of objects, then a functor $T: A \to B$ is additive if and only if $T$ carries biproduct diagrams in $\~A$ to biproduct diagrams in $\~B$.
\end{prop}

\section{Abelian categories}

\begin{defn}
    An \emph{abelian category} $\~A$ is an $\Ab$-category that satisfies the following conditions:
    \begin{enumerate}[label=(\roman*)]
        \item $\~A$ has a null object.
        \item $\~A$ has binary biproducts.
        \item Every arrow in $\~A$ has a kernel and cokernel.
        \item Every monic arrow is a kernel, and every epi is a cokernel.
    \end{enumerate}
\end{defn}

The first two conditions say that $\~A$ is an additive category. The third condition, when combined with the first two, implies that $\~A$ has finite limits. For if $f, g: A \to B$ are arrows in $\~A$, we can construct an equalizer of $f$ and $g$ as the kernel of $f - g$:
\[\ker(f-g) = \lim(X \rightrightarrows Y),\]
and finite products and equalizers give all finite limits. The final condition is the strongest of the four. For instance, it implies that any arrow which is both monic and epic is an isomorphism. For if $f: A \mono B$ is monic, then $f = \ker g$ for some $g$, hence $fg = 0 = f0$. But since $f$ is an epimorphism, we can cancel $f$ to get $g = 0: B \to C$. And $\ker(0_{BC}) = 1_B$, hence $f$ is an isomorphism.

\begin{prop}[Mono-epi factorization]
    In an abelian category $\~A$, every arrow $f:A \to B$ has a factorization $f = em$, with $e: A \epi C$ epic and $m: C \mono B$ monic and
    \[m = \ker(\coker f), \quad e = \coker(\ker f).\] 
    Furthermore, given any other factorization $f' = m'e'$ with $m'$ monic and $e'$ epic, any commutative square
    \[\begin{tikzcd}[column sep=scriptsize]
    \cdot & \cdot \\
    \cdot & \cdot
    \arrow["f", from=1-1, to=1-2]
    \arrow["g"', from=1-1, to=2-1]
    \arrow["h", from=1-2, to=2-2]
    \arrow["{f'}"', from=2-1, to=2-2]
\end{tikzcd}\]
extends to a commutative rectangle of the form
\[\begin{tikzcd}[column sep=scriptsize]
    \cdot & \cdot & \cdot \\
    \cdot & \cdot & \cdot
    \arrow["e", two heads, from=1-1, to=1-2]
    \arrow["f", shift left=5, from=1-1, to=1-3]
    \arrow["g"', from=1-1, to=2-1]
    \arrow["m", hook, from=1-2, to=1-3]
    \arrow["k", from=1-2, to=2-2]
    \arrow["h", from=1-3, to=2-3]
    \arrow["{e'}"', two heads, from=2-1, to=2-2]
    \arrow["{f'}"', shift right=5, from=2-1, to=2-3]
    \arrow["{m'}"', hook, from=2-2, to=2-3]
\end{tikzcd}\]
\end{prop}

\chapter{Derived functors}

In this chapter we will further develop the theory of abelian categories by placing more sophisticated constructions from homological algebra in the context of an abelian category. 
% Two of the most fundamental constructions in commutative algebra are the  hom-functor $\Hom$ and the tensor product $\otimes$ functor. These functors have many desirable properties 

\chapter{Triangulated categories}
Our goal in this chapter is to introduce triangulated categories, which are one of the most widely used types of ``categories with extra structure''. While these categories arise in some very sophisticated mathematical contexts, such as $K$-theory and the theory of motives, the definition of a triangulated category, while lengthy, requires no special categorical machinery to state. Triangulated categories were introduced by Verdier\footnote{When discussing triangulated categories, Verdier is more commonly known as ``Grothendieck's student Verdier''.} in his thesis \cite{verdier_thesis}. Much of the modern theory of triangulated categories was developed by Amnon Neeman, whose book on the subject \cite{neeman_triangulated} serves as the main reference for this chapter.
\section{Basic definitions} 
The full definition of a triangulated category requires a lot of data, so it will be useful to first describe the simpler notion of a \emph{pre-triangulated category}. Before giving the definition proper, we introduce some terminology that will be used throughout the rest of this chapter. Let $\~T$ be an additive category (Definition \ref{def:additive category}) and let $\Sigma: \~T \to \~T$ be an additive functor which is also an equivalence (not surprisingly, such a functor is called an \emph{additive auto-equivalence}). A diagram of the form
\[A \to B \to C \to \Sigma A\]
will be called a \emph{triangle in $\~T$}. We may also depict these triangles using a diagram of the form
\[\begin{tikzcd}
    A && B \\
    & C
    \arrow[from=1-1, to=1-3]
    \arrow[from=1-3, to=2-2]
    \arrow[squiggly, from=2-2, to=1-1]
\end{tikzcd}\]
where $C \rightsquigarrow A$ denotes the morphism $C \to \Sigma A$. While this is the convention we adopt in this text, we caution the reader that there is no standardized way to represent triangles in $\~T$, and different authors may have different notational preferences (for instance, Bridgeland draws his triangles using dotted arrows instead of squiggly arrows). Given two triangles $A \to B \to C \to \Sigma A$ and $A' \to B' \to C' \to \Sigma A'$ in $\~T$, a \emph{morphism of triangles} is a collection of four maps $(u, v, w, \Sigma u)$ which make the diagram
\[\begin{tikzcd}
    A & B & C & {\Sigma A} \\
    {A'} & {B'} & {C'} & {\Sigma A'}
    \arrow[from=1-1, to=1-2]
    \arrow["u"', from=1-1, to=2-1]
    \arrow[from=1-2, to=1-3]
    \arrow["v", from=1-2, to=2-2]
    \arrow[from=1-3, to=1-4]
    \arrow["w", from=1-3, to=2-3]
    \arrow["{\Sigma u}", from=1-4, to=2-4]
    \arrow[from=2-1, to=2-2]
    \arrow[from=2-2, to=2-3]
    \arrow[from=2-3, to=2-4]
\end{tikzcd}\]
commute.
\begin{defn}
    A \emph{triangulated category} is a triple $(\~T, \Sigma, \Delta)$ where $\~T$ is additive, $\Sigma: \~T \to \~T$ is an additive auto-equivalence, and $\Delta$ is a collection of triangles in $\~T$. This data is subject to the following axioms:
    \begin{itemize}[leftmargin=1.5cm]
        \item[(TR0)] Any triangle which is isomorphic to a triangle in $\Delta$ must also be in $\Delta$.
        \item[(TR1)]The triangle $A \To{1_A} A \to 0 \to \Sigma A$ is in $\Delta$.
    \end{itemize}
\end{defn}

% \chapter{Fibred categories}

% \chapter{Bicategories, 2-categories, and beyond}

% \part{Categories of Sheaves}

% \chapter{Sheaves \& Sites}

\chapter{Topoi}

% The third and final chapter of this text is intended to be an introduction to the area of \emph{topos theory}. Roughly speaking, a \emph{topos} is a special category that behaves like the category $\Set$ of sets, or, more generally, the category of sheaves of sets on a topological space. 
\section{Elementary topoi}

\section{Sheaves} 
 
Let $\~C$ be a category. Recall that a \textit{presheaf} on $\~C$ is a functor$\~F: \~C^\op \to \Set$. This means that a presheaf on $\~C$ contains the following data:
\begin{itemize}
\item For each $A \in \~C$, a set $\~F(A)$.
\item For each arrow $B \to A$ in $\~C$, a morphism $\rho_{AB}: \~F(A) \to \~F(B)$.
\end{itemize}
This data is subject to the following conditions:
\begin{enumerate}[label=(\roman*)]
\item The map $\rho_{AA}: \~{F}(A) \to \~{F}(A)$ is the identity map $1_{\~{F}(A)}: \~{F}(A) \to \~{F}(A)$ ($\rho_{AA} = 1_{\~{F}(A)}$).
\item Given a commutative diagram
\[\begin{tikzcd}
        C && {B} \\
        & {A}
        \arrow[from=1-1, to=1-3]
        \arrow[from=1-3, to=2-2]
        \arrow[from=1-1, to=2-2]
    \end{tikzcd}\]
in $\~C$, the diagram
\[\begin{tikzcd}
    {\~{F}(C)} && {\~{F}(B)} \\
    & {\~{F}(A)}
    \arrow["{\rho_{AC}}", from=2-2, to=1-1]
    \arrow["{\rho_{BC}}"', from=1-3, to=1-1]
    \arrow["{\rho_{AB}}"', from=2-2, to=1-3]
\end{tikzcd}\]
commutes in $\Set$, i.e., $\rho_{AC} = \rho_{BC} \circ \rho_{AB}$.
\end{enumerate}
We write $\PSh(\~C) = \Fun(\~C^\op, \Set)$ for the category of presheaves on $\~C$. If $f \in \~{F}(A)$, we call $f$ a \emph{section of $\~{F}$ over $A$} and we sometimes denote $\rho_{AB}(f)$ by $f|_B$. The maps $\rho_{AB}$ are called \emph{restriction maps}; other authors prefer to use the notation $\mathrm{res}_{AB}$ instead of $\rho_{AB}$ for obvious reasons. The reason for this terminology is explained in Example \ref{ex:presheaf_cont_functions} below.

We now fix a topological space $X$, and denote by $\Open(X)$ the partially ordered set of open subsets of $X$, ordered by inclusion. View $\Open(X)$ as a posetal category in the usual manner, by drawing a unique arrow $V \to U$ in $\Open(X)$ when $V \subseteq U$ in $X$. A \emph{presehaf on $X$} is a presheaf on the category $\Open(X)$. 

\begin{ex}[Presheaf of continuous $\R$-valued functions]
\label{ex:presheaf_cont_functions}
We can define a presheaf of rings $C$ on $X$ as follows: The assignment of an open subset $U \subseteq X$ to a ring $C(U)$ is 
    \[C(U) = \{f: U \to \R : f \text{ continuous} \}\]
where each $C(U)$ is endowed with the structure of a commutative ring by defining addition and multiplication pointwise:
\[(f + g)(x) = f(x) + g(x) \quad \text{and} \quad fg(x) = f(x)g(x) \quad \text{for all } f, g \in C(U), \, x \in U\]
We now have a ring $C(U)$ for every open set in $X$, so to ensure that $C$ defines a sheaf on $X$, we just need to specify where $C$ sends a map $V \to U$ (i.e., an inclusion $V \subseteq U$). Suppose we have $V \subseteq U$ in $\Open(X)$. Then every continuous function $U \to \R$ is can be restricted to a continuous function $V \to \R$, however there may be functions defined on $V$ but not $U$. For instance, the real-valued function $f(x) = 1/x$, is not defined on the entire real line $\R$, but is well-defined on the open subset $\R \setminus \{0\} \subset \R$. More generally, if $V \subseteq U$ in $\Open(X)$, then there are more continuous functions $V \to \R$ than there are continuous functions $U \to \R$. We therefore obtain a \emph{reverse inclusion of sets}
\[C(V) \supseteq C(U)\] 
Given that $C(U)$ is entirely contained within $C(U)$, there is a very natural choice of ring homomorphism $C(U) \to C(V)$, given by
\begin{align*}
    \rho_{UV}: C(U) &\to C(V) \\
    f &\mapsto f|_V
\end{align*}
To spell this out, if $f: U \to \R$ is in $C(U)$, then the map $\rho_{UV}$ sends $f: U \to \R$ to $f|_V: V \subseteq U \to \R$, where $f|_V = \rho_{UV}(f)$ has the same definition as $f$, only with a smaller domain (namely, the subset $V$ of the original domain $U$). The map $f|_V$ is called the \emph{restriction of $f$ to $V$}; this is why we refer to the maps $\rho_{UV}$ as ``restriction maps'', and use the notation $f|_V$. We can now easily check that $\rho_{UV}$ defines a ring homomorphism $C(U) \to C(V)$. Finally, to show functoriality of $C$, we first note that if $f \in C(U)$ then clearly $\rho_{UU}(f) = f|_U = f$, so we have  
\[\rho_{UU} = 1_{C(U)}\]
so $\~F$ preserves identities. For composition, suppose we have a chain of inclusions $W \subseteq V \subseteq U$ in $\Open(X)$. It is clear that the restriction $\left(f|_V \right)|_W$ is exactly the same as the restriction $f|_W$, so we have a commutative triangle: 
\[\begin{tikzcd}
    {C(W)} && {C(V)} \\
    & {C(U)}
    \arrow["{\rho_{UW}}", hook', from=2-2, to=1-1]
    \arrow["{\rho_{VW}}"', hook', from=1-3, to=1-1]
    \arrow["{\rho_{UV}}"', hook, from=2-2, to=1-3]
\end{tikzcd}\]
Hence $C$ defines a presheaf of rings on $X$.
\end{ex}

A related piece of terminology that we will frequently make use of is the following:
\begin{defn}
    Let $\~F: \~C^\op \to \Set$ be a presheaf on the category $\~C$. A \emph{subfunctor} of $\~F$ is a subobject of $\~F$ in $\PSh(\~C)$ A presheaf $\~G: \~C^\op \to \Set$ is called a \emph{subfunctor} of $\~F$ (written $\~G \subseteq \~F$) if $\~G(A) \subseteq \~F(A)$ for all $A \in \~C$, and for all arrows $f: B \to A$ in $\~C$, $\~G(f)$ is the restriction of $\~F(f)$ to $\~G(C)$.
\end{defn}

Recall that for an open set $U$ in $X$, an \emph{open cover of $U$} is a collection $\{U_i\}_{i \in I}$ of open subsets of $U$ such that $ \bigcup_{i \in I} U_i = U$. In this situation we say that $\{U_i\}_{i \in I}$ \emph{covers $X$}. To turn a presheaf on $X$ into a sheaf on $X$, we will need to specify how functions defined on an open cover $\{U_i\}_{i \in I}$ of each open $U \subseteq X$ relate to functions on the whole space $U$, as well as what should happen to a function $f$ defined on the intersection $U_i \cap U_j$ of two sets in this open cover. 
\begin{defn}\label{defn:sheaf}
A presheaf $\~F$ on a topological space $X$ is called a \emph{sheaf} if it satisfies the following additional axioms for all 
open subsets $U$ of $X$
\begin{enumerate}[label=(\roman*)]
        \item\label{sheaf:identity} (\emph{Identity}) If $\{U_i\}_{i \in I}$ is an open cover of $U$ and $f, g \in \~F(U)$ are such that $f|_{U_i}= g|_{U_i}$ for all $i$, then $f = g$.
        \item\label{sheaf:gluing} (\emph{Gluing}) If $\{U_i\}_{i \in I}$ is an open cover of $U$ and we have sections $f_i \in \~F(U_i)$ for each $i$, with the property that for each $i, j \in I$ we have $f_i|_{U_i \cap U_j} = f_j|_{U_i \cap U_j}$, then there is an element $f \in \~F(U)$ such that $f|_{U_i} = f_i$ for all $i$ (note that condition (i) implies $f$ is unique).
\end{enumerate}
\end{defn}
\begin{ex}
\hfill
\begin{enumerate}[label=(\roman*)]
    \item The presheaf (of rings) $C$ in Example \ref{ex:presheaf_cont_functions} is also a sheaf of rings, called the \emph{sheaf of continuous functions} on $X$. When $X$ is a smooth manifold, we may similarly define the sheaf of smooth functions on $X$, and if $X$ is a complex manifold (for instance, a Riemann surface) we can define the sheaf of holomorphic functions on $X$. 

    \item Suppose $X$ is a topological space with $x \in X$, and let $S$ be any set. Let $i_x: \{x\} \to X$ be the inclusion. Then $i_{x, *}S$ defined by
    \[i_{x, *}S(U) = \begin{cases} S & \text{if } x \in U;\\
                            \{\ast\} & \text{if } x \not \in U
                     \end{cases}\]
    defines a sheaf, called the \emph{skyscraper sheaf}. There is an analogous definition for more general categories: the set $S$ can be replaced with whatever object we want, and the singleton $\{\ast\}$ should be replaced with the terminal object. The restriction maps are the obvious suspects: if $U \subseteq V$ and $x \in U$, then $x \in V$ as well, so $i_{x, *}S(U) = i_{x, *}S(V) = S$, and we set $\rho_{VU} = 1_S$, the identity map on $S$. Otherwise, if $U \subseteq V$ and $x \not \in U$, then $i_{x, *}S(U) = \{\ast\}$, and we set $\rho_{VU}: i_{x, *}(V) \to \{\ast\}$ equal to the only map it can possibly be. The name ``skyscraper sheaf'' is meant to suggest that all of the data of this sheaf occurs at the the point $x \in X$; the sheaf casts aside any open subsets which do not contain $x$ and assigns all of the subsets that do contain $x$ to the set $S$, as if $S$ is a skyscraper placed at the point $x$. 
\end{enumerate}
\end{ex}

\begin{ex} There are many presheaves that fail to satisfy the sheaf axioms. Our next two examples illustrate this.
    \begin{enumerate}[label=(\roman*)]
        \item For $U \subseteq \C$, define $C^b(U)$ to be the collection of bounded functions $f : U \to \C$. Then we have an assignment of a set to each open $U \subseteq \C$, and given $V \subseteq U \subseteq \C$, we may define a morphism $C^b(f): C^b(U) \to C^b(V)$ via the usual restriction of functions. In this way, $C^b$ defines a presheaf on the category $O(\C)$ of open subsets of $\C$, i.e., we have a functor $C^b: O(\C)^\op \to \Set$.

        We'll now show that $C^b$ does not satisfy gluing, and can therefore not be a sheaf. To that end, consider the covering $\C = \bigcup_{n \in \N} U_n$, where $U_n = \{z \in \C : |z| < n\}$ is the open disk of radius $n$. For each $n \in \N$, define
            \[f_n: U_n \to \C, \quad f_n(z) = z.\]
        Then certainly each $f_n$ is bounded on $U_n$, and for $m < n$, we have $U_m \subseteq U_n$, and 
            \[f_n|_{U_m}(z) = z = f_m(z),\]
        so we have $f_i|_{U_i \cap U_j} = f_j|_{U_i \cap U_j}$. Now, suppose there exists a global bounded holomorphic function $f \in C^b(\C)$ such that $f|_{U_n} = f_n$ for all $n$. Then each $U_i$ is a connected open subset of $\C$, so by the identity principle for holomorphic functions, we must have $f(z) = z$. But now $f(z) = z$ is \emph{not} bounded on all of $\C$, so we cannot glue the sections $f_n$ together to get a global section, hence $C^b$ cannot be a sheaf. 

        \item For $U \subseteq \C$, define $\~F(U)$ to be the set of holomorphic functions which admit a square root on $U$.
        \[\~F(U) = \{f: U \to \C :f = g^2 \text{ for some holomorphic function } g: U \to \C\}\] 
        Then we have an assignment of a set to each open $U \subseteq \C$, and given $V \subseteq U \subseteq \C$, we may define a morphism $\~F(f): \~F(U) \to \~F(V)$ via the usual restriction of functions. In this way, $\~F$ defines a presheaf on the category $O(\C)$ of open subsets of $\C$, i.e., we have a functor $\~F: O(\C)^\op \to \Set$.

        We again claim that the presheaf $\~F$ fails the gluing axiom. Let $U = \C \setminus \{0\}$, and consider the open cover $U = U_1 \cup U_2$ where $U_1 = \C \setminus \R_{\le 0}$ and $U_2 = \C \setminus \R_{\ge 0}$. Define $f: U \to \C$ by $f(z) = z$. On $U_1$, $f$ has a holomorphic square root $g_1(z)  =e^{\frac{1}{2} \log z}$ where $\log z$ is a principal branch of the logarithm that is holomorphic on $U_1$. On $U_2$, $f$ has a holomorphic square root $g_2 = e^{\frac{1}{2} (\log z \pi i)}$. Then $g_1|_{U_1 \cap U_2} = -g_2|_{U_1 \cap U_2}$, but both $g_1$ and $-g_2$ are valid square roots of $f(z) = z$. If $\~F$ were a sheaf, we'd be able to glue $g_1$ and $g_2$ into a single holomorphic square root $g \in \~F(U)$ with $g|_{U_1} = g_1$ and $g|_{U_2} = g_2$. But no such $g$ exists, since we just saw that $g_1$ and $g_2$ differ by a sign on $U_1 \cap U_2$. Hence $\~F$ does not satisfy the axioms of a sheaf.
    \end{enumerate}
\end{ex}

We now unpack the axioms in Definition \ref{defn:sheaf} more carefully. To keep the discussion concrete, we first consider the sheaf $C$ of Example \ref{ex:presheaf_cont_functions}. Suppose we have an open set $U \subseteq X$ with a covering $U = \bigcup_{i \in I} U_i$ and an $I$-indexed family of functions $\{f_i: U_i \to \R\}_{i \in I}$. We can describe the family of maps $f_i: U_i \to \R$ as an element of the product $\prod_{i \in I}C(U_i)$ (since an element in this product is precisely an $I$-indexed family of maps on the sets $U_i$). Likewise, the assignments $\left(f_i\right)_{i \in I} \mapsto \left(f_i|_{U_i \cap U_j}\right)_{i, j \in I}$ and $\left(f_i\right)_{i \in I} \mapsto \left(f_j|_{U_i \cap U_j}\right)_{i, j \in I}$ define maps
\[\begin{tikzcd}
    {\prod_{i \in I}C(U_i)} & {\prod_{i,j \in I} C(U_i \cap U_j)}
    \arrow["p", shift left, from=1-1, to=1-2]
    \arrow["q"', shift right, from=1-1, to=1-2]
\end{tikzcd}\]
Now define a map $e: C(U) \to \prod_{i \in I} C(U_i)$ by sending each section $f \in \~F(U)$ to the family of maps $\left(f|_{U_i}\right)_{i \in I}$. The axioms in Definition \ref{defn:sheaf} are then equivalent to the assertion that the following diagram is an equalizer for all coverings $U = \bigcup_{i \in I} U_i$:
\begin{equation}\label{eq:sheaf_equalizer1}\begin{tikzcd}
    {C(U)} & {\prod_{i \in I}C(U_i)} & {\prod_{i,j \in I} C(U_i \cap U_j)}
    \arrow["e", from=1-1, to=1-2]
    \arrow["p", shift left, from=1-2, to=1-3]
    \arrow["q"', shift right, from=1-2, to=1-3]
\end{tikzcd}\end{equation}
this means that for any set $A$ map $a: A \to C(U)$, there is a unique arrow $\bar{e}: A \to C(U)$ such that $e \circ \bar{e} = a$, as in the diagram
\[\begin{tikzcd}
    {C(U)} & {\prod_{i \in I}C(U_i)} & {\prod_{i,j \in I} C(U_i \cap U_j)} \\
    A
    \arrow["e", from=1-1, to=1-2]
    \arrow["p", shift left, from=1-2, to=1-3]
    \arrow["q"', shift right, from=1-2, to=1-3]
    \arrow["{\bar{e}}", dashed, from=2-1, to=1-1]
    \arrow["a"', from=2-1, to=1-2]
\end{tikzcd}\]
Let $(g_i)_{i \in I}$ be an element in the image of $a$. The existence of the map $\bar{e}: A \to C(U)$ tells us that if $g_i|_{U_i \cap U_j} = g_j|_{U_i \cap U_j}$ for all $i, j \in I$ there exists an element $g \in C(U)$ such that $(e \circ \bar{e})(g) = (g_i)_{i \in I}$, i.e., $g|_{U_i} = g_i$ for all $i \in I$. This is precisely the gluing axiom of Definition \ref{defn:sheaf}. Moreover, the uniqueness of $\bar{e}$ tells us that the element $g$ is unique, which is precisely the identity axiom.

There was nothing special about the sheaf of continuous $\R$-valued functions; any sheaf of sets on $X$ can be described using an equalizer diagram of the form (\ref{eq:sheaf_equalizer1}). This motivates our second definition of a sheaf of sets, which is equivalent to Definition \ref{defn:sheaf}:

\begin{defn}\label{defn:sheaf2}
    Let $X$ be a topological space and suppose $\~F$ is a presheaf on $X$. Then $\~F$ is a sheaf if the following diagram is an equalizer for every open cover $\~U = \{U_i\}_{i \in I}$ of every open subset $U$ of $X$:
    \begin{equation}\label{eq:sheaf_equalizer2}
    \begin{tikzcd}
        {\~F(U)} & {\prod_i \~F(U_i)} & {\prod_{i,j} \~F(U_i \cap U_j)}
        \arrow["e",from=1-1, to=1-2]
        \arrow["p",shift left, from=1-2, to=1-3]
        \arrow["q"', shift right, from=1-2, to=1-3]
    \end{tikzcd}
    \end{equation}
    Here the map $e$ is the product of the restrictions $\~F(U) \to \~F(U_i)$ over all $i \in I$, whose $i^{th}$ component is given by $f|_{U_i}$ for all $f \in \~F(U)$. The map $p$ is the product of the restrictions $\~F(U_i) \to \~F(U_i \cap U_j)$ over all $i, j \in I$, whose $i, j$ component is given by $f_i|_{U_i \cap U_j}$ for all $f_i \in \~F(U_i)$, and similarly, the map $q$ is the product of the restrictions $\~F(U_j) \to \~F(U_i \cap U_j)$. Omitting the labels on arrows, this condition can be written as 
    \begin{equation}
    \label{eq:sheaf_limit}
        \~F(U) = \lim \left(\prod_i \~F(U_i) \rightrightarrows \prod_{i,j} \~F(U_i \cap U_j) \right)
    \end{equation}
\end{defn}

\begin{rmk}
    Many texts include the assumption that $\~F(\varnothing) = \{\ast\}$ (the terminal object in $\Set$). However, this axiom actually follows from Definition \ref{defn:sheaf2}: $\varnothing$ has an empty open covering, and a product over an empty index set is a singleton $\{\ast\}$, therefore the equalizer (\ref{eq:sheaf_equalizer2}) becomes $F(\varnothing) \to \{\ast\} \rightrightarrows \{\ast\}$, so that $\~F(\varnothing) = \{\ast\}$.
\end{rmk}

\begin{prop}
    If $\~F$ is a sheaf of sets on a topological space $X$ and $U, V \subseteq X$ are disjoint subsets of $X$, then $\~F(U \cup V) = \~F(U) \times \~F(V)$.
\end{prop}

\begin{proof}
    Since $U, V \subseteq X$ are open, so is $U \cup V$, and clearly $\{U, V\}$ is an open cover of $U \cup V$. Then using (\ref{eq:sheaf_limit}) above and the fact that $U \cap V = \varnothing$, we can write
    \begin{align*}
        \~F(U \cup V) &= \lim \left(\~F(U) \times \~F(V) \rightrightarrows \~F(U \cap V) \right) \\
        &= \lim \left(\~F(U) \times \~F(V) \rightrightarrows \~F(\varnothing) \right)\\
        &= \lim \left(\~F(U) \times \~F(V) \rightrightarrows 1\right)
    \end{align*}
    Now, $\lim \left(\~F(U) \times \~F(V) \rightrightarrows 1\right)$ is just the equalizer of the diagram $\~F(U) \times \~F(V) \rightrightarrows 1$, and this is clearly equal to $\~F(U) \times \~F(V)$, so we find that $\~F(U \cup V) = \~F(U) \times \~F(V)$.
\end{proof}

So far, we have explored some basic properties of sheaves, but we have not yet described the morphisms between these objects. Recall that the category $\PSh(\~C)$ of presheaves on $\~C$ is given by the functor category $\Fun(\~C^\op, \Set)$. Therefore a morphism from a presheaf $\~F$ to a presheaf $\~G$ in this category is simply a natural transformation from the functor $\~F$ to the functor $\~G$:
\[\begin{tikzcd}[column sep = scriptsize]
        {\~C^\op} & \Set
        \arrow[""{name=0, anchor=center, inner sep=0}, "\~F", bend left=50, from=1-1, to=1-2]
        \arrow[""{name=1, anchor=center, inner sep=0}, "\~G"', bend right=50, from=1-1, to=1-2]
        \arrow[shorten <=4pt, shorten >=4pt, Rightarrow, from=0, to=1]
    \end{tikzcd}\]
This is also a reasonable definition of morphism when $\~F$ and $\~G$ are both sheaves. 

\begin{defn}
If $\~F$ and $\~G$ are presheaves on a topological space $X$, then a \emph{morphism} $\vp: \~F \to \~G$ consists of a morphism $\vp_U: \~F(U) \to \~G(U)$ for each open set $U$, such that whenever $V \subseteq U$ in $X$, the diagram
\[\begin{tikzcd}
    {\~F(U)} & {\~F(V)} \\
    {\~G(U)} & {\~G(V)}
    \arrow["{\rho_{UV}}", from=1-1, to=1-2]
    \arrow["{\vp_U}"', from=1-1, to=2-1]
    \arrow["{\vp_V}", from=1-2, to=2-2]
    \arrow["{\rho'_{UV}}"', from=2-1, to=2-2]
\end{tikzcd}\]
commutes, where $\rho$ and $\rho'$ denote the restriction maps in $\~F$ and $\~G$ respectively. When $\~F$ and $\~G$ are sheaves we use exactly the same definition for a morphism. The morphism $\vp: \~F \to \~G$ is an \emph{isomorphism} if it has a two-sided inverse. We will of course write $\~F \isom \~G$ or $\~F \cong \~G$ when $\~F$ and $\~G$ are isomorphic. We write $\PSh(X)$ for the category of presheaves on $X$, and $\Sh(X)$ for the category of sheaves on $X$.
\end{defn}

 Recall that a \emph{basis} or \emph{base} for $X$ is a collection $\{B_i\}_{i \in I}$ of open subsets of $X$ with the property that any open set $U \subseteq X$ can be written as $U = \bigcup_{j \in J} B_j$ for some $J \subseteq I$. We have seen that studying sheaves on $X$ involves keeping track of data attached to each open set $U \subseteq X$, and working with this data can become quite complicated. If $\{B_i\}_{i \in I}$ is a basis for $X$ and $\~F \in \Sh(X)$, we can write any open $U \subseteq X$ as $U = \bigcup_{j \in J} B_j$ for some $J \subseteq I$. It is reasonable to ask if we can recover all of the information of $\~F(U)$ by passing to the covering $U = \bigcup_{j \in J} B_j$ and looking at the sets $\{\~F(B_j)\}_{j \in J}$. Indeed, we can, and in this section we make this notion precise. 

Given a sheaf $\~F$ on $X$ and a basis $\{B_i\}$, consider the sets $\~F(B_i)$, with their associated restriction maps
\[\rho_{ij} := \rho^{B_i}_{B_j}: \~F(B_i) \to \~F(B_j).\]
For each open $U \subseteq X$, write $U = \bigcup_{j \in J}B_j$ for some basic open sets $B_j$, and define a map $\iota_U$ by
\begin{align*}
    \iota_U: \~F(U) &\to \prod_{i \in I} \~F(B_i) \\ 
    s &\mapsto \left(s|_{B_i}\right)_{i \in I}.
\end{align*}
Clearly this map is injective, since $\vp(f) = \vp(g)$ means that $f|_{B_i} = g|_{B_i}$ on the cover $\{B_i\}_{i \in I}$, hence $f = g$ by locality. The image of $\iota_U$ is given by
\[\im(\iota_U) = \left\{(f_i)_{i \in I} \in  \prod_{i \in I} \~F(B_i) : f_i|_{B_k} = f_j|_{B_k} \text{ for all } i, j \in I \text{ and all } B_k \subseteq B_i \cap B_j \right\}.\]
We therefore identify $\~F(U)$ with its image under $\iota_U$, which allows us to think of $\~F(U)$ as sitting inside the product $\prod_{i \in I} \~F(B_i)$. This allows us to recover the value of $\~F$ on each open set $U$ by studying the values of $\~F$ on the basis $\{B_i\}$. This motivates the following construction: 
\begin{defn}
    Let $X$ be a topological space with a basis $\{B_i\}_{i\in I}$. A \emph{presheaf} $F$ \emph{on the basis $\{B_i\}_{i\in I}$} consists of a set $F(B_i)$ for each $i\in I$, together with restriction maps
    \[\rho_{ij} := \rho^{B_i}_{B_j}: F(B_i) \to F(B_j)\]
    whenever $B_j \subseteq B_i$, such that $\rho_{ii} = \id_{F(B_i)}$ for all $i$ and $\rho_{jk} \circ \rho_{ij} = \rho_{ik}$ whenever $B_k \subseteq B_j \subseteq B_i$. We say $\~F$ is a \emph{sheaf on the basis $\{B_i\}_{i\in I}$} if the following two axioms are satisfied:
    \begin{enumerate}[label=(\roman*)]
        \item (\emph{Base identity}) If $B = \bigcup_{i \in I}B_i$ and we have $f, g \in F(B)$ such that $f|_{B_i} = g|_{B_i}$ for all $i \in I$, then $f = g$.
        \item (\emph{Base gluability}) If $B = \bigcup_{i \in I} B_i$ and we have sections $f_i \in F(B)$ such that $f_i|_{B_k} = f_j|_{B_k}$ for all $B_k \subseteq B_i \cap B_j$, there exists $f \in F(B)$ such that $f|_{B_i} = f$ for all $i \in I$.
    \end{enumerate}
    We write $\Sh(B)$ for the category of sheaves on the basis $B$.
\end{defn}

\begin{rmk}
    If $X$ is a topological space with a basis $\{B_i\}_{i \in I}$, we write $F$ for a sheaf on this basis, and $\~F$ for the associated sheaf on the whole space $X$. 
\end{rmk}

If $B$ is a basis for $X$, it is clear that every sheaf $\~F$ on $X$ may be restricted to obtain a sheaf on the basis $B$. This gives a restriction functor
\[r: \Sh(X) \to \Sh(B),\]
which is an equivalence of categories (see Theorem 2.1.3 in \cite{SGL}).

\section{Grothendieck topoi} Suppose $\~C$ is a small category, and recall that $\PSh(\~C) = \Fun(\~C^\op, \Set)$ denotes the category of presheaves on $\~C$. Recall that $H_\bullet: \~C \to \PSh(\~C)$ denotes the Yoneda embedding ($H_\bullet(A) = \~C(-, A)$ for all $A \in \~C$). 
\begin{defn}
    A \emph{sieve} $S$ on an object $A \in \~C$ is a subfunctor $S \subseteq H_\bullet(A)$. More concretely, a sieve is a family $S$ of morphisms in $\~C$, all of which have codomain $A$, such that for all $f: B \to A$ in $S$ and all arrows $g: C \to B$, we have $f \circ g \in S$. 
\end{defn}
Hence a sieve is a collection of arrows with codomain $A$ which is closed under precomposition. Sieves therefore play a similar role to right ideals in ring theory. If $S$ is a sieve on $A \in \~C$ and $h: B \to A$ is any arrow in $\~C$, we can define a sieve $h^\ast(S)$ on $B$ by 
\[h^\ast(S) := \left\{g : \cod(g) = B, \, h\circ g \in S \right\}.\]
Moreover, for any object $A \in \~C$, the collection of all possible maps into $A$ clearly defines a sieve $t_A$, called the \emph{maximal sieve}:
\[t_A := \left\{f : \cod(f) = A\right\}.\]
Notice that if $S$ is a sieve on $A \in \~C$, we have 
\[ S = t_A \iff 1_A \in S, \]
this is another instance of the similarity between sieves and right ideals.

\begin{defn}\label{defn:GT}
    A \emph{Grothendieck topology} on a category $\~C$ is a function $J$ which assigns to each object $A \in \~C$ a collection $J(A)$ of sieves on $\~C$, such that:
    \begin{enumerate}
        \item \label{GT:1} The maximal sieve $t_A$ is in $J(A)$. 
        \item \label{GT:2}(\emph{Stability axiom}) If $S \in J(A)$, then $h^\ast(S) \in J(B)$ for any arrow $h: B \to A$.   
        \item \label{GT:3}(\emph{Transitivity axiom}) If $S \in J(A)$ and $R$ is a sieve on $A$ such that $f^*(R) \in J(B)$ for all $f: B \to A$ in $S$, then $R \in J(A)$.
    \end{enumerate}
    A \emph{site} is a pair $(\~C, J)$ consisting of a small category $\~C$ and a Grothendieck topology $J$ on $\~C$.
\end{defn}

Suppose $(\~C, J)$ is a site. Then if $S \in J(A)$ for some $A \in \~C$, any larger sieve $R \supseteq S$ on $A$ must also lie in $J(A)$. To see this, suppose we have an arrow $h: B \to A$ in $S$. Then certainly $1_B \in h^*(S)$, so $h^*(S) = t_B$, the maximal sieve on $B$. As $h^*(S) \subseteq h^*(R)$, we must also have $h^*(R) = t_B$, hence $h^*(R) \in J(D)$ by Definition \ref{defn:GT}. This is true for all $h: B \to A$ in $S$, so transitivity gives $R \in J(A)$, as claimed.

We can rephrase the axioms of a Grothendieck topology using the more suggestive terminology of ``coverings''. If $S \in J(A)$, we call $S$ a \emph{covering sieve} and say that $S$ \emph{covers} $A$. We say that $S$ \emph{covers an arrow $f: B \to A$} if $f^*(S)$ covers $B$ (therefore $S$ covers $A$ if and only if $S$ covers the identity arrow on $A$).


%\printglossary[title={Index of Notation}]

\printbibliography

\end{document}